\documentclass[11pt,oneside,openany]{book}
\usepackage[T1]{fontenc}
\usepackage{lmodern}
\usepackage{microtype}
\usepackage[a4paper,margin=28mm,headheight=14pt]{geometry}
\usepackage{amsmath,amssymb,amsthm,mathtools}
\usepackage[hidelinks,bookmarksopen,bookmarksnumbered]{hyperref}
\hypersetup{pdftitle={Foundations of Algebraic Topology and Riemannian Geometry},
  pdfauthor={MathGenome},pdfsubject={Unified notes, P2-02}}
\allowdisplaybreaks[1]
\setlength{\emergencystretch}{2em}
\theoremstyle{plain}
\newtheorem{theorem}{Theorem}[chapter]
\newtheorem{lemma}[theorem]{Lemma}
\newtheorem{proposition}[theorem]{Proposition}
\newtheorem{corollary}[theorem]{Corollary}
\theoremstyle{definition}
\newtheorem{definition}[theorem]{Definition}
\newtheorem{example}[theorem]{Example}
\theoremstyle{remark}
\newtheorem{remark}[theorem]{Remark}
\newcommand{\R}{\mathbb R}
\newcommand{\C}{\mathbb C}
\newcommand{\Z}{\mathbb Z}
\newcommand{\Q}{\mathbb Q}
\newcommand{\N}{\mathbb N}
\newcommand{\norm}[1]{\lVert #1\rVert}
\newcommand{\pair}[2]{\langle #1,#2\rangle}
\DeclareMathOperator{\id}{id}
\DeclareMathOperator{\im}{im}
\DeclareMathOperator{\Tor}{Tor}
\DeclareMathOperator{\Ext}{Ext}
\DeclareMathOperator{\supp}{supp}
\DeclareMathOperator{\vol}{vol}
\DeclareMathOperator{\Ric}{Ric}
\DeclareMathOperator{\Scal}{Scal}
\DeclareMathOperator{\grad}{grad}
\DeclareMathOperator{\Div}{div}
\DeclareMathOperator{\tr}{tr}

\begin{document}
\frontmatter
\hypersetup{pageanchor=false}
\begin{titlepage}
\thispagestyle{empty}
\vspace*{27mm}
{\large MathGenome \enspace / \enspace P2-02\par}
\vspace{8mm}
\hrule height 0.8pt
\vspace{11mm}
{\Huge\bfseries Foundations of Algebraic Topology\\[2mm]
and Riemannian Geometry\par}
\vspace{12mm}
{\large Unified notes\par}
\vspace{6mm}
{\normalsize Coverings, chain models, and duality\\
Connections, geodesics, curvature, and global geometry\par}
\vfill
{\large 8 November 2026\par}
\end{titlepage}
\hypersetup{pageanchor=true}
\tableofcontents
\chapter*{Scope and conventions}
\addcontentsline{toc}{chapter}{Scope and conventions}

Part I develops ordinary algebraic topology from paths and coverings to
singular, simplicial, and cellular homology. It then develops ordinary
cohomology, products, duality, and the de Rham comparison. Part II develops
the metric and connection on a finite-dimensional manifold, its geodesics
and curvature, and the comparison and global theorems that follow.
The two parts share covering theory and the Euler characteristic.

For ordinary topology we write coefficient rings or groups at the
point of use. A pair is a space with a specified subspace. Ordinary
singular chains are finite sums, and cochains are their algebraic
duals with the stated coefficients. Where compact supports are needed,
the support condition is included explicitly. We use point-set
topology from \cite{D10}. The algebraic coefficient and K\"unneth
results used here are proved in \cite[Chapter 11]{P103}; see also
\cite{Weibel}. The prerequisites for smooth manifolds, differential
forms, integration, and Stokes' theorem are given in \cite{P104},
with \cite{LeeSM} as a reference.

Riemannian manifolds have smooth positive-definite metrics and are
finite-dimensional, Hausdorff and second countable. They have no
boundary unless a boundary is specified. Distance is finite on each
connected component. Our curvature convention is
\[
 R(X,Y)Z=\nabla_X\nabla_Y Z-\nabla_Y\nabla_X Z-\nabla_{[X,Y]}Z,
 \qquad
 K(X,Y)=\frac{\langle R(X,Y)Y,X\rangle}
 {\langle X,X\rangle\langle Y,Y\rangle-\langle X,Y\rangle^2}.
\]
The convention makes the sphere positively curved. We use
\(\Delta_g=\operatorname{div}_g\operatorname{grad}_g\) and
\(E(\gamma)=\tfrac12\int\langle\dot\gamma,\dot\gamma\rangle\,dt\).
Differential equations and flows use their actual maximal domains
\cite{D13,P104}.

Hatcher is the main topological reference, and Munkres--Krantz--Parks
is a second route through chain models and duality. Munkres provides
a comparison for coverings. Lee and do Carmo provide complementary
routes through the geometry. We follow the corrected hypotheses in
the reference editions and give the proof interfaces at their uses.

\mainmatter
\part{Algebraic Topology}
\chapter{Fundamental Groups and Covering Spaces}

Throughout this chapter, paths are parametrized by \(I=[0,1]\).
The product \(ab\) traverses \(a\) first and \(b\) second, and
\(\bar a(t)=a(1-t)\). Path homotopies hold both endpoints fixed
unless specified otherwise. All covering maps are surjective.

\section{Paths and the fundamental group}\label{p202:T01}

\begin{definition}[Path classes and the fundamental group]
For paths \(a\) and \(b\) with \(a(1)=b(0)\), define \(ab\) by
running each path on one half of \(I\). A homotopy of paths is
a continuous map \(I\times I\to X\) with prescribed first and
last paths and constant values on the two endpoint sides.
At \(x\in X\), \(\pi_1(X,x)\) consists of classes of loops
based at \(x\), with product induced by concatenation.
\end{definition}

\begin{proposition}[Functoriality and change of basepoint]
\label{prop:p202-path-classes}
Path classes form a groupoid, so the based loops at each point
form a group. A map of pointed spaces induces a homomorphism.
Based homotopic maps induce the same homomorphism, and a
pointed homotopy equivalence induces an isomorphism. A path
\(d:x\to y\) gives an isomorphism
\(\pi_1(X,y)\to\pi_1(X,x)\) by \([b]\mapsto[d b\bar d]\).
Changing \(d\) changes the isomorphism by an inner automorphism.
\end{proposition}

\begin{proof}
Concatenating homotopies on the two half squares proves that
the product descends to classes. A continuous nondecreasing
reparametrization \(\phi:I\to I\) with the same endpoints
changes a path by the homotopy \(a((1-s)t+s\phi(t))\).
In particular, after choosing a common three-piece path,
different assignments of time to its three pieces agree on
classes. This proves associativity. The same argument removes
an initial or terminal constant segment. To contract \(a\bar
a\), let \(\lambda(t)=2\min(t,1-t)\) and use
\(a((1-s)\lambda(t))\). The reverse product contracts by using
\(\bar a\). Postcomposition respects these formulas and sends
a based homotopy applied to any loop to a homotopy relative to
endpoints. These facts prove functoriality and the homotopy
assertions. The basepoint formula has inverse defined using
\(\bar d\). For \(d'\) with the same endpoints, the new formula
equals the old one followed by conjugation by \([d'\bar d]\).
\end{proof}

For an unbased homotopy \(F:X\times I\to Y\) from \(f\) to \(g\),
the same square for a loop based at \(x\) gives
\([f\circ a]=[d(g\circ a)\bar d]\), where \(d(t)=F(x,t)\).
This is the basepoint correction needed when using an ordinary
rather than a pointed homotopy. See \cite[Sections 0 and 1.1]{Hatcher}.

\section{Path and homotopy lifting}\label{p202:T05}

\begin{definition}[Covering map]
A surjective map \(p:E\to B\) is a covering if each \(b\in B\)
has an open neighborhood \(U\) such that \(p^{-1}U\) is a disjoint
union of open subsets on each of which \(p\) is a homeomorphism
onto \(U\). We call such a \(U\) evenly covered and the open
subsets its sheets.
\end{definition}

\begin{theorem}[Lifting with prescribed initial values]
\label{thm:p202-cover-lifting}
Each path in \(B\) has a unique lift starting at a specified
point over its initial point. More generally, if
\(F:Z\times I\to B\) is continuous and \(F(\cdot,0)\) has a
given continuous lift \(h:Z\to E\), there is a unique
continuous lift of \(F\) extending \(h\). The space \(Z\)
can be arbitrary.
\end{theorem}

\begin{proof}
For a path, pull back the evenly covered neighborhoods to
\(I\) and use a Lebesgue number to choose a finite subdivision
such that each closed subinterval maps to one neighborhood.
Starting from the given lift, on each interval take the
inverse of the sheet that contains the preceding endpoint.
The pieces paste. For uniqueness, two lifts agree on an open
and closed subset of \(I\): locally both are obtained using
inverse branches, and branches over the same open set either
agree or have disjoint images. The agreement set contains the
initial point.

For the homotopy, use the path construction at every \(z\).
To prove continuity near a given \(z_0\), choose a finite
subdivision for \(F(z_0,\cdot)\). Compactness of each closed
subinterval and the tube lemma give a neighborhood \(W\) of
\(z_0\) on which the image of that whole subinterval stays in
the same evenly covered neighborhood. Intersect these
finitely many neighborhoods. Shrink once at the first endpoint
to put \(h(W)\) in the chosen first sheet. On the first slab
the inverse branch now gives a continuous formula. The lifted
right endpoint is continuous, so shrink again to put it in
the second sheet. Continue through the finitely many slabs.
The resulting local lift on \(W\times I\) is continuous by
pasting and agrees with all the pointwise lifts by uniqueness.
\end{proof}

A lift of a constant path is constant. In particular, the
endpoints of the lifts of a path homotopy relative to endpoints
are fixed if the initial lift is fixed. This applies to the
argument square even when the total space is not assumed
connected. See \cite[Section 1.3]{Hatcher} and
\cite[Sections 53--54]{MT}.

\section{Seifert--van Kampen}\label{p202:T03}

\begin{theorem}[The two-open-set pushout]
\label{thm:p202-van-kampen}
Let \(X=U\cup V\) where \(U,V\) are open and path connected,
and \(U\cap V\) is nonempty and path connected. With a basepoint
\(x\in U\cap V\), the inclusions exhibit \(\pi_1(X,x)\) as
the group pushout of \(\pi_1(U,x)\leftarrow\pi_1(U\cap V,x)
\to\pi_1(V,x)\).
\end{theorem}

\begin{proof}
Let \(P\) be the free product of the two groups modulo the
normal closure of the relations that identify the two images
of each element from the intersection. The natural
homomorphism \(P\to\pi_1(X,x)\) is defined. We construct its
inverse. Choose a path \(d_y\) from \(x\) to each \(y\in X\).
Choose it in \(U\cap V\) for \(y\in U\cap V\), in the appropriate
open set otherwise, and choose \(d_x\) constant.

For a path \(a\), subdivide it into closed subintervals whose
images each lie in \(U\) or \(V\). For a piece from \(y\) to
\(z\) take the class in \(P\) of the based loop \(d_y a_i\bar
d_z\), using the group of the chosen open set. The product of
these classes is unchanged by refinement, since adjacent
connecting paths cancel within the same group. Any two
partitions have a common refinement. When the label of a
piece changes on this refinement, that piece and its two
connecting paths lie in \(U\cap V\), and the imposed relation
gives equality. Thus the word is well defined.

Suppose two based loops are homotopic. Pull the two open sets
back to the homotopy square and use a sufficiently fine
rectangular grid such that every closed small rectangle lies
in one inverse image. For each oriented edge define its word
as above. On a small rectangle, all four words can be computed
in the same group, where their boundary product is trivial
because the homotopy fills that rectangle. Equivalently, the
word along its bottom and right edges equals the word along
its left and top edges. Sweeping across a row and cancelling
the intervening vertical words gives this equality for the
row. Sweeping through the rows gives it for the entire square.
The vertical sides give identity words because the homotopy
is based. The two loop words therefore agree.

Concatenation gives multiplication of words, so this defines
the required homomorphism \(\pi_1(X,x)\to P\). Under
\(P\to\pi_1(X,x)\) the connecting paths cancel and recover the
original loop class. In the other direction a loop in \(U\)
or \(V\) may be taken as a single piece, yielding its own
generator in \(P\). These generators generate the quotient.
\end{proof}

This argument makes explicit the open-cover and intersection
conditions used to identify relations. See \cite[Section 1.2]{Hatcher}
and \cite[Sections 70--71]{MT}.

\section{Circle, products, and spheres}\label{p202:T02}

\begin{theorem}[Basic fundamental groups]
\label{thm:p202-basic-pi1}
The winding number gives \(\pi_1(S^1,1)\cong\mathbb Z\).
For a product of pointed spaces the fundamental group is the
product of the fundamental groups. For \(n\ge2\),
\(S^n\) is simply connected.
\end{theorem}

\begin{proof}
The map \(q:\mathbb R\to S^1\), \(q(t)=\exp(2\pi i t)\),
is a covering: a proper sufficiently short open arc lifts to
a disjoint union of intervals differing by integer
translations. Lift a loop at \(1\) starting at \(0\). Its
endpoint is an integer and is unchanged by based homotopy
by Theorem~\ref{thm:p202-cover-lifting}. Concatenation adds
endpoints, since a translated lift has the same projection.
Every integer \(m\) occurs from \(q(mt)\). Two lifts with the
same endpoint are homotopic relative to endpoints by their
straight-line interpolation in \(\mathbb R\), and projecting
gives a based homotopy. This proves the first assertion.

A map into a product is continuous exactly when its
coordinate maps are continuous. Apply this fact to loops and
homotopies and use coordinatewise concatenation. For the sphere,
use the complements of its two opposite poles. Each is
homeomorphic to \(\mathbb R^n\) and their intersection is
homeomorphic to \(S^{n-1}\times\mathbb R\), which is path
connected for \(n\ge2\). Theorem~\ref{thm:p202-van-kampen}
makes the group trivial.
\end{proof}

\section{The lifting criterion and monodromy}\label{p202:T06}

\begin{theorem}[Lifting a map]
\label{thm:p202-cover-criterion}
Let \(p:(E,e)\to(B,b)\) be a covering and
\(f:(Y,y)\to(B,b)\) a pointed map with \(Y\) path connected
and locally path connected. A lift taking \(y\) to \(e\)
exists exactly when
\[
 f_*\pi_1(Y,y)\subset p_*\pi_1(E,e).
\]
It is unique. Moreover \(p_*\) is injective.
\end{theorem}

\begin{proof}
For injectivity, lift a based null homotopy of the projection
of a loop in \(E\), starting with that loop. Uniqueness on the
constant boundary edges fixes the lifted endpoints, and the
terminal lift of the constant loop is constant. The subgroup
condition is necessary by composition. For sufficiency, to
define the lift at \(z\in Y\), choose a path \(a:y\to z\) and
take the endpoint of the lift of \(fa\) starting at \(e\).
A based loop at \(b\) lifts to a loop at \(e\) exactly when
its class lies in \(p_*\pi_1(E,e)\), by lifting an endpoint
homotopy to any projected representative. Thus two choices
of \(a\) give the same endpoint by applying the subgroup
condition to \(a\bar a'\) and using uniqueness of the lifts
of reversed paths. For continuity at \(z\), choose a path
connected open neighborhood of \(z\) inside the preimage
of an evenly covered set. Extending the chosen path to any
point of this neighborhood by a path within it gives exactly
the same local inverse branch of \(p\). This proves
continuity, and lifting along each chosen path proves
uniqueness.
\end{proof}

\begin{corollary}[Action on the fiber]
\label{cor:p202-cover-action}
For \(G=\pi_1(B,b)\) the endpoint of a lifted based loop
defines a right \(G\)-action on \(p^{-1}(b)\). If \(E\) and
\(B\) are path connected, this action is transitive and,
for \(H=p_*\pi_1(E,e)\), the fiber is \(G\)-equivariantly
identified with \(H\backslash G=\{Hg:g\in G\}\).
The number of sheets is the index of \(H\).
\end{corollary}

\begin{proof}
Homotopy lifting gives a well-defined endpoint on loop
classes, and consecutive lifts give the action law. A path
in \(E\) between two fiber points projects to a based loop,
proving transitivity. The loops that fix \(e\) are exactly
those in \(H\), as in the preceding proof. The assignment
\(Hg\mapsto e\cdot g\) is therefore a bijection. Lifting
a path between two base points and its reverse gives
inverse bijections of their fibers. Each evenly covered
neighborhood has one sheet for each point in its fiber.
\end{proof}

\section{Universal and classified coverings}\label{p202:T07}

A space is semilocally simply connected when each point has
an open neighborhood whose based loops at that point are
null homotopic in the whole space. With local path
connectedness such neighborhoods may be taken path
connected, and then their loops at any point are null
homotopic in the whole space, by change of basepoint.

\begin{theorem}[Classification of connected coverings]
\label{thm:p202-cover-classification}
Let \(B\) be path connected, locally path connected, and
semilocally simply connected. Fix \(b\in B\) and put
\(G=\pi_1(B,b)\). Isomorphism classes of connected pointed
coverings of \(B\) correspond to subgroups of \(G\), by
taking the image of the fundamental group. Forgetting the
chosen lift gives the classification by conjugacy classes
of subgroups. The trivial subgroup gives a simply connected
covering, which maps by a unique pointed lift to each
connected pointed covering.
\end{theorem}

\begin{proof}
For a subgroup \(H\), consider paths from \(b\).
Identify \(\alpha,\beta\) with the same endpoint exactly if
\([\alpha\bar\beta]\in H\). This is an equivalence relation
by the path groupoid law. Denote the class by
\([\alpha]_H\), and project to its endpoint. For a path
connected open \(U\) whose loops contract in \(B\), a class
\([\alpha]_H\) over \(z\in U\) defines a subset consisting of
all \([\alpha\eta]_H\) for paths \(\eta\) in \(U\) starting at
\(z\). For each endpoint in \(U\) the resulting class is
unique, since the loop formed from two choices of \(\eta\)
contracts in \(B\). Such subsets cover the preimage of \(U\)
and are disjoint or equal, as one sees by extending any
common class along paths in \(U\) and by reversing these
paths. Declare them to be basis sets. At an intersection,
take a smaller neighborhood of the same kind inside the
intersection of the base neighborhoods and extend the
common class along paths there. This verifies the basis
axiom. The projection maps each basic sheet homeomorphically
onto its open base neighborhood: inverse images of smaller
neighborhoods are unions of basic sets, and every open
subset of a sheet projects to an open subset of \(U\).
It follows that the resulting space \(E_H\) is a covering.

The prefix of a path \(\alpha\), stopped at parameter \(t\),
defines a continuous path of classes starting at the class
of the constant path and ending at \([\alpha]_H\). Continuity
follows locally from the sheet construction, since nearby
endpoints differ by a short piece of \(\alpha\) inside an
eligible \(U\). Hence \(E_H\) is path connected. The lifted
endpoint of a loop \(a\) at \(b\) is \([a]_H\), so the loop
lifts closed at the distinguished point exactly for classes
in \(H\). This identifies the image subgroup. The total
space is locally path connected because it is locally
homeomorphic to \(B\); consequently, a connected covering
over this base is path connected.

Any two such coverings with the same subgroup admit pointed
maps in both directions over \(B\) by
Theorem~\ref{thm:p202-cover-criterion}, applied to the
projection of the source. Uniqueness makes the composites
identities. Choosing the distinguished point \(e\cdot g\)
instead of \(e\) changes the stabilizer to \(g^{-1}Hg\) by
Corollary~\ref{cor:p202-cover-action}. This proves the
unpointed classification. The covering \(E_{\{1\}}\) has
trivial fundamental group by injectivity, and the lifting
criterion gives its pointed universal property.
\end{proof}

\section{Deck transformations and regular coverings}\label{p202:T08}

\begin{theorem}[The deck group]
\label{thm:p202-deck-group}
Under the base hypotheses of
Theorem~\ref{thm:p202-cover-classification}, a connected
covering with subgroup \(H\le G\) has deck group
\(N_G(H)/H\), where a deck transformation is a
homeomorphism of the total space over the base.
The deck action on the fiber is transitive exactly when
\(H\) is normal. In that case the covering is called
regular.
\end{theorem}

\begin{proof}
Work in \(E_H\). If the loop class \(g\) belongs to
\(N_G(H)\), prepending its loop representative to every
path defines \(F_g([\alpha]_H)=[g\alpha]_H\). The equivalence
condition is respected because
\([g\alpha\bar\beta\bar g]\in H\) exactly when
\([\alpha\bar\beta]\in H\). The formula preserves the
projection and carries a basic sheet to a basic sheet.
It is continuous with inverse \(F_{g^{-1}}\). Prepending
shows \(F_g\circ F_h=F_{gh}\), and \(F_g\) fixes the
distinguished point exactly for \(g\in H\).
Uniqueness of a lift implies that a deck transformation is
determined by its value on that point. If this value is
\(e\cdot g\), its stabilizer under monodromy is both \(H\)
and \(g^{-1}Hg\), since the deck transformation commutes with
path lifting. Thus \(g\in N_G(H)\), and every deck
transformation has the displayed form. The orbit of \(e\)
is \(H\backslash N_G(H)\); equality with \(H\backslash G\)
is equivalent to \(N_G(H)=G\).
\end{proof}

\begin{example}[Circle and real projective space]
The \(m\)-sheeted circle cover \(z\mapsto z^m\), for \(m\ge1\),
has subgroup \(m\mathbb Z\) by winding number. Its deck
transformations are multiplications by \(m\)-th roots of
unity, which realize \(\mathbb Z/m\mathbb Z\).
For \(n\ge2\) the quotient of \(S^n\) by the antipodal
involution is \(\mathbb{RP}^n\). A small open ball disjoint
from its antipodal image maps homeomorphically to its image,
whose full inverse image is the two balls, so this is a
two-sheeted covering. Since \(S^n\) is simply connected, its
subgroup is trivial, and the classification or the fiber
action gives \(\pi_1(\mathbb{RP}^n)\cong\mathbb Z/2\mathbb Z\).
\end{example}


\chapter{Cell Structures and Approximation}

\section{Cells, skeleta, and simplices}\label{p202:T09}

\begin{definition}[CW and simplicial structures]
A CW complex is a Hausdorff space built from disks by
successive attachments. The \(n\)-skeleton \(X^n\) is the
pushout of \(X^{n-1}\) and the disjoint union of the closed
\(n\)-disks, identifying the boundary spheres by the
attaching maps. A characteristic map is the induced
map from a closed disk, and its restriction to the interior
is a homeomorphism onto an open cell. This cell need not be
open in the full space. The closure of a cell is contained
in a finite union of cells, and the whole space has the weak
topology with respect to the closed characteristic disks.
A subcomplex consists of cells together with their
closures.

An abstract simplicial complex on a vertex set is a
collection of nonempty finite subsets closed under taking
nonempty subsets. In its geometric realization each
simplex is represented by nonnegative barycentric coordinates
with sum one and support in that simplex. Faces agree by
setting the other coordinates to zero. We use the weak
realization topology. A \(\Delta\)-complex instead glues
ordered geometric simplices by order-preserving face maps,
allowing distinct faces to have the same image and different
simplices to share the same vertices.
\end{definition}

\begin{proposition}[Basic properties of weak cell structures]
\label{prop:p202-cw-basics}
Subcomplexes are closed and carry the subspace topology.
Every compact subset of a CW complex is contained in a
finite subcomplex. The geometric realization of a simplicial
complex is a CW complex with one cell per simplex interior.
The same construction gives the associated CW structure
of a \(\Delta\)-complex.
\end{proposition}

\begin{proof}
The attachment topology says a subset is open if its inverse
image in the old skeleton and each closed disk is open.
An open subset of the old skeleton extends to the next one:
in each disk add an outer collar over the inverse image of
that subset on the sphere. This collar can use \(r>1/2\) in
radial coordinates. Iterating proves that each skeleton has
its original topology as a subspace. Disjoint open sets of
the old skeleton can be extended disjointly by this construction.
An interior point of a new disk can be separated from a
distinct point using an interior ball, and the outer collars
can be made thin enough to miss that ball's closure. These
observations also prove that the weak topology of a cell
attachment to a Hausdorff skeleton is Hausdorff.

For a subcomplex \(A\), the inverse image of \(A\) in the
disk of a cell contained in \(A\) is the whole disk. In the
disk of a cell outside \(A\) it is a subset of the boundary
and is closed by induction on dimension, using the attaching
map. The weak topology makes \(A\) closed. Similarly, an
open set of \(A\) in its own weak topology extends to \(X\):
inductively extend its old-skeleton part, use the given set
in the disks belonging to \(A\), and extend by outer collars
in the other disks. This verifies the topology assertion.

If a compact set met infinitely many open cells, choose a
point in each of a countably infinite collection of these
cells. Every subset of the chosen points is closed: its
intersection with the image of each characteristic disk is
finite by closure finiteness and hence is closed there.
Thus the chosen points form an infinite closed discrete
subspace of the compact set, contradicting compactness.
The finitely many cells met by the compact set lie in a
finite subcomplex by repeatedly adjoining all cells in
their boundaries, a finite process with decreasing
dimension.

In either simplex realization the characteristic maps are
the closed simplex maps. At dimension \(n\) their boundary
identifications are precisely the attachment to the
\((n-1)\)-skeleton. Each simplex has finitely many faces,
also when some are identified. Distinct simplex interiors
are disjoint, and the weak topology is the one used in
the definition. The Hausdorff argument above applies to
these successive attachments.
\end{proof}

We use products with \(I\) in the ordinary product topology.
Here the weak-topology test still works. If \(q:P\to Q\)
is quotient, then \(q\times\mathrm{id}_I\) is quotient.
Indeed, suppose the inverse image of \(O\subset Q\times I\)
is open. For \((y,t)\in O\) and a lift of \(y\), choose
a compact relative interval \(K\) with \(t\) in its relative
interior and \(\{y\}\times K\subset O\). The set of
\(p\in P\) whose product with \(K\) lies in the inverse
image is open by the tube lemma and is saturated for \(q\).
Its image is an open neighborhood of \(y\), whose product
with the relative interior of \(K\) lies in \(O\).
Apply this to the quotient from the disjoint union of the
characteristic disks. These observations justify the weak
pasting arguments with homotopies below. For arbitrary
products of infinite CW complexes the analogous assertion
requires attention to the topology. See \cite[Chapter 0 and
Appendix]{Hatcher} and \cite[Sections 1.2--1.5]{MKP}.

\section{Homotopy extension and good pairs}\label{p202:T10}

\begin{theorem}[Homotopy extension for a CW pair]
\label{thm:p202-cw-hep}
Let \(A\) be a subcomplex of \(X\). Given a map \(X\to Y\)
and a homotopy on \(A\) starting at its restriction, the
homotopy extends over \(X\).
\end{theorem}

\begin{proof}
The disk cylinder retracts onto
\(D^n\times\{0\}\cup\partial D^n\times I\).
One explicit retraction for \(n\ge1\) follows the ray
from \((0,2)\) through \((x,t)\) to its last point in the
cylinder. Its ray parameter, with \((x,t)\) at parameter
one, is the minimum of \(2/(2-t)\) and \(1/\lVert x\rVert\),
where the second term is omitted at \(x=0\).
This parameter is continuous and gives the prescribed
identity on the bottom and side. For \(n=0\) use the
constant retraction onto the bottom. At each new cell
outside \(A\), pull back the prescribed map on the bottom
and the previously extended homotopy on the side, then
compose with this retraction. In the cells in \(A\) use
the given homotopy. The resulting maps agree on all
boundaries and are continuous on \(X\times I\) by the
weak-topology test.
\end{proof}

\begin{proposition}[Neighborhood and quotient consequences]
\label{prop:p202-cw-neighborhood}
A subcomplex \(A\) is a strong deformation retract of an
open neighborhood. If \(A\) is nonempty and contractible,
\(X\to X/A\) is a homotopy equivalence. If attaching maps
\(S^{n-1}\to X\) are changed by a homotopy, the attached
spaces are homotopy equivalent relative to \(X\).
\end{proposition}

\begin{proof}
Construct the neighborhood by dimension, starting with
the vertices in \(A\). Suppose \(W\) in the old skeleton
deforms to its part of \(A\) by \(h\). Include every new
disk belonging to \(A\). In any other new disk with
attaching map \(a\), include the open outer collar
\(r>1/2\) over all directions \(\theta\) with
\(a(\theta)\in W\). Together with \(W\) these form an open
set in the new skeleton. At \((r,\theta)\) in such a
collar spend time \(1-r\) moving radially to the boundary.
For the remaining time use
\(h(a(\theta),(t-(1-r))/r)\). At \(r=1\) use \(h\) from
time zero. This formula pastes continuously along the
boundary, since the initial radial interval tends to
zero and the later time tends to \(t\). On cells in \(A\)
use the identity. The deformations are compatible on all
skeleta and fix \(A\), hence paste in the weak topology
to the desired deformation.

For the quotient assertion, extend a contraction of \(A\)
inside \(A\) to a homotopy \(H\) of the identity of \(X\)
using Theorem~\ref{thm:p202-cw-hep}. The terminal map is
constant on \(A\) and factors as \(j\circ q\) for
\(q:X\to X/A\). Also \(qH\) is constant on \(A\) at
each time, so it factors through a homotopy on
\((X/A)\times I\) by the product quotient property.
This gives \(qj\simeq\mathrm{id}\) as well as
\(jq\simeq\mathrm{id}\).

For the attachment assertion, let \(a_t\) be a homotopy
from \(a_0\) to \(a_1\). Define a map between the
attachments, fixed on \(X\), by sending the inner
half-radius disk to the entire new disk by rescaling.
On the outer annulus send \((r,\theta)\) to
\(a_{2(1-r)}(\theta)\) in \(X\). At the outer boundary
this is \(a_0\), and at the inner interface it agrees
with the boundary of the rescaled disk. The reverse
construction uses \(a_{1-t}\). Each composite sends an
inner disk to the whole disk and uses the original
attaching-map homotopy followed by its reverse on the
remaining annuli. Contract this out-and-back homotopy
through such annulus formulas, using a linear contraction
of the time excursion. With a constant homotopy on the
annulus, increasing the inner radius to one gives the
identity. This proves the claimed homotopy equivalence.
\end{proof}

\section{Graphs and attaching presentations}\label{p202:T04}

\begin{theorem}[Graph groups and attached cells]
\label{thm:p202-graph-attach}
The fundamental group of a connected graph is free on its
edges outside a maximal tree, with one orientation chosen
per such edge. Attaching a \(2\)-cell to a path connected
space divides its fundamental group by the normal closure
of the attaching loop. Attaching a cell of dimension
at least \(3\) leaves the group unchanged. For a connected
finite CW complex this gives the presentation using the
edges outside a maximal tree of the \(1\)-skeleton as
generators and the \(2\)-cell attaching loops as relations.
\end{theorem}

\begin{proof}
A maximal acyclic subgraph containing all vertices exists
by the union argument for chains of forests and Zorn's
lemma. If it had two components, an edge on a graph path
joining components could be added without creating a cycle.
Thus it is a spanning tree. Fix a root in the tree. Give
each edge length one, and move each point along its unique
path to the root to the fraction \(1-s\) of its original
distance at time \(s\). On every edge the paths lie in
a finite subtree and the formula is continuous there, so
weak pasting gives a contraction. Collapsing the tree is a
homotopy equivalence by Proposition~\ref{prop:p202-cw-neighborhood}
and gives a bouquet of circles. For a finite bouquet,
van Kampen gives the free product of the groups of the
circles. For instance, to add one circle, use as one open
set the old bouquet with a contractible arc around the
basepoint in the new circle and as the other open set the
new circle with such arcs in each old circle. Their
intersection is a contractible star of arcs. An arbitrary
loop or null homotopy in an infinite bouquet has compact
image, hence lies in a finite subbouquet. The finite free
product calculation then proves the free-group description
also in the infinite case.

For one attached \(n\)-disk, cover the adjunction by the
interior of the disk and the complement of a smaller
closed central disk. The latter retracts onto the old
space by pushing the annular collar out radially, and
their intersection retracts to \(S^{n-1}\). For \(n\ge2\)
these sets and their intersection are path connected.
Theorem~\ref{thm:p202-van-kampen} gives the quotient by the
normal closure of the image of \(\pi_1(S^{n-1})\). Under
the deformation to the old space this image comes from the
attaching map; a path to the chosen basepoint only
conjugates it. The image is generated by the attaching
loop for \(n=2\) and is trivial for \(n\ge3\) by
Section~\ref{p202:T02}. In a connected finite CW complex
the \(1\)-skeleton is connected, since a cell of dimension
at least two attaches along a connected sphere into one
component and cannot join components. Applying the
one-cell calculation successively proves the presentation.
\end{proof}

\begin{example}[A surface group]
For the polygon model of a closed orientable surface of
genus \(g\ge1\), the edge identifications give one vertex,
edges \(a_1,b_1,\ldots,a_g,b_g\), and one disk whose
boundary reads \(\prod_{i=1}^g a_i b_i a_i^{-1}b_i^{-1}\).
Thus the group has these generators and this one relation.
For \(g=1\) this is \(\mathbb Z^2\), agreeing with the
product calculation for \(S^1\times S^1\).
\end{example}

\section{Simplicial approximation}\label{p202:T11}

\begin{theorem}[Finite-source simplicial approximation]
\label{thm:p202-simplicial-approx}
For a continuous map \(f:|K|\to |L|\) with \(K\) a finite
simplicial complex, a sufficiently iterated barycentric
subdivision \(K'\) admits a simplicial map \(g:K'\to L\)
whose realization is homotopic to \(f\).
\end{theorem}

\begin{proof}
Barycentric subdivision uses as vertices the barycenters
of nonempty faces, with simplices the strictly nested
chains of faces. Within each simplex, ordering the
barycentric coordinates of a point identifies a chain of
faces whose barycenters span a simplex containing the
point. Ties lie on the corresponding shared faces, so the
subdivisions agree on common faces. If \(F\subset G\) are
nonempty faces of a simplex of dimension at most \(n\),
the barycenter of \(G\) is the convex combination of the
barycenter of \(F\) with weight \(\#F/\#G\) and of the
barycenter of the complementary vertices. Thus the
distance between the two barycenters is at most
\(n/(n+1)\) times the old diameter. The mesh tends to
zero for the finite complex, since this bound holds for
all vertices of each new simplex.

For a vertex \(w\) in \(L\), its open star is the set of
points with positive \(w\)-coordinate. Such stars are
open in the weak topology, as seen on each closed
simplex. Their preimages cover the compact metric space
\(|K|\), and have a Lebesgue number. Choose a subdivision
fine enough that every open star of a new vertex \(v\)
has diameter smaller than this number. This follows from
the mesh bound, since points in the star are within one
mesh of \(v\). Choose \(g(v)\) with
\(f(\operatorname{st}(v))\subset\operatorname{st}(g(v))\).
If \(v_0,\ldots,v_j\) span a source simplex, the image of
any point of its relative interior belongs to all the
stars of \(g(v_i)\), so the latter vertices span a
simplex. This defines \(g\). More generally for each
\(x\), the support vertices of \(|g|(x)\) belong to
the carrier simplex of \(f(x)\), by applying the same
criterion to the support vertices of \(x\). Their straight
interpolation therefore stays in that carrier simplex.
The compact image of \(f\) is contained in a finite
subcomplex by Proposition~\ref{prop:p202-cw-basics}.
The chosen image vertices occur in this subcomplex,
and straight interpolation is continuous in its finite
barycentric realization. This proves the homotopy.
\end{proof}


\section{Relative cellular approximation}\label{p202:T12}

\begin{lemma}[Avoiding an interior point of a higher cell]
\label{lem:p202-avoid-cell}
Let \(f:D^n\to Z\) be a map, where a \(k\)-cell of the
Hausdorff space \(Z\) is open and its complement is \(W\).
If \(k>n\), there is a homotopy fixed on \(f^{-1}(W)\)
to a map missing some point in that cell.
\end{lemma}

\begin{proof}
Identify the open cell with \(\mathbb R^k\).
Choose concentric open balls \(B_0,B_1,B_2\) with
successive compact containments in the cell. The compact
set \(f^{-1}(\overline{B_1})\) lies in the open set
\(f^{-1}(B_2)\), because \(\overline{B_1}\) is compact
and \(Z\) is Hausdorff. View the source disk as a finite
polyhedron. A sufficiently fine triangulation gives a
finite subpolyhedron \(P\) lying in \(f^{-1}(B_2)\)
and containing \(f^{-1}(\overline{B_1})\) in its relative
interior. Indeed, the distance of the compact set to
the complement of \(f^{-1}(B_2)\) is positive when that
complement is nonempty. Take the union of closed simplices
meeting a smaller neighborhood of the compact set.
On \(P\) approximate the cell-coordinate map by its linear
interpolation at the vertices. Uniform continuity on a
compact neighborhood inside \(f^{-1}(B_2)\) makes the
approximation error smaller than the distance from
\(\overline B_0\) to the complement of \(B_1\).
Choose a continuous cutoff on \(D^n\) equal to one on
\(f^{-1}(\overline{B_1})\) and zero off the relative
interior of \(P\), using the distance to these two
closed sets. Interpolate in the cell coordinates
between \(f\) and the piecewise linear map using this
cutoff, and use \(f\) outside \(P\). The homotopy is
continuous by pasting, remains in the cell on \(P\),
and fixes \(f^{-1}(W)\).

Any point that the resulting map sends into \(B_0\)
must have been mapped by \(f\) into \(B_1\). There the
new map is piecewise linear. Its image on \(P\) lies
in a finite union of affine subsets of dimension at
most \(n\), none of which contains a nonempty open subset
of \(\mathbb R^k\). Thus some point of \(B_0\) is missed.
If the compact preimage used to choose \(P\) is empty,
\(f\) already misses \(B_0\).
\end{proof}

\begin{theorem}[Cellular approximation relative to a subcomplex]
\label{thm:p202-cellular-approx}
Let \(f:X\to Y\) be a map of CW complexes. If \(f\) is
cellular on a subcomplex \(A\), then it is homotopic
relative to \(A\) to a cellular map.
\end{theorem}

\begin{proof}
Assume the map already sends \(X^{n-1}\) into
\(Y^{n-1}\). On a characteristic disk of a source
\(n\)-cell outside \(A\), its image is contained in a
finite target subcomplex. In that finite subcomplex
choose a cell of maximal dimension \(k>n\). It is
open there. Lemma~\ref{lem:p202-avoid-cell} moves the
map on the disk to miss a point of this cell and
does not move the source boundary, since the boundary
already maps to \(Y^{n-1}\).
In the characteristic disk of the target cell, radial
projection away from the missed interior point gives
a deformation of the punctured disk onto its boundary
relative to the sphere. Passing to the adjunction,
and using the identity on its other cells, gives the
corresponding deformation of the punctured finite target
subcomplex. The product quotient property applies on
the saturated open complement of the missed point,
so this deformation is continuous in the adjunction
topology. Compose to remove this cell from the image.
Repeating through finitely many higher cells puts the
image into \(Y^n\). The construction also works for
\(n=0\); its source boundary is empty.

For the fixed \(n\), carry out this construction on
every \(n\)-cell outside \(A\), and use the constant
homotopy on \(A\cup X^{n-1}\). The cell homotopies
paste on \(A\cup X^n\). Extend the homotopy to \(X\)
by Theorem~\ref{thm:p202-cw-hep} for this subcomplex.
Perform the successive dimensions on time intervals
whose right endpoints tend to one. On each
characteristic disk the restriction is eventually
constant, so the resulting homotopy is continuous at
the terminal time on every such disk. The weak-topology
test on products gives continuity on \(X\times I\).
The final map is cellular and the homotopy fixes \(A\).
\end{proof}

In particular, maps from a CW complex of dimension \(n\)
can be moved into \(Y^n\), and homotopies between maps
already cellular on the ends can be taken in the
\((n+1)\)-skeleton when the source cylinder has that
dimension. Apply the theorem to the cylinder relative
to its two ends, using the cells of \(I\) with both
endpoints as vertices. This gives the skeletal control
of based homotopy classes used below. Compare
\cite[Section 4.1]{Hatcher}.

\section{Whitehead's theorem}\label{p202:T13}

For \(n\ge1\), \(\pi_n(Y,y)\) uses maps
\((I^n,\partial I^n)\to(Y,y)\), with homotopies fixing
the boundary. Concatenation in one cube coordinate
gives the group law, using the same reparametrization
and reversal as for paths. For change of basepoint along
a path \(d\) from \(y_0\) to \(y_1\), put a rescaled map
based at \(y_0\) on a smaller concentric cube. On the
surrounding collar run \(d\) from its inner to its outer
boundary, using distance from the cube center in the
maximum norm. This produces a map based at \(y_1\).
The reversed path gives the inverse by contracting the
out-and-back collar and rescaling. If a homotopy moves
its basepoint along \(d\), use the same construction at
each time with the remaining segment of \(d\) in the
collar. This makes the whole homotopy based at \(y_1\)
and proves the usual change-of-basepoint identity.
For a pair \((Y,X)\) the relative disk condition permits
the boundary to map
into \(X\). In dimension one, a relative class is a
path starting at a specified basepoint and ending
somewhere in \(X\), with homotopies allowing the
terminal point to move in \(X\). We use its pointed
set interpretation.

\begin{lemma}[Disk replacement]
\label{lem:p202-whitehead-disk}
Suppose \(X\subset Y\) induces a bijection on path
components and isomorphisms on all \(\pi_n\), at
all basepoints in \(X\). Every map \(f:D^n\to Y\)
with boundary in \(X\) is homotopic relative to
the boundary to a map into \(X\).
\end{lemma}

\begin{proof}
For \(n=0\) use surjectivity on path components.
For \(n=1\), the endpoints lie in the same path
component of \(X\) by injectivity on path components,
so choose a path \(g\) between them in \(X\).
The class of \(f\bar g\) in \(Y\) is the image
of a loop \(\ell\) in \(X\) at the initial point.
Thus \(f\) is homotopic relative to the endpoints
to \(\ell g\).

For \(n\ge2\), the boundary of \(f\) is null as a
based sphere in \(Y\), since \(f\) fills it. The
based assertion follows from any free contraction
by following the path of its basepoint, which
transports the zero class to zero. Injectivity on
\(\pi_{n-1}\) implies there is a filling \(g\) in
\(X\) with the same boundary. Gluing \(f\) and
\(g\) on the two oppositely oriented faces of a
sphere gives a based class in \(\pi_n(Y)\). It
comes from a sphere in \(X\) by surjectivity.
Insert a representative of this latter sphere
in the interior of the filling \(g\), with the
sign that cancels the difference. To make this
explicit, make \(g\) constant on a small inner
ball by a radial collapse of that ball to its
center, and use the inserted sphere on the ball
with its boundary collapsed. Transport the chosen
sphere basepoint along the image under \(g\) of
a path from the boundary basepoint to the center.
The gluing sphere then represents the old difference
minus the inserted class: pinch the separating
sphere to a point to read the two terms as the
two concatenated cube maps. Denote the new
filling by \(g'\). The glued sphere for \(f,g'\)
is null. On \(\partial(D^n\times I)\), put \(f\)
and \(g'\) on the two disk faces and the common
boundary map, constant in time, on the side.
This sphere map represents the same glued class,
since the side can be collapsed in the time
coordinate. Its null homotopy extends it over
the ball \(D^n\times I\), giving the required
homotopy relative to the boundary.
\end{proof}

\begin{theorem}[Whitehead for CW complexes]
\label{thm:p202-whitehead}
A map of CW complexes that induces a bijection on
path components and isomorphisms on all based
homotopy groups in positive dimensions is a
homotopy equivalence.
\end{theorem}

\begin{proof}
First assume the map is the inclusion of a
subcomplex \(X\subset Y\). Deform the identity
of \(Y\) into \(X\) one skeleton at a time,
fixing \(X\). At the \(n\)-th step, the
current map on every \(n\)-cell outside \(X\)
has boundary in \(X\). Lemma~\ref{lem:p202-whitehead-disk}
moves it into \(X\) relative to the boundary.
Paste on \(X\cup Y^n\), then extend to \(Y\)
using homotopy extension. The time intervals
accumulate at one as in the proof of cellular
approximation, and every characteristic disk is
eventually fixed. Thus the homotopy is continuous
and ends at a retraction. This gives a strong
deformation retract.

For a general map \(f:X\to Y\), first replace it
by a homotopic cellular map. Unbased homotopy
transports the homotopy groups along the paths of
basepoints, so the hypotheses are preserved.
Form its mapping cylinder by identifying
\((x,1)\) with \(f(x)\) in \(X\times I\sqcup Y\).
It is a CW complex with both \(X\times\{0\}\)
and \(Y\) subcomplexes: in addition to their
cells, use one cylinder-interior cell of
dimension \(n+1\) for each \(n\)-cell of \(X\).
Its bottom and top attach to the \(n\)-skeleta
because \(f\) is cellular, and its side lies in
the lower cylinder cells. The product quotient
test gives the same weak topology as the
mapping-cylinder adjunction. Sliding upward in
the cylinder deformation retracts it onto \(Y\),
and the restriction to \(X\times\{0\}\) is \(f\).
The inclusion of this copy of \(X\) therefore
satisfies the hypotheses of the inclusion case.
Composing its homotopy equivalence with the
cylinder retraction proves the result.
\end{proof}


\chapter{Singular Chains, Excision, and Gluing}

We take \(R\) to be a commutative ring with identity unless
an abelian coefficient group is explicitly indicated. All
chains are finite sums. In particular, no local-finiteness
hypothesis on the space is needed for ordinary singular
homology.

\section{Singular chains and degree zero}\label{p202:T14}

\begin{definition}[Singular and augmented chains]
For \(n\ge0\), the standard simplex \(\Delta^n\) has ordered
vertices \(v_0,\ldots,v_n\). A singular \(n\)-simplex is a
continuous map \(\sigma:\Delta^n\to X\). Let \(C_n(X;R)\)
be the free \(R\)-module with this basis. Let
\(\delta_i:\Delta^{n-1}\to\Delta^n\) be the affine
inclusion that omits \(v_i\). For \(n>0\), put
\[
 \partial\sigma=\sum_{i=0}^n(-1)^i\sigma\delta_i ,
\]
and let \(\partial:C_0\to 0\) be zero. The augmentation
\(\epsilon:C_0\to R\) sends each point simplex to \(1\).
We use ordinary homology \(H_n(X;R)\) of \(C_*(X;R)\)
and reduced homology \(\widetilde H_n(X;R)\) of the
augmented complex with \(R\) in degree \(-1\).
With this convention \(\widetilde H_{-1}(\emptyset;R)=R\).
\end{definition}

\begin{proposition}[The singular chain functor]
\label{prop:p202-singular-functor}
The boundary squares to zero, and postcomposition defines
a functor of spaces to chain complexes. For an abelian
group \(A\), coefficients in \(A\) are obtained by
\(C_*(X;\mathbb Z)\otimes_{\mathbb Z} A\), with the same
boundary formula. For \(R\)-module coefficients \(M\),
one can equivalently form \(C_*(X;R)\otimes_R M\).
The zero-dimensional group is the direct sum of the
coefficient group over path components. Reduced degree
zero is the kernel of the sum map from this direct sum
to the coefficient group.
\end{proposition}

\begin{proof}
For two deleted vertices \(i<j\), the two orders of deletion
give the same face by the affine vertex formulas and
signs \((-1)^{j+i}\) and \((-1)^{i+j-1}\). They cancel
in the double boundary. The two endpoints of an edge have
opposite signs, so the augmentation kills boundaries.
Postcomposition commutes with every face map and with
the augmentation. Tensoring a free group or module on
the simplex basis gives the asserted chain groups and
boundaries with coefficients.

In degree zero, each path gives the relation that its
endpoints represent the same class. These are exactly the
relations needed to identify all the generators in each
path component. More explicitly, send a point to the
basis element of its component. To invert the induced
map on the quotient, choose a point in each component
and take its class. This class is independent of the
choice, and the composites are identities. The augmentation
is the sum on the direct sum, so taking its kernel gives
the reduced assertion. The empty space has no simplex
basis and the conventions follow directly from the
augmentation.
\end{proof}

The point has a single simplex in each nonnegative degree.
Its differential in positive degree is \(1\) in even degrees
and \(0\) in odd degrees, by the alternating face sum. Hence
its ordinary homology is \(R\) in degree zero and zero in
positive degrees, and its reduced homology is zero.
For \(S^0\) the two points give two such complexes,
so its reduced \(H_0\) is \(R\) and all its positive groups
vanish. Compare \cite[Section 2.1]{Hatcher}
and \cite[Sections 4.2 and 6.3]{MKP}.

\section{The prism operator}\label{p202:T15}

\begin{theorem}[Homotopy invariance]
\label{thm:p202-prism}
A homotopy from \(f\) to \(g\) induces a chain homotopy
from \(f_\#\) to \(g_\#\). If the homotopy is a homotopy
of pairs, it induces the same identity on relative
chains. In particular, based or unbased homotopic maps
induce the same ordinary and reduced homology maps,
and a nonempty contractible space has zero reduced homology.
\end{theorem}

\begin{proof}
Write \(v_i^0=(v_i,0)\) and \(v_i^1=(v_i,1)\).
For a singular \(n\)-simplex \(\sigma\) and a homotopy
\(F\), let \(P(\sigma)\) be the image by
\(F\circ(\sigma\times\mathrm{id})\) of the signed sum
of affine simplices in the prism with vertex lists
\[
 \sum_{i=0}^{n}(-1)^i
 [v_0^0,\ldots,v_i^0,v_i^1,\ldots,v_n^1].
\]
The top face appears when the first bottom vertex in the
\(i=0\) term is deleted, with positive sign. The bottom
face appears when the last top vertex in the \(i=n\) term
is deleted, with negative sign. For \(i>0\), deleting
\(v_i^0\) in term \(i\) gives the same face as deleting
\(v_{i-1}^1\) in term \(i-1\), with opposite signs.
The other faces come from the prisms of the faces of
\(\sigma\). Deleting \(v_j^0\) for \(j<i\) has sign
\((-1)^{i+j}\) and corresponds to term \(i-1\) of the
\(j\)-th face prism. Deleting \(v_j^1\) for \(j>i\)
has sign \((-1)^{i+j+1}\) and corresponds to term \(i\)
of that face prism. Both signs are the negatives of the
signs in \(P\partial\). Thus
\(\partial P+P\partial=g_\#-f_\#\).
In reduced degree \(-1\) the two maps are the identity
of \(R\), and we extend \(P\) by zero there. A homotopy
preserving a specified subspace sends each prism of
a simplex in that subspace into the target subspace, so
the identity descends to the relative quotient. Applying
homotopy invariance to the homotopy equivalence with
a point proves the contractible case.
\end{proof}

\section{Relative chains and exactness}\label{p202:T16}

\begin{definition}[Relative complex]
For a pair \((X,A)\), put
\(C_*(X,A;R)=C_*(X;R)/C_*(A;R)\), with induced boundary.
The inclusion is injective on simplex bases and the
quotient has a basis of the remaining simplices.
Its homology is \(H_*(X,A;R)\).
\end{definition}

\begin{lemma}[The connecting map for a short exact sequence]
\label{lem:p202-chain-les}
A degreewise short exact sequence
\(0\to K\xrightarrow{i}L\xrightarrow{p}M\to0\)
of chain complexes has a natural long exact homology
sequence. For a cycle \(z\in M_n\), choose a lift
\(c\in L_n\) and define the connecting map by
\([z]\mapsto [i^{-1}\partial c]\in H_{n-1}(K)\).
\end{lemma}

\begin{proof}
The boundary of a lift lies in the image of \(i\)
and is a cycle there because \(i\) is injective and
\(\partial^2=0\). Changing the lift changes this
cycle by a boundary in \(K\); changing \(z\) by a
boundary is handled by adding the boundary of a lift
of its primitive to \(c\). Hence the formula descends.
If \([c]\in H_n(L)\) maps to zero, subtract the
boundary of a lift of a primitive for \(p(c)\).
The resulting cycle belongs to \(i(K)\). A cycle
\(z\in M\) has zero connecting class exactly when
its lift can be corrected by an element of \(i(K)\)
to become a cycle. Finally a cycle in \(K\) maps
to a boundary in \(L\) exactly when it is the boundary
of a chain lifting a cycle in \(M\). This gives all
three exactness checks. Applying a map of short exact
sequences to any chosen lift gives the naturality.
\end{proof}

\begin{theorem}[The long exact sequence of a pair]
\label{thm:p202-relative-les}
For any pair \((X,A)\) there is a natural exact
sequence
\[
 \cdots\longrightarrow H_n(A;R)\longrightarrow H_n(X;R)
 \longrightarrow H_n(X,A;R)
 \xrightarrow{\partial}H_{n-1}(A;R)\longrightarrow\cdots .
\]
The connecting map sends the class of a relative cycle
represented by \(c\) to the class of its ordinary
boundary in \(A\).
\end{theorem}

\begin{proof}
Apply Lemma~\ref{lem:p202-chain-les} to the short exact
sequence defining relative chains. This also proves the
reduced version when using the same augmentation term
on \(A\) and \(X\). In degree \(-1\) the induced
relative complex is zero, and in every nonnegative
degree it remains the displayed relative complex.
\end{proof}

\section{Subdivision, small chains, and excision}\label{p202:T17}

For an open cover \(\mathcal U\) of \(X\), write
\(C_*^{\mathcal U}(X;R)\) for the subcomplex spanned by
simplices with image in a single member of the cover.
For a subspace \(A\) use the restricted cover, discarding
empty intersections. The support of a chain is the union
of the images of its finitely many simplices with nonzero
coefficient.

\begin{lemma}[Support-controlled subdivision]
\label{lem:p202-subdivision}
Barycentric subdivision defines a natural chain map
\(S:C_*(X;R)\to C_*(X;R)\) chain homotopic to the identity.
The chain homotopy \(K\) and \(S\) do not enlarge the
support of a simplex.
\end{lemma}

\begin{proof}
In \(\Delta^n\) take the signed sum of the affine maps
of the full-dimensional simplices of its barycentric
subdivision. Order their vertices by increasing face
dimension in the flag, with each sign given by agreement
with the orientation of \(\Delta^n\). Define \(S(\sigma)\) by
postcomposing this sum with \(\sigma\).
Every internal codimension-one face occurs twice with
opposite induced orientations. A boundary face has
precisely its own induced barycentric subdivision and
inherits the sign of the original ordered face.
These facts give \(\partial S=S\partial\).
They can be read directly from the flags of faces:
an internal missing rank in a flag has two possible
insertions, whereas a flag in an original face has one
extension by the top simplex.

Construct \(K\) inductively using chains inside
\(\Delta^n\). In degree zero take \(K=0\), since
subdivision fixes a vertex. If \(K\) has already been
defined in lower degrees, write \(\iota_n\) for the
identity simplex of \(\Delta^n\). The affine chain
\[
 z=S(\iota_n)-\iota_n-K\partial\iota_n
\]
has zero boundary by the inductive identity. Cone each
affine simplex in \(z\) to the barycenter \(b\), inserting
\(b\) as the first vertex. The boundary of the resulting
cone is \(z-b*\partial z=z\), by the alternating face
formula. For degree zero the cone identity uses the
augmentation. Put \(K(\iota_n)=b*z\), and define
\(K(\sigma)\) by postcomposition. This defines a natural
operator with \(\partial K+K\partial=S-\mathrm{id}\).
All affine maps used before postcomposition have image in
the original standard simplex. This proves support control.
\end{proof}

\begin{theorem}[The small-chain comparison]
\label{thm:p202-small-chains}
The inclusion of \(\mathcal U\)-small chains into ordinary
chains induces a homology isomorphism. It is a chain
homotopy equivalence over \(R\). For any subspace \(A\),
the corresponding inclusion of small relative chains
into \(C_*(X,A;R)\) has the same properties.
\end{theorem}

\begin{proof}
For a singular simplex, pull back the cover to its compact
domain and use a Lebesgue number. The mesh estimate from
Section~\ref{p202:T11} shows that sufficiently many
subdivisions make each affine sub-simplex small. This
applies to any finite chain by taking the maximum of
finitely many numbers of subdivisions. The telescoping
chain homotopy between \(S^m\) and the identity is
\(\sum_{j=0}^{m-1} KS^j\), and it preserves smallness
whenever the input is already small. Thus every ordinary
cycle is homologous to a small cycle. If a small cycle
bounds \(c\) in ordinary chains, take \(m\) so \(S^m c\)
is small. The small cycle is homologous within small
chains to its \(m\)-fold subdivision, which now bounds
there. This gives injectivity. The same argument in
the relative quotient works because \(S\) and \(K\)
preserve chains in \(A\).

For the stronger chain assertion, in either case let
\(Q\) be the quotient of the large chain complex by
the small subcomplex. The inclusion is of a subset
of the given simplex basis, also after removing
simplices in \(A\). Thus \(Q\) is degreewise free.
It is acyclic by the just-proved isomorphism and
Lemma~\ref{lem:p202-chain-les}. A bounded-below
acyclic complex of free modules admits a contraction
\(h\): starting at degree zero, split each surjection
\(Q_n\to Z_{n-1}(Q)\), since \(Z_{-1}=0\) and the
previous splitting makes the next cycle module
projective. In the decomposition
\(Q_n=Z_n(Q)\oplus Z_{n-1}(Q)\), send the first
summand to its chosen lift in \(Q_{n+1}\) and send
the second to zero. This gives
\(\partial h+h\partial=\mathrm{id}\).
If \(s\) is a graded section of the quotient map,
\(s'=\partial s h+s h\partial\) is a chain section:
its composite with the quotient map is the identity,
and the two chain-map expressions both equal
\(\partial s h\partial\).
Hence the large complex splits as the small complex
plus a contractible summand, proving the chain homotopy
equivalence. Choices in this argument do not affect the
naturality of the inclusion on homology.
\end{proof}

\begin{theorem}[Excision]
\label{thm:p202-excision}
If \(Z\subset A\subset X\) and
\(\overline{Z}\subset\operatorname{int}_X A\), the
inclusion of \((X\setminus Z,A\setminus Z)\) into
\((X,A)\) induces a homology isomorphism.
\end{theorem}

\begin{proof}
Use the open cover \(U=X\setminus\overline Z\) and
\(V=\operatorname{int}_X A\) of \(X\). Since
\(V\subset A\), the small relative complex is naturally
\(C_*(U;R)/C_*(U\cap A;R)\). In fact every small
simplex in \(V\) is zero in the relative complex, and
the only generators in \(U\) that become zero are those
with image in \(U\cap A\). On \(X\setminus Z\) use
the restricted open cover by \(U\) and
\(V\cap(X\setminus Z)\). The same calculation gives
the same quotient, and inclusion is the identity on
it. The small-chain comparison now proves the result.
\end{proof}

\begin{corollary}[The quotient of a good pair]
\label{thm:p202-good-pair}
If \(A\ne\emptyset\) is closed in \(X\) and is a strong
deformation retract of an open neighborhood \(W\), then
the quotient projection induces an isomorphism
\[
 H_n(X,A;R)\ \cong\ \widetilde H_n(X/A;R)
 \qquad (n\ge0).
\]
\end{corollary}

\begin{proof}
Write \(q:X\to X/A\) and \(*=q(A)\).
The maps from \((X,A)\) to \((X,W)\) and from
\((X/A,\{*\})\) to \((X/A,W/A)\) induce
isomorphisms by the long exact sequences and homotopy
invariance. For the second map, \(W/A\) is contractible,
since the retraction homotopy on \(W\) fixes \(A\)
and descends by the product quotient property.
Now excise \(A\) from \((X,W)\), and \(\{*\}\)
from \((X/A,W/A)\). The closed sets being excised
are contained in the respective open second terms.
The resulting pairs are homeomorphic under \(q\):
off \(A\) this quotient is the identity identification
of an open saturated complement. Thus the quotient
map is an isomorphism on relative homology.
The reduced pair sequence for \((X/A,\{*\})\) identifies
its relative homology with \(\widetilde H_n(X/A;R)\)
also in degree zero. For \(A=\emptyset\) the relative
complex is the ordinary complex of \(X\); this case
uses the usual ordinary/reduced distinction rather than
the nonempty-subspace formula.
\end{proof}

\section{Mayer--Vietoris and the ordinary axioms}\label{p202:T18}

\begin{theorem}[Open-cover Mayer--Vietoris]
\label{thm:p202-mv}
If \(X=U\cup V\) with \(U,V\) open, there is a natural
long exact sequence
\[
 \cdots\to H_n(U\cap V;R)\to H_n(U;R)\oplus H_n(V;R)
 \to H_n(X;R)\to H_{n-1}(U\cap V;R)\to\cdots .
\]
For a subspace \(A\subset X\), the relative version
uses pairs with second terms \(A\cap U\cap V\),
\(A\cap U\), \(A\cap V\), and \(A\), respectively.
\end{theorem}

\begin{proof}
There is a short exact sequence of chain complexes
\[
 0\to C_*(U\cap V;R)
 \xrightarrow{(i,-j)} C_*(U;R)\oplus C_*(V;R)
 \xrightarrow{+} C_*^{\{U,V\}}(X;R)\to0 .
\]
Indeed the target basis consists precisely of the
union of the two bases, and duplicate entries are
exactly the simplices with image in the intersection.
Apply Lemma~\ref{lem:p202-chain-les} and then use
Theorem~\ref{thm:p202-small-chains}. If the small
cycle decomposes as \(a+b\), with \(a\) in \(U\) and
\(b\) in \(V\), its connecting class is
\([\partial a]=[-\partial b]\) in the intersection.
This is the sign determined by \((i,-j)\).
For relative chains, remove the bases of simplices
lying in \(A\) throughout the short exact sequence.
The same basis argument applies, and the small
relative comparison finishes the proof. All these maps
are natural for maps sending the two cover members
to their respective cover members and respecting the
second terms of the pairs.
\end{proof}

Ordinary singular homology satisfies functoriality,
homotopy invariance, the exact sequence of a pair,
and excision by the preceding constructions. The
dimension axiom follows from the direct point
calculation in Section~\ref{p202:T14}.
For disjoint unions a simplex has connected domain
and its image lies in one summand. Thus the chain
complex is the direct sum of the summand complexes,
which gives the usual additivity statement. The
constructions in this chapter are for ordinary homology
with the specified coefficients. They provide the
relative and small-chain interfaces used in the
comparisons below. See \cite[Sections 2.1--2.3]{Hatcher}
and \cite[Sections 3.2--3.6 and 4.2--4.4]{MKP}.
\chapter{Chain Models and Cohomological Operations}

Unless a coefficient is specified, a chain complex in this chapter has
coefficients in a commutative ring \(R\) with unit. Integral chains are
the case \(R=\mathbb Z\). A change to coefficients in an \(R\)-module
\(A\) uses tensor products for chains and \(\operatorname{Hom}_R(-,A)\)
for cochains. We use \(H_{-1}=0\) for ordinary homology; in reduced
formulas involving degree zero we use the augmented complex. An
oriented standard simplex has vertex order \(0,\ldots,n\). The
boundary of a positive-dimensional simplex is its signed sum of
ordered faces, and the ordinary boundary of a vertex is zero.

\section{Simplicial chains and comparison}\label{p202:T19}

For an abstract simplicial complex, a generator \([v_0,\ldots,v_q]\)
is a simplex with an orientation, with permutations acting by sign.
The affine face boundary is compatible with this relation. A
\(\Delta\)-complex instead has specified parametrizations of its cells
whose faces are identified by order-preserving affine maps. In that
case distinct faces of a cell can represent the same cell of the
complex. Its chain differential is the signed sum counting each face
with its multiplicity. Both boundary squares vanish, since deleting
two indices in either order yields opposite signs.

\begin{theorem}[Simplicial--singular comparison]\label{thm:p202-simplicial-comparison}
For a simplicial complex \(K\) with its weak realization topology,
the map sending each oriented simplex, written in a chosen increasing
vertex order, to its affine singular parametrization induces
\[
 H_q^{\mathrm{simp}}(K;A)\simeq H_q(|K|;A)
\]
for every \(q\geq0\). This is natural for simplicial maps. The
corresponding parametrization map for a \(\Delta\)-complex has the
same conclusion, naturally for maps preserving the ordered simplex
parametrizations. These comparisons also apply to subcomplex pairs.
\end{theorem}

\begin{proof}
The increasing representatives are a basis, and each face inherits
the increasing order. Hence this is a chain map, and the same face
identity works with specified \(\Delta\)-parametrizations.

We recall the elementary relative calculation used in the skeletal
proof. The long exact sequence of a disk and its boundary, with a
contractible disk and \(H_0\) computed by path components, identifies
\[
 H_j(D^n,S^{n-1};A)=
 \begin{cases} A & j=n,\\0&j\ne n.\end{cases}
\]
Indeed the sphere computation follows by induction from the
two-hemisphere Mayer--Vietoris sequence, starting with the two
components of \(S^0\). For \(n=0\) use the pair of a point and the
empty space. In the simplex model the integral relative class of
the affine \(n\)-simplex is a generator, and its multiples by
elements of \(A\) give the claimed copy of \(A\).
Its boundary is the signed boundary
cycle on the simplex boundary. To check the sign and primitivity
inductively, split that boundary into the face opposite \(v_0\) and
the cone from \(v_0\) over the boundary of that face. After using
collars to thicken the pieces, the Mayer--Vietoris connecting map
sends the signed cycle to the oriented boundary of the opposite
face, a generator by induction. In dimension one the signed
boundary is the difference of the two endpoints. This proves the
relative calculation with the indicated generator.

The \(n\)-skeleton modulo the \((n-1)\)-skeleton is a weak wedge of
\(n\)-spheres when the lower skeleton is nonempty and \(n>0\).
The pair is a good pair by Section~\ref{p202:T10}. Excision and the
long exact sequence, or an excisive cover of a finite wedge followed
by finite support, identify its relative homology as the direct sum
of the preceding relative groups. The zero skeleton is discrete
and gives the same direct-sum description directly. The
simplicial relative chain complex of consecutive skeleta is
concentrated in degree \(n\). The parametrization map sends its
basis to the just-identified relative generators. Starting with
the zero skeleton, compare the long exact sequences of the
skeletal pairs and apply exactness and the five lemma in each
degree. This gives comparison on every skeleton and on every
finite subcomplex. A finite singular chain has compact image,
hence lies in a finite subcomplex by Section~\ref{p202:T09};
the same holds for any proposed bounding chain. Thus the finite
subcomplex comparison passes to the whole complex. Comparison for
subcomplex pairs then follows from the exact sequences of the pairs.

For an arbitrary simplicial map, the simplicial image of a
simplex is its signed oriented image, or zero when vertices
collapse. The resulting square of chain maps need not commute on
the nose with the parametrization maps, but both routes on a
simplex have support in its image simplex. They agree on vertices.
Inductively, after a homotopy \(h\) has been defined on all lower
faces, the difference of the two images minus \(h\partial\) is a
cycle in that image simplex. It bounds there by the cone
contraction to a vertex. Filling it defines \(h\) on the simplex.
This gives the required chain homotopy and works in a subcomplex
pair because the carrier of a simplex of the subcomplex stays
there. The same argument compares two vertex orders, inserting the
orientation sign before applying the second parametrization.
For maps preserving \(\Delta\)-parametrizations the square
commutes directly.
\end{proof}

See \cite[Section 2.1]{Hatcher} and \cite[Section 4.7]{MKP}.

\section{Cellular chains and attaching degrees}\label{p202:T20}

\begin{theorem}[Cellular homology]\label{thm:p202-cellular-comparison}
For a CW complex \(X\) put \(X^{-1}=\varnothing\), and set
\[
 C_n^{\mathrm{cell}}(X;A)=H_n(X^n,X^{n-1};A).
\]
With \(d_0=0\), define \(d_n\) for \(n>0\) by the connecting map
to \(H_{n-1}(X^{n-1};A)\) followed by the relative projection
to \(C_{n-1}^{\mathrm{cell}}(X;A)\). These are chain groups with
one copy of \(A\) per oriented \(n\)-cell, using a direct sum,
and their homology is naturally \(H_n(X;A)\) for cellular maps.
For \(n>1\), the coefficient from an \(n\)-cell to an
\((n-1)\)-cell is the degree of the attaching map
\(S^{n-1}\to X^{n-1}\) followed by collapsing the other cells and
\(X^{n-2}\) to the basepoint of the selected sphere. For \(n=1\)
the coefficient is the signed incidence of the two endpoints.
\end{theorem}

\begin{proof}
The relative disk computation in Section~\ref{p202:T19}, with
the good skeletal pair and its wedge quotient, identifies all
\(H_j(X^n,X^{n-1};A)\), with the only possible nonzero group
in degree \(n\). The projection to a wedge summand gives its
coordinate. Naturality of the connecting homomorphism for a
characteristic map \((D^n,S^{n-1})\to (X^n,X^{n-1})\) sends
the relative generator to the attaching sphere class. Projecting
to the selected coordinate gives precisely its degree. An
oriented interval has endpoint boundary final minus initial.
Compactness of the attaching sphere and the finite-subcomplex
property give finite support even for an infinite CW complex.

The pair exact sequence gives \(d_{n-1}d_n=0\), since a class
coming from absolute homology has zero connecting image. Induction
on \(m\) with the same exact sequence yields \(H_j(X^m;A)=0\)
for \(j>m\). In particular
\(H_n(X^n;A)\to C_n^{\mathrm{cell}}(X;A)\) is injective.
For \(n>0\), the map
\(H_{n-1}(X^{n-1};A)\to C_{n-1}^{\mathrm{cell}}(X;A)\) is
injective by the same vanishing, so the image of
\(H_n(X^n;A)\) is \(\ker d_n\); for \(n=0\) this is immediate.
The pair \((X^{n+1},X^n)\) gives
\[
 H_n(X^{n+1};A)=
 H_n(X^n;A)/\operatorname{im}\bigl(H_{n+1}(X^{n+1},X^n;A)
               \longrightarrow H_n(X^n;A)\bigr).
\]
Consequently this is \(\ker d_n/\operatorname{im}d_{n+1}\).
Higher cell attachments preserve \(H_n\) by the relative
vanishing, and passage to the union is justified by the
finite-subcomplex property on cycles and bounding chains.
All maps used in this identification arise from the skeletal
pair sequences and commute with cellular maps.
\end{proof}

\section{Representative homology computations}\label{p202:T21}

\begin{proposition}[Spheres, wedges, and standard finite CW models]\label{prop:p202-homology-computations}
For \(n\geq0\), \(\widetilde H_j(S^n;A)\) is \(A\) if \(j=n\)
and zero otherwise in nonnegative degrees. A point or any
nonempty contractible space has zero reduced homology. A finite
wedge of based CW complexes at vertices, and more generally a
finite wedge of good pointed pairs with the usual wedge topology,
has reduced homology the direct sum of the reduced homologies of
the summands. With integral coefficients, for \(m\geq1\),
\[
\begin{aligned}
 H_j(T^m;\mathbb Z)&\simeq \mathbb Z^{\binom m j},\\
 H_j(\mathbb{CP}^m;\mathbb Z)&\simeq
   \begin{cases}\mathbb Z & j=0,2,\ldots,2m,\\0&\text{otherwise},\end{cases}\\
 H_j(\mathbb{RP}^m;\mathbb Z)&\simeq
   \begin{cases}
      \mathbb Z & j=0,\ \text{or }j=m\text{ odd},\\
      \mathbb Z/2 & 0<j<m\text{ odd},\\
      0&\text{otherwise}.
   \end{cases}
\end{aligned}
\]
A closed oriented genus-\(g\) surface in its standard polygon
model has \(H_0,H_2\simeq\mathbb Z\) and
\(H_1\simeq\mathbb Z^{2g}\). The standard nonorientable genus-\(k\)
model, \(k\geq1\), has \(H_0\simeq\mathbb Z\),
\(H_1\simeq\mathbb Z^{k-1}\oplus\mathbb Z/2\), and \(H_2=0\).
All the listed spaces have zero homology above their cell dimension.
\end{proposition}

\begin{proof}
Homotopy invariance gives the contractible calculation, and the
sphere computation was proved in Section~\ref{p202:T19}. For a
wedge of two good pointed pairs, enlarge each summand by a
contractible-to-the-basepoint neighborhood from the other.
The resulting open cover retracts onto the two summands and
its intersection retracts to the basepoint. Apply reduced
Mayer--Vietoris and then induct.

In the product cubical decomposition of \(T^m=(S^1)^m\)
there is one cell for each choice of active circle coordinates.
The opposite faces in a coordinate have the same image and
opposite signs, so every differential vanishes. For
\(\mathbb{CP}^m\), use its nested projective subspaces. In
\(\mathbb{CP}^j\setminus\mathbb{CP}^{j-1}\) the last nonzero
homogeneous coordinate can be set to \(1\), giving \(\mathbb C^j\).
A closed ball in \(\mathbb C^j\) parametrizes the cell by
\((z,\sqrt{1-\lvert z\rvert^2})\) in homogeneous coordinates.
The boundary lies in the preceding subspace. Thus there is
one cell in each even dimension and none in odd dimensions.

The analogous real model gives one cell in every dimension.
The attaching sphere for the \(j\)-cell maps by the antipodal
quotient \(S^{j-1}\to\mathbb{RP}^{j-1}\). The two hemispheres
cover the open \((j-1)\)-cell once each; their oriented
contributions differ by the degree \((-1)^j\) of the antipodal
map on \(S^{j-1}\). Thus the cellular coefficient is
\(1+(-1)^j\), with the \(j=1\) value also given by the coincident
endpoints. For the degree sign, a single coordinate
reflection interchanges the two thickened hemispheres and
fixes their equator. The reduced Mayer--Vietoris boundary
isomorphism changes sign when the two open sets are
interchanged, by the difference convention of the small-chain
sequence. Thus the reflection has degree \(-1\) in positive
sphere dimension, and the assertion in dimension zero is
the interchange of the two generators. The antipodal map
on \(S^{j-1}\) is the composite of \(j\) such reflections.
Taking homology gives the displayed groups.

For the oriented polygon the attaching word of the sole
two-cell is a product of commutators in \(2g\) edges, so every
signed degree is zero. For the nonorientable polygon the word
is \(a_1^2\cdots a_k^2\), so the column of \(d_2\) is
\(2(1,\ldots,1)\). Integer row operations send it to
\((2,0,\ldots,0)\), giving the result. There is one vertex
in both models.
\end{proof}

\begin{corollary}[Euler--Poincar\'e]\label{cor:p202-euler}
For a finite CW complex \(X\) and any field \(F\),
\[
 \chi(X)=\sum_q(-1)^q \#\{\text{\(q\)-cells}\}
       =\sum_q(-1)^q\dim_F H_q(X;F).
\]
\end{corollary}
\begin{proof}
Each finite-dimensional cellular chain group has dimension
\(\dim H_q+\dim B_q+\dim B_{q-1}\), by the cycle and
boundary exact sequences. The boundary dimensions cancel
in the alternating sum. In particular \(\chi\) is independent
of the chosen finite CW structure.
\end{proof}

\section{Degree and the first homology group}\label{p202:T22}

\begin{proposition}[Degree and no retraction]\label{prop:p202-degree}
For a continuous map \(f:S^n\to S^n\), its degree is the integer
by which it acts on \(\widetilde H_n(S^n;\mathbb Z)\simeq\mathbb Z\).
It is homotopy invariant and multiplicative under composition.
There is no retraction \(D^{n+1}\to S^n\), and every continuous
self-map of \(D^{n+1}\) has a fixed point.
\end{proposition}
\begin{proof}
The degree assertions follow from homotopy invariance and
functoriality. For \(n=0\) the generator is the difference
of the vertices, so a self-map of \(S^0\) has degree \(1\),
\(-1\), or \(0\) according to whether it is the identity,
the interchange, or constant. A retraction would make the
identity of the reduced sphere group factor through the
zero reduced group of the disk, which is impossible. If
a self-map \(F\) of the closed ball has no fixed point,
the ray from \(F(x)\) through \(x\) meets the unit sphere at
its forward exit point. This defines a continuous function:
put \(v=x-F(x)\ne0\), and take the larger root of
\(\lvert F(x)+tv\rvert^2=1\). It is at least \(1\) and
equals \(1\) for \(x\) on the sphere, since
\(\langle x,x-F(x)\rangle>0\) there when \(F(x)\ne x\).
The root depends continuously on \(x\) by the quadratic
formula. This would give the prohibited retraction.
\end{proof}

\begin{theorem}[Degree-one Hurewicz map]\label{thm:p202-hurewicz-one}
If \(X\) is path connected, the map taking a based loop to
its singular one-cycle induces
\(\pi_1(X)_{\mathrm{ab}}\simeq H_1(X;\mathbb Z)\).
\end{theorem}
\begin{proof}
Endpoint-fixed homotopy gives homologous chains by the prism
identity. For paths \(a,b\) with matching endpoints, the
difference \(a+b-(a*b)\) is the boundary of a singular
two-simplex. To obtain it, map the edges of the triangle to
\(a,b,a*b\) in their indicated order. This boundary loop
retraces the concatenation and is nullhomotopic, hence
extends over the triangle. A constant path is the boundary
of the constant two-simplex. Thus the map from loops is a
homomorphism into an abelian group.

Fix paths \(\ell_y\) from the basepoint to each \(y\in X\),
taking the basepoint path constant. Assign to an oriented
singular edge \(a:y\to z\) the class of
\(\ell_y*a*\ell_z^{-1}\) in the abelianization and extend
linearly to one-chains. The signed boundary of any
two-simplex maps to zero, since its two consecutive edges
are homotopic rel endpoints to the remaining edge by the
image of the triangle. This assignment therefore descends
on cycles to homology. On a based loop it is the original
abelianized loop. In the other direction its image on an
edge differs from the edge chain by the two connector paths,
modulo boundaries, by the concatenation identity. In a
cycle their signed occurrences cancel at every vertex.
The descended maps are inverse.
\end{proof}

\section{Ordinary and relative cohomology}\label{p202:T23}

For an abelian group \(A\), put
\[
 C^q(X,Y;A)=\operatorname{Hom}_{\mathbb Z}
       (C_q(X;\mathbb Z)/C_q(Y;\mathbb Z),A),
 \qquad (\delta a)(c)=a(\partial c).
\]
Equivalently these are the cochains on \(X\) vanishing on every
simplex in \(Y\). Write \(H^q(X,Y;A)\) for their cohomology.

\begin{lemma}[Dualizing free chain comparisons]\label{lem:p202-free-chain-dual}
A quasi-isomorphism between bounded-below complexes of free
modules over a ring is a chain homotopy equivalence. Consequently
applying \(\operatorname{Hom}(-,A)\) gives a cohomology
isomorphism for every coefficient module \(A\).
\end{lemma}
\begin{proof}
Its mapping cone is a bounded-below acyclic complex of free
modules. In the lowest possible degree its boundary group
is the whole chain group. Inductively, if \(B_{q-1}\) is
projective, the surjection \(C_q\to B_{q-1}\) splits, and its
kernel \(B_q\) is projective as a direct summand. The splittings
give a contraction by sending the boundary summand of \(C_q\)
to its chosen section in \(C_{q+1}\) and sending the complementary
summand to zero. For clarity, write the cone of a map
\(f:C\to D\) in degree \(q\) as \(D_q\oplus C_{q-1}\) with
differential \((y,x)\mapsto(\partial y+fx,-\partial x)\).
The \(D_q\to C_q\) block of the contraction is a chain map
\(g:D\to C\), by the lower-left block of
\(\partial h+h\partial=1\). Its upper-left and lower-right
blocks give respectively \(fg\simeq 1_D\) and \(gf\simeq1_C\).
Dualizing these chain homotopies proves the final assertion.
\end{proof}

\begin{proposition}[Cohomology of pairs and open covers]\label{prop:p202-cohomology-pairs}
Ordinary cohomology is contravariant and homotopy invariant.
Every pair has a natural long exact sequence in cohomology.
Excision holds for the same closure-inside-interior
hypothesis as in singular homology, and an open two-set cover
gives the ordinary cohomological Mayer--Vietoris sequence.
\end{proposition}
\begin{proof}
Pullback of cochains is precomposition with the induced chain
map, and precomposition with the prism homotopy proves
homotopy invariance. Restriction \(C^q(X;A)\to C^q(Y;A)\)
is onto because a cochain can be extended by zero on simplices
outside \(Y\). Its kernel is \(C^q(X,Y;A)\). For the connecting
map, lift a cocycle on \(Y\) to a cochain on \(X\) and take its
coboundary, which vanishes on \(Y\). Replacing the lift changes
this by a relative coboundary. Lifting a primitive, or restricting
a proposed primitive, gives exactness at each of the three
terms. This description also proves naturality. Excision on
chains is a quasi-isomorphism between free relative complexes,
so Lemma~\ref{lem:p202-free-chain-dual} applies.

For \(X=U\cup V\), the chain sequence for the small complex
\(C_*(U)+C_*(V)\) is split degreewise by distributing the
basis simplices in \(U\cap V\) to one of the two factors.
Its dual is therefore short exact. Apply the cochain long
exact sequence and compare the dual small complex with
\(C^*(X;A)\) by the same lemma and the small-chains theorem.
This derives Mayer--Vietoris before any use of a dual of
homology.
\end{proof}

\section{Product chains and the K\"unneth sequence}\label{p202:T26}

For a singular \(n\)-simplex \(\sigma\), denote restriction
to its ordered face with vertices \(i_0,\ldots,i_j\) by
\(\sigma[i_0,\ldots,i_j]\). The tensor differential is
\(\partial(a\otimes b)=\partial a\otimes b+
(-1)^{\deg a}a\otimes\partial b\).

\begin{theorem}[Eilenberg--Zilber equivalence]\label{thm:p202-ez}
For arbitrary spaces \(X,Y\), there are natural chain homotopy
inverse maps between \(C_*(X\times Y;R)\) and
\(C_*(X;R)\otimes_R C_*(Y;R)\). The map to the tensor product is
\[
 \operatorname{AW}(x,y)=
 \sum_{i=0}^n x[0,\ldots,i]\otimes y[i,\ldots,n].
\]
The inverse sends \(\sigma\otimes\tau\), of bidegree \((p,q)\),
to the oriented shuffle triangulation of
\(\sigma\times\tau:\Delta^p\times\Delta^q\to X\times Y\).
The latter map induces the homological external product.
\end{theorem}

\begin{proof}
In \(\partial\operatorname{AW}(x,y)\), deleting vertex \(j<i\)
from the first face and deleting \(j>i\) from the second face
are exactly the terms of \(\operatorname{AW}\partial(x,y)\).
Their sign is \((-1)^j\), including the tensor differential
sign in the second case. The terms deleting the shared vertex
from the first face at split \(i\) and from the second face at
split \(i-1\) have opposite signs and cancel. The boundary of
a zero-face is zero, which removes the two extraneous endpoint
terms. This proves the chain identity.

A shuffle is a path from \((0,0)\) to \((p,q)\) using \(p\)
horizontal and \(q\) vertical steps. Its simplex has consecutive
vertices \((v_i,w_j)\) along the path. Assign it the sign of
the permutation of the \(p+q\) steps relative to all horizontal
steps followed by all vertical steps. These simplices triangulate
the product. In barycentric cumulative coordinates, ordering
the \(p\) coordinates from the first factor and the \(q\) from
the second factor selects one shuffle region; equality of
coordinates gives the shared faces. Subtract consecutive
vertex columns in the affine determinant. The resulting
columns are successive edge vectors from the two factors, so
the determinant sign is precisely the shuffle sign. In the
boundary of the signed sum, the face omitting a vertex between
steps of different kinds occurs twice, by interchanging the
two steps, with opposite coefficients. The surviving faces
belong to \(\partial\Delta^p\times\Delta^q\) and
\(\Delta^p\times\partial\Delta^q\), including the initial
and terminal vertex omissions. Their orientations have the
ordinary boundary signs on the first factor and the extra
\((-1)^p\) on the second factor: moving its outward vector
past the \(p\) oriented first-factor coordinates accounts for
this sign. Pushing forward proves the chain identity for the
shuffle map \(\operatorname{sh}\).

Here are the homotopies. A convex model has an augmented
singular contraction by coning to a fixed vertex; the
boundary of the cone is the original chain minus the cone
on its boundary, with the augmentation term in degree zero.
The tensor of two convex-model complexes is similarly
acyclic above degree zero. Explicitly, if \(P_C,P_D\) are
the augmentation-to-vertex projections and \(h_C,h_D\) the
contractions, use
\[
 h(c\otimes d)=h_C(c)\otimes d+
               (-1)^{\deg c}P_C(c)\otimes h_D(d).
\]
Then \(\partial h+h\partial=1-P_C\otimes P_D\).

For the composite \(\operatorname{sh}\operatorname{AW}\) and
the identity, the universal generator in degree \(n\) is the
diagonal simplex in \(\Delta^n\times\Delta^n\). The difference
of the two chain maps vanishes in degree zero. Inductively,
subtract the already defined homotopy on the boundary from
their difference on this universal generator. Its boundary
is zero by the lower-degree identity. Fill by the contraction
in the convex product, and define the homotopy on an arbitrary
generator by pushing forward that universal filling under
\(x\times y\). This preserves the inductive identity on its
faces because the lower-degree choices were made on universal
simplices. For \(\operatorname{AW}\operatorname{sh}\) the
universal generator is the pair of identities on
\(\Delta^p,\Delta^q\); apply the same induction in total
degree using the tensor contraction. These choices prove
both chain homotopies and their naturality.
\end{proof}

\begin{corollary}[Homological K\"unneth over a PID]\label{cor:p202-topological-kunneth}
If \(R\) is a PID, for arbitrary \(X,Y\) there is a natural
short exact sequence
\[
 \begin{split}
 0&\longrightarrow \bigoplus_{i+j=n}H_i(X;R)\otimes_R H_j(Y;R)
 \longrightarrow H_n(X\times Y;R)\\
 &\longrightarrow \bigoplus_{i+j=n-1}
   \operatorname{Tor}_1^R(H_i(X;R),H_j(Y;R))
 \longrightarrow0 .
 \end{split}
\]
The first map is the external product. The sequence splits
noncanonically, and the last sum is zero in degree zero.
\end{corollary}
\begin{proof}
The singular groups are free in every degree, and their
boundary submodules are free over a PID by \cite[Lemma 11.7]{P103}.
Apply the algebraic K\"unneth theorem of \cite[Chapter 11]{P103}
to the nonnegative tensor complex,
then transport by Theorem~\ref{thm:p202-ez}. The first map
sends tensor products of cycle classes to the corresponding
tensor cycles before applying the shuffle map. The splitting
uses choices of sections of boundary projections and is
not asserted to be natural.
\end{proof}

The two chain constructions can also be restricted to smooth
singular simplices, since all affine faces and shuffle
parametrizations are smooth. If \(\sigma\) and \(\tau\) are
smooth, the shuffle chain is the oriented triangulation
of their product used in integration as well.
See \cite[Sections 3.2 and 3.B]{Hatcher} and
\cite[Sections 7.7--7.11]{MKP}.

\section{The cup product}\label{p202:T24}

\begin{theorem}[The singular cohomology ring]\label{thm:p202-cup}
For a commutative ring \(R\), \(H^*(X;R)\) is a natural
unital associative graded ring, and homogeneous classes
satisfy \(a\smile b=(-1)^{pq}b\smile a\) for
\(a\) of degree \(p\) and \(b\) of degree \(q\).
\end{theorem}
\begin{proof}
On cochains of degrees \(p,q\) use
\[
 (a\smile b)(\sigma)=
 a(\sigma[0,\ldots,p])\,b(\sigma[p,\ldots,p+q]).
\]
It is evaluation on the \((p,q)\) component of the chain
diagonal \(D=\operatorname{AW}\Delta_\#\), where
\(\Delta:X\to X\times X\) is the diagonal. The chain
identity gives
\(\delta(a\smile b)=\delta a\smile b+(-1)^p a\smile\delta b\).
Evaluating three factors cuts the same simplex at vertices
\(p\) and \(p+q\) in either bracketing, proving associativity
already on cochains. The constant \(0\)-cochain \(1\) is a
unit, and pullback commutes with the displayed formula.

Let \(T(x\otimes y)=(-1)^{\deg x\,\deg y}y\otimes x\).
Both \(D\) and \(TD\) are natural chain maps, and agree on
vertices. On the universal generator \(\operatorname{id}_{\Delta^n}\)
their difference, minus the homotopy already defined on its
boundary, is a cycle in \(C_*(\Delta^n)\otimes C_*(\Delta^n)\).
Fill by the tensor contraction in the proof of
Theorem~\ref{thm:p202-ez} and push forward. This proves
\(D\simeq TD\). Evaluating with cocycles makes the
homotopy difference a coboundary and gives the graded sign
on classes.
\end{proof}

\section{Products for relative pairs}\label{p202:T25}

\begin{proposition}[Relative cup product]\label{prop:p202-relative-cup}
For open subspaces \(U,V\subset X\), there is a natural
product
\[
 H^p(X,U;R)\otimes_R H^q(X,V;R)
     \longrightarrow H^{p+q}(X,U\cup V;R)
\]
compatible with passage to absolute cohomology. The long
exact sequence of an individual pair requires no openness
hypothesis.
\end{proposition}
\begin{proof}
The front face of a simplex contained in \(U\) still lies
in \(U\); the back face of a simplex contained in \(V\)
lies in \(V\). Thus the displayed cup formula vanishes on
\(S=C_*(U)+C_*(V)\). It descends to cohomology of the dual
of \(C_*(X)/S\) by the coboundary identity. The inclusion
\(S\hookrightarrow C_*(U\cup V)\) is a quasi-isomorphism
by small chains. Comparison of quotient exact sequences
makes the surjection \(C_*(X)/S\to C_*(X,U\cup V)\) a
quasi-isomorphism as well. Both complexes are degreewise free
since their defining subgroups are spanned by subsets of the
singular basis. Lemma~\ref{lem:p202-free-chain-dual} gives
the required cohomology comparison with arbitrary coefficients.
All maps are induced by inclusion or pullback; inverting
the induced isomorphism gives naturality for maps respecting
the two subspaces. The last assertion was proved in
Proposition~\ref{prop:p202-cohomology-pairs}.
\end{proof}

\section{Universal coefficients for singular chains}\label{p202:T27}

\begin{theorem}[Singular universal coefficients over a PID]\label{thm:p202-topological-uct}
For a PID \(R\), an \(R\)-module \(A\), and any space \(X\),
there are natural exact sequences
\[
 \begin{split}
 0&\longrightarrow H_n(X;R)\otimes_R A
 \longrightarrow H_n(X;A)
 \longrightarrow \operatorname{Tor}_1^R(H_{n-1}(X;R),A)
 \longrightarrow0,\\
 0&\longrightarrow \operatorname{Ext}^1_R(H_{n-1}(X;R),A)
 \longrightarrow H^n(X;A)
 \longrightarrow \operatorname{Hom}_R(H_n(X;R),A)
 \longrightarrow0 .
 \end{split}
\]
They split noncanonically. The last arrow in the second
sequence is evaluation on cycles.
\end{theorem}
\begin{proof}
The singular modules are free, and \cite[Lemma 11.7]{P103} gives
free boundary modules. Thus both the flatness and projectivity hypotheses of the
algebraic universal coefficient theorem in
\cite[Chapter 11]{P103} hold. The same applies to the
relative complex, whose basis consists of simplices outside
the chosen subspace. Tensoring the free singular groups
with \(A\) is precisely the \(A\)-valued chain construction;
\(\operatorname{Hom}_R(-,A)\) gives the corresponding
cochains. The algebraic theorem gives the sequences with
\(H_{-1}=0\). To see the role of the splitting, choose
sections of \(C_n\to B_{n-1}\). This separates the summands
\(Z_n\) and \(B_{n-1}\) in \(C_n\), reducing the differential
to the inclusion \(B_{n-1}\hookrightarrow Z_{n-1}\).
Applying tensor or Hom to these two-term complexes gives
respectively the Tor or Ext summands. Evaluation on cycles
is independent of these choices, whereas the sections are
not canonical. Over a field evaluation identifies cohomology
with the full linear dual of homology, without a
finite-dimensional hypothesis.
\end{proof}

\chapter{Orientation, Duality, and Comparison}

Topological manifolds in this chapter are Hausdorff and second
countable. In the duality theorems they are finite-dimensional;
a boundary is used only when explicitly indicated. We keep the
coefficient module and support convention visible in each duality
map. None of the homology groups without a support superscript
uses locally finite chains.

\section{Local orientation and the cap product}\label{p202:T29}

For \(x\) in a boundaryless \(n\)-manifold \(M\), excising outside
a coordinate ball identifies
\[
 H_j(M,M\setminus\{x\};R)
  \simeq H_j(\mathbb R^n,\mathbb R^n\setminus\{0\};R).
\]
This is \(R\) for \(j=n\) and zero otherwise, by the punctured
ball's radial retraction and the sphere calculation. In dimension
zero the local pair is a point and the empty space.
An \(R\)-orientation chooses a generator \(o_x\) of the local
degree-\(n\) module at each \(x\), with the choice in a coordinate
ball coming from a single oriented chart generator times a
locally constant unit. Restricting an oriented closed ball class
to any interior point gives the same chart orientation: its
boundary sphere surrounds that point and can be radially
parametrized from it, with outward orientation. Thus the
agreement condition is independent of the chosen sufficiently
small balls.

The orientation cover has as its fiber the two generators of the
integral local group. A coordinate chart and its signed local
generators trivialize it over the chart. On overlapping charts
the sign is locally constant, since restricting a ball class to
points of a smaller ball compares both local choices. The
trivializations therefore glue to a two-sheeted covering; a
section is exactly an integral orientation. Over \(\mathbb F_2\)
there is a unique local generator, so every manifold is
canonically oriented with those coefficients.

\begin{theorem}[Fundamental classes near compact sets]\label{thm:p202-compact-orientation}
If \(M\) has an \(R\)-orientation and \(K\subset M\) is compact,
there is a unique class
\(\mu_K\in H_n(M,M\setminus K;R)\) restricting to \(o_x\)
at every \(x\in K\). The classes are compatible for \(K\subset L\).
Moreover \(H_j(M,M\setminus K;R)=0\) for \(j>n\), and
restriction to the local groups is injective in degree \(n\).
For compact boundaryless \(M\), \(\mu_M\) is its fundamental
class in ordinary homology.
\end{theorem}

\begin{proof}
Write \(H_j(M\mid K)=H_j(M,M\setminus K;R)\). Two compacts
\(K,L\) give an exact sequence
\[
 \begin{aligned}
 \cdots&\to H_{j+1}(M\mid K\cap L)
 \to H_j(M\mid K\cup L)\\
 &\to H_j(M\mid K)\oplus H_j(M\mid L)
 \to H_j(M\mid K\cap L)\to\cdots .
 \end{aligned}
\]
Indeed, with \(P=M\setminus K\), \(Q=M\setminus L\), use the
short exact sequence of chain quotients by
\(C(P\cap Q)\), by \(C(P)\) and \(C(Q)\), and by
\(C(P)+C(Q)\). Its maps are diagonal and difference, and
the last quotient computes the pair with \(P\cup Q\) by
small chains since \(P,Q\) are open. In degrees above \(n\)
this sequence gives the claimed vanishing for \(K\cup L\)
if it holds for the three other compacts. In degree \(n\),
vanishing of the term \(H_{n+1}(M\mid K\cap L)\) makes the
middle map injective. Compatible classes on \(K,L\) lift
because they agree on the intersection by its local
injectivity. This proves that the three assertions pass to
a union under those hypotheses.

For the local step take an open coordinate ball identified
with all of \(\mathbb R^n\). By excision the assertions hold
for a compact convex set in this chart. For a nonempty such set and any point
\(x\) of it, projection along rays from \(x\) to a surrounding
sphere gives a homotopy equivalence on the complement.
Every ray beyond a point of the complement stays outside
the convex set, and the radial deformation to a sphere
containing the set stays outside it. This identifies the
top relative group with \(R\), makes restriction at \(x\)
an isomorphism, and supplies the generator determined by
the chart orientation. The empty set gives zero groups.
Finite unions of convex compact sets in one chart follow
by induction, since intersection with the last convex set
is a union of fewer convex compact sets.

For an arbitrary compact \(K\) in that chart, take a relative
cycle and let \(B\) be the compact image of its boundary
chain. It is disjoint from \(K\). A finite union \(L\) of
sufficiently small closed balls centered in \(K\), still inside
the chart, contains \(K\) and misses \(B\). The chain defines
a class for \(L\). If the degree is greater than \(n\), it
vanishes there. If the degree is \(n\) and every local value
on \(K\) is zero, restriction of this class to any of the
balls vanishes by its value at the center. Hence the local
values on \(L\) all vanish, and the class for \(L\) is zero.
Existence of the orientation class on \(K\) follows by
restricting the class for a larger closed ball after identifying
the chart with \(\mathbb R^n\). In dimension zero the same
assertions follow by summing over the finitely many points
of a compact subset of a discrete manifold.

Finally cover \(K\) by finitely many compact pieces in
coordinate charts, obtained by intersecting it with closures
of smaller precompact chart neighborhoods. On adding the last
piece the intersection is a union of fewer pieces, each in a
chart. The first induction therefore proves the assertions
for the whole compact. Compatibility under \(K\subset L\)
follows from the uniqueness with prescribed local values.
\end{proof}

\begin{proposition}[Cap identities]\label{prop:p202-cap}
For \(a\in C^p(X;R)\) and a singular \(q\)-simplex \(\sigma\)
with \(q\ge p\), put
\[
 a\cap \sigma=a(\sigma[0,\ldots,p])\,\sigma[p,\ldots,q],
\]
and use zero for \(q<p\). Then
\[
 \partial(a\cap c)=(-1)^p\bigl(a\cap\partial c-(\delta a)\cap c\bigr),
 \qquad
 \langle b,a\cap c\rangle=\langle a\smile b,c\rangle .
\]
For a continuous \(f:X\to Y\) one has
\(f_\#((f^*a)\cap c)=a\cap f_\#c\).
In particular there are pair cap products
\[
 H^p(X;R)\otimes H_q(X,B;R)\to H_{q-p}(X,B;R),
 \qquad
 H^p(X,B;R)\otimes H_q(X,B;R)\to H_{q-p}(X;R),
\]
where tensor products in this chapter's operation formulas
are over \(R\).
\end{proposition}
\begin{proof}
For faces with deleted index \(j>p\) in \(\sigma\), the
first formula reproduces the back-face boundary signs after
multiplication by \((-1)^p\). The deleted indices \(j\le p\)
combine with \((\delta a)(\sigma[0,\ldots,p+1])\) and leave
only the term deleting the first vertex of the back face.
This proves the boundary formula, including \(q=p\) where
all terms of negative degree are zero. Evaluation of \(b\)
gives exactly the front-back cup formula, and composing
faces with \(f\) proves naturality. Chains supported in \(B\)
remain supported in \(B\) under cap, and are killed when
\(a\) vanishes on \(B\). These facts, with the boundary
identity, give the two products.
\end{proof}

\section{Poincar\'e duality and compact supports}\label{p202:T30}

Define
\[
 H_c^p(M;R)=\varinjlim_{K\subset M\text{ compact}}
                       H^p(M,M\setminus K;R).
\]
The maps for increasing \(K\) extend the allowed supports,
and every equality in the limit is witnessed on some larger
compact set. Equivalently one may take the cohomology of the
union of relative cochain complexes, since a cocycle or a
primitive lies in one of these complexes.

\begin{theorem}[Poincar\'e duality]\label{thm:p202-poincare}
For an \(R\)-oriented boundaryless \(n\)-manifold \(M\), the
maps \(a\mapsto a\cap\mu_K\) induce isomorphisms
\[
 D_M:H_c^p(M;R)\longrightarrow H_{n-p}(M;R).
\]
For compact \(M\) this is cap with \(\mu_M\) on ordinary
cohomology. The target in the noncompact case still uses
finite singular chains.
\end{theorem}

\begin{proof}
If \(K\subset L\), a cochain supported on \(K\) kills
chains in \(M\setminus K\), and \(\mu_L\) maps to \(\mu_K\).
Thus the formulas agree when supports are enlarged. If
\(U\subset M\) is open and \(K\subset U\), excision identifies
the groups near \(K\) computed in \(U\) and in \(M\).
This defines extension of compact supports and, by naturality
of cap and local orientation classes, \(D\) commutes with
open inclusion on the homology side.

We verify the connecting map needed for gluing. Suppose
\(M=U\cup V\) and let \(K\subset U\), \(L\subset V\) be
compact. Write \(P=M\setminus K\), \(Q=M\setminus L\).
A cocycle \(a\) vanishing on \(P\cap Q\) can be written
\(a=a_P+a_Q\), where \(a_P\) vanishes on \(P\) and \(a_Q\)
on \(Q\). This follows by assigning its values separately
on the two basis subsets, which agree at zero on the
intersection. The cohomological connecting class is
represented by \(\delta a_P=-\delta a_Q\), which vanishes
on \(C(P)+C(Q)\). Small chains identify this relative
complex with supports on \(K\cap L\) as in
Section~\ref{p202:T25}. The resulting compact-support
Mayer--Vietoris sequence follows from this short exact
sequence and the filtered limit. In the limit the
sets \(K\cup L\) are cofinal among compacts of \(M\):
intersect a given compact with a finite cover by closed
pieces subordinate to \(U,V\). The sets \(K\cap L\) are
cofinal among compacts of \(U\cap V\).

Represent \(\mu_{K\cup L}\) by a relative cycle \(c\)
subdivided subordinate to \(U\setminus L,\ U\cap V,\
V\setminus K\). These open sets cover \(M\). Write the
corresponding decomposition as \(c=c_U+c_W+c_V\).
Then \(c_W\) represents \(\mu_{K\cap L}\) in \(U\cap V\),
and \(c_U+c_W\) represents \(\mu_K\) in \(U\), by excision
and compatibility. In particular its boundary is supported
away from \(K\). The homology connecting map of \(a\cap c\)
is represented by \(\partial(a\cap c_U)\), since the other
two pieces lie in \(V\). For \(p=\deg a\), the cap identity
gives
\[
 \partial(a\cap c_U)
 =(-1)^p a_P\cap\partial c_U
 =(-1)^{p+1}a_P\cap\partial c_W .
\]
Here \(a_Q\) kills \(\partial c_U\), and \(a_P\) kills
\(\partial(c_U+c_W)\). Modulo a boundary in \(U\cap V\),
\(a_P\cap\partial c_W\) equals \((\delta a_P)\cap c_W\),
by the same identity. This proves commutation of the
connecting square up to the indicated unit sign. The
remaining squares commute by open inclusion and linearity.
In both sequences the map from the intersection has opposite
signs in the two components, and the map from the two pieces
is their sum. A connecting square commuting up to a unit
sign suffices for the exact diagram chase, since a sign
does not alter image or kernel. Hence duality for \(U,V\)
and \(U\cap V\) implies duality for \(M\) by the five lemma.

For \(\mathbb R^n\), take a cofinal sequence of large closed
balls. The pairs have cohomology only in degree \(n\),
where it is \(R\). This follows by the relative disk
calculation and its dual, or by the exact sequence and
radial retraction of the complement. In that degree
cap sends a cocycle evaluating to \(1\) on the ball
class to a zero-chain whose augmentation is \(1\);
the cap pairing with the constant \(0\)-cochain says
exactly this. It therefore gives the isomorphism to
\(H_0(\mathbb R^n;R)\). The dimension-zero calculation
is the same on a point.

Duality passes to increasing open unions. Finite chains and
their bounding chains have compact image; compact supports
likewise lie in one member of an increasing open cover.
For the latter assertion use the open-inclusion excision
identification at each compact. Thus both groups are the
corresponding direct limits and \(D\) is their compatible
map. An open set in \(\mathbb R^n\) has a countable cover
by open boxes contained in it, whose nonempty finite
intersections are boxes. A finite union of these satisfies
duality by the two-set step, inducting also on the finite
union of intersections with the last box. The increasing
union gives duality on any open subset of a chart. Finally
take a countable chart cover of \(M\). Intersections of
charts are open subsets of charts. Apply the same finite
and increasing-union arguments to conclude.
\end{proof}

\section{Poincar\'e--Lefschetz duality}\label{p202:T31}

\begin{theorem}[Duality for a compact manifold with boundary]\label{thm:p202-lefschetz}
Let \(M\) be a compact \(R\)-oriented \(n\)-manifold with
collared boundary \(B=\partial M\), with the induced boundary
orientation. There is a relative fundamental class
\(\mu\in H_n(M,B;R)\) restricting to the orientation at
interior points and with \(\partial\mu=\mu_B\).
The cap maps
\[
 H^p(M,B;R)\longrightarrow H_{n-p}(M;R),\qquad
 H^p(M;R)\longrightarrow H_{n-p}(M,B;R)
\]
are isomorphisms.
\end{theorem}
\begin{proof}
Put \(U=M\setminus B\) and choose a collar \(B\times[0,2\epsilon)\).
A compact core \(K\subset U\) is obtained by removing
\(B\times[0,\epsilon)\). The removed portion retracts to
\(B\). The pair sequence and excision identify \(H_n(M,B)\)
with \(H_n(U,U\setminus K)\). The latter group has its
compatible orientation class. Cores of this form with
\(\epsilon\) tending to zero exhaust \(U\) and contain
every compact subset eventually. By the compatibility
of compact orientation classes this construction is
independent of the core. To compute its boundary at
\(y\in B\), choose a closed collar box \(C\) centered on
the boundary face at \(y\), and put \(A=M\setminus
\operatorname{int}_M C\). In the triple \((M,A\cup B,A)\)
the connecting map from \(H_n(M,A\cup B)\) to
\(H_{n-1}(A\cup B,A)\) is natural from the triple
\((M,B,\varnothing)\). Excision into a larger open collar
box and coordinate retraction onto \(C\) identify these
groups with the relative group of the box with its entire
boundary and the relative group of its bottom face with
the boundary of that face. The localized class is the
relative box class, since its value at an
interior point is the given orientation generator and the
relative box group maps isomorphically to that local group.
The boundary of this product along
the face at \(0\) is the disk orientation with the outward
normal first. This follows from the oriented shuffle
boundary formula in Section~\ref{p202:T26}; order the
interval first and use its lower-end boundary sign.
Thus \(\partial\mu\) has the prescribed local value
at each boundary point, and uniqueness of the fundamental
class on \(B\) gives \(\partial\mu=\mu_B\). The empty-boundary
case recovers Theorem~\ref{thm:p202-poincare}.

For the first map, restriction to the interior and the
same core construction identify \(H^p(M,B)\) with
\(H_c^p(U)\). More explicitly, for each core the maps
\(H^p(M,B)\leftarrow H^p(M,M\setminus K)\to
H^p(U,U\setminus K)\) are isomorphisms by the collar
retraction and excision. The inclusion \(U\hookrightarrow M\)
is a homotopy equivalence by pushing the collar a
positive distance inward. The cap formula is therefore
the interior duality map under these identifications.

Compare the cohomology exact sequence of the pair with
the homology exact sequence in complementary degrees:
\[
\begin{gathered}
 H^{p-1}(B)\longrightarrow H^p(M,B)\longrightarrow H^p(M)
 \longrightarrow H^p(B)\longrightarrow H^{p+1}(M,B),\\
 H_{n-p}(B)\longrightarrow H_{n-p}(M)\longrightarrow H_{n-p}(M,B)
 \longrightarrow H_{n-p-1}(B)\longrightarrow H_{n-p-1}(M).
\end{gathered}
\]
The vertical maps in order are cap with \(\mu_B,\mu,\mu,\mu_B,\mu\).
The map into relative homology commutes with cap by the quotient
description. If \(c\) represents \(\mu\), the boundary
of \(a\cap c\) for an absolute cocycle \(a\) is
\((-1)^p a|_B\cap\partial c\). For the cohomology
connecting arrow lift a boundary cocycle \(b\) to \(a\);
then \((\delta a)\cap c\) differs from
\(b\cap\partial c\) by a boundary, by the same cap identity.
These are the two other types of squares, with possible
unit signs. The four surrounding vertical maps are
isomorphisms by the first map and closed-boundary duality.
Exactness and the five lemma prove the middle isomorphism.
\end{proof}

\section{Alexander duality for compact locally contractible subsets}
\label{p202:T32}

\begin{lemma}[Euclidean neighborhood retraction]\label{lem:p202-enr}
Let \(K\subset\R^n\) be compact. Suppose that for every \(x\in K\)
and every relative neighborhood \(U\) of \(x\), there is a smaller
relative neighborhood \(V\) of \(x\) whose inclusion into \(U\)
is nullhomotopic. Then some open neighborhood of \(K\) in
\(\R^n\) retracts onto \(K\).
\end{lemma}

\begin{proof}
The empty set and the zero-dimensional ambient space are immediate.
For \(\epsilon>0\), compactness and the hypothesis give
\(\delta(\epsilon)>0\) with the following property. A continuous
map from the boundary of any positive-dimensional simplex to \(K\)
whose image has diameter less than \(\delta(\epsilon)\) extends
over the simplex with image of diameter at most \(\epsilon\).
Indeed, for each \(y\in K\) choose a relatively open neighborhood
of \(y\) that contracts in \(K\cap B(y,\epsilon/3)\).
Take a Lebesgue number for a finite subcover of these neighborhoods.
A boundary image smaller than this number lies in one of the
contracting neighborhoods, and the contraction extends the map
over the cone on the boundary.

Put \(c=\sqrt n\). At dyadic scale \(s=2^{-j}\), \(j\in\Z\), take
the balls \(B(z,cs)\) for \(z\in s\Z^n\) satisfying
\(2cs\le d(z,K)\le9cs\). If \(x\) belongs to such a ball, then
\(cs<d(x,K)<10cs\). Thus at any \(x\) at most five scales can
occur, and at each scale at most \((\lceil 2c\rceil+2)^n\) centers
can occur. The family is locally finite in the complement of \(K\):
on a sufficiently small bounded neighborhood away from \(K\), the
distance is bounded above and below by positive constants, so
only finitely many scales and grid points meet it.
For \(x\notin K\), choose \(s\) with
\(4cs\le d(x,K)<8cs\) and a nearest point \(z\) in \(s\Z^n\).
Then \(|x-z|\le cs/2\) and \(2cs\le d(z,K)\le9cs\).
In particular, keeping only balls with \(s\le\tau\) still covers
all \(x\) with \(0<d(x,K)<4c\tau\).

The nerve has a vertex for each retained ball and a simplex for
each finite family having nonempty intersection. Its dimension is
at most the fixed integer \(m=5(\lceil 2c\rceil+2)^n\), regardless
of the upper cutoff on the scales. For every \(k\ge0\) choose bounds
\(\epsilon_{j,k}>0\), \(1\le j\le m\), starting from
\(\epsilon_{m,k}=2^{-k}\) and choosing the lower dimensions
successively so that
\[
 \epsilon_{j,k}\le 2^{-k},\qquad
 2\epsilon_{j,k}<\delta(\epsilon_{j+1,k})\quad(j<m).
\]
Choose them nonincreasing in \(k\) for each \(j\), by shrinking
each new lower-dimensional choice if necessary. Choose positive
\(\tau_k\downarrow0\), strictly decreasing, such that
\(20c\tau_k<\delta(\epsilon_{1,k})\). Retain scales \(s\le\tau_0\).
For a simplex \(\sigma\), let \(k(\sigma)\) be the largest \(k\)
such that every scale in \(\sigma\) is at most \(\tau_k\).
Taking a face can only increase \(k(\sigma)\).

Send a vertex for \(B(z,cs)\) to a closest point \(p_z\in K\).
If the balls of a simplex contain a common point \(x\), then
\(|p_z-x|\le10cs\) for each vertex. Hence the two vertex images
of an edge have distance at most \(20c\tau_{k(\sigma)}\).
Extend over every edge by the small-extension property.
Inductively require the image of each \(j\)-simplex \(\sigma\)
to have diameter at most \(\epsilon_{j,k(\sigma)}\).
For \(j\ge2\), the images of its codimension-one faces have
diameter at most \(\epsilon_{j-1,k(\sigma)}\) and intersect
pairwise, since the faces themselves do. Their union has
diameter at most \(2\epsilon_{j-1,k(\sigma)}\).
The same extension property completes the induction.
The chosen maps agree on faces and give a map \(f\) on the nerve
that is continuous on every finite subcomplex.

On \(0<d(x,K)<4c\tau_0\), assign each retained ball \(B_i\)
the weight \(d(x,\R^n\setminus B_i)\), normalized by the sum of
the weights. This sum is locally finite and positive, and the
nonzero weights are coordinates in a simplex of the nerve.
The resulting map, followed by \(f\), is continuous on this
punctured neighborhood, since locally only finitely many
coordinates occur. Define \(r(x)\) to be its value there and
put \(r(x)=x\) on \(K\). As \(x\) tends to \(y\in K\), every
ball contributing a weight has \(s<d(x,K)/c\), so
\(k(\sigma)\to\infty\). Its vertex image satisfies
\(|p_z-x|\le10cs\), and \(r(x)\) is within \(2^{-k(\sigma)}\)
of that vertex image. This also holds for a single vertex,
where the distance is zero. Thus \(r(x)\to y\), proving
continuity and the claimed retraction.
\end{proof}

The small-extension argument is the one used in
\cite[Appendix, Theorem A.7]{Hatcher}. Here the bounded-dimensional
nerve supplies the skeleta on which to extend. The hypothesis
places no local-flatness condition on the embedding.

\begin{theorem}[Alexander duality]\label{thm:p202-alexander}
If \(K\) is a nonempty proper compact locally contractible
subset of \(S^n\), \(n\geq1\), then
\[
 \widetilde H_i(S^n\setminus K;\mathbb Z)
   \simeq \widetilde H^{\,n-i-1}(K;\mathbb Z)
 \quad(i\geq0).
\]
Here reduced cohomology in a negative degree is zero for
nonempty \(K\). The same conclusion holds with coefficients in
any fixed field in place of \(\Z\).
\end{theorem}
\begin{proof}
Write \(U=S^n\setminus K\). The orientation of the sphere
orients \(U\). By compact-support duality and excision,
\[
 H_i(U;R)\simeq
 \varinjlim_{N\supset K} H^{n-i}(S^n,N;R),
\]
where \(N\) ranges over open neighborhoods of \(K\):
take the compact support \(S^n\setminus N\) in \(U\).
Here \(R\) is \(\Z\) or the chosen field.

We check that the colimit on the right is \(H^{n-i}(S^n,K)\).
Remove one point outside \(K\) to work in a Euclidean chart.
By Lemma~\ref{lem:p202-enr} there is a retraction
\(r:N_0\to K\) on an open neighborhood. Restriction
\(\varinjlim_N H^j(N)\to H^j(K)\) is surjective by \(r\).
For any neighborhood \(N\) of \(K\), a smaller one
\(N'\subset N\cap N_0\) can be chosen so that the line
segment from \(x\) to \(r(x)\) lies in \(N\) for every
\(x\in N'\). Indeed the distance from \(K\) to the
complement of \(N\) in the Euclidean chart has a positive
lower bound after shrinking around compact \(K\), and
\(\lvert x-r(x)\rvert\) is uniformly small near \(K\).
Thus restriction from \(N\) to \(N'\) factors in
cohomology through \(K\). This proves injectivity in the
colimit. Compare the pair long exact sequences and use
exactness of filtered colimits to obtain the relative
isomorphism.

For \(0<i<n\), the reduced pair sequence of the sphere
now gives the displayed shift directly. For \(i=0\) the
kernel of \(H_0(U)\to H_0(S^n)\) is reduced \(H_0(U)\),
since \(U\) is nonempty. Naturality of cap with extension
of compact supports identifies this map with
\(H^n(S^n,K)\to H^n(S^n)\); the kernel in the reduced
pair sequence is exactly \(\widetilde H^{n-1}(K)\).
For \(i=n\) the relative degree-zero group is zero
since \(S^n\) is connected and \(K\) is nonempty.
Higher degrees on the homology side are zero by duality
with negative-degree cohomology. This handles all endpoints.
\end{proof}

\section{Cohomology rings of basic spaces}\label{p202:T28}

\begin{lemma}[Classes of transverse coordinate submanifolds]
\label{lem:p202-projective-thom}
Let \(N\subset M\) be a compact embedded smooth submanifold
with normal bundle of rank \(d\) oriented over \(R\).
There is a unique class in \(H^d(M,M\setminus N;R)\)
restricting to the orientation generator on each normal
fiber at its origin. For finitely many such submanifolds
intersecting transversely, the relative product of their
classes restricts to the corresponding product orientation
generator on a normal slice of the intersection.
\end{lemma}
\begin{proof}
Use a tubular neighborhood and excision. In a
trivialization over an open subset \(W\subset N\) the
relative pair is
\((W\times\mathbb R^d,W\times(\mathbb R^d\setminus0))\).
The maps in Theorem~\ref{thm:p202-ez} respect a subspace
in the second factor, including their homotopies by
naturality. The relative fiber complex is chain homotopy
equivalent to \(R\) concentrated in degree \(d\): send
the generator to a relative cycle representing its
orientation class, use the relative disk calculation,
and apply Lemma~\ref{lem:p202-free-chain-dual} to the
free complexes. Tensoring the homotopy with the first
factor gives relative cohomology \(H^{*-d}(W;R)\).
In particular there is no cohomology below \(d\), and
in degree \(d\) a class is determined by its local
fiber values. This also holds for any open subset of a
trivialization. Now glue over a finite trivializing
cover of \(N\): the Mayer--Vietoris sequence of relative
cochains identifies degree \(d\) with compatible choices,
since degree \(d-1\) on the overlap is zero. This
vanishing itself propagates by the same sequence.
Induct on the number of open sets, applying the
same argument to their intersections. The orientation
makes the local generators compatible.

For the product, in coordinates at a transverse point
the individual normal equations can be used simultaneously
as coordinate blocks, by the submersion theorem.
Their relative classes become the pullbacks of the
generators in the individual disk pairs. The shuffle
map on the product of oriented relative disk cycles,
and the front-back formula, evaluate their external
product to \(1\). Indeed \(\operatorname{AW}\operatorname{sh}\)
is homotopic to the identity relative to all the disk
faces, so evaluation is the product of the individual
evaluations. Pullback by the normal-coordinate
identification gives the claimed local product orientation.
\end{proof}

\begin{proposition}[Tori, projective spaces, and surfaces]
\label{prop:p202-cohomology-rings}
The integral cohomology ring of \(T^m\) is exterior on the
\(m\) degree-one circle classes. Moreover
\[
 H^*(\mathbb{RP}^m;\mathbb F_2)
       \simeq \mathbb F_2[u]/(u^{m+1}),\quad \deg u=1,
 \qquad
 H^*(\mathbb{CP}^m;\mathbb Z)
       \simeq \mathbb Z[v]/(v^{m+1}),\quad \deg v=2 .
\]
For the closed oriented surface of genus \(g\), evaluation
of the cup product on the orientation class gives a
unimodular alternating form on \(H^1\), with a symplectic
basis of \(g\) pairs.
\end{proposition}
\begin{proof}
The homological external products of the circle
generators and point classes form an integral basis
by K\"unneth, since all groups are free. The product
of pulled-back circle cohomology classes evaluates
on the corresponding cycle basis by the product of
evaluations, from
\(\operatorname{AW}\operatorname{sh}\simeq1\).
Consequently the square-free monomials are a cohomology
basis by universal coefficients. Each individual square
is zero by pullback from the circle, and degree-one
classes anticommute by Theorem~\ref{thm:p202-cup}.
This gives the exterior ring over \(\mathbb Z\), including
the integral square relation.

For \(m=0\) each projective space is a point and its ring
has the asserted presentation. For \(m\geq1\) let \(a\)
be the image of the
relative class of a hyperplane, using its complex normal
orientation in the complex case and the canonical mod-two
choice in the real case. Coordinate hyperplanes have
the same class: they are carried to one another by
orthogonal rotations or unitary transformations joined
to the identity. Intersect \(i\) such hyperplanes with
a standard projective \(i\)-plane so that their common
intersection in that plane is one point. They meet
transversely there in \(i\) independent coordinate blocks.
By Lemma~\ref{lem:p202-projective-thom}, the \(i\)-th
power restricts to the local orientation cohomology
generator at that point, mapped into the ordinary
top cohomology of the projective \(i\)-plane.
Its evaluation on the fundamental class is \(1\) after
choosing the complex orientation, or \(1\) mod two,
because the global class restricts to the local
orientation class. The cellular calculation and universal
coefficients show that each of the indicated degrees is
cyclic of rank one over the given ring, with no further
degrees. The unit evaluation proves that every power up
to \(m\) is a generator, giving both presentations.

On the oriented surface, duality in degree one and
evaluation by universal coefficients show that the
cup pairing is perfect over \(\mathbb Z\). It is
skew-symmetric, and \(H^2\) is torsion-free, so it is
alternating. For a primitive vector \(a\), the induced
functional is primitive by perfection. Thus there is
a \(b\) pairing to \(1\) with \(a\). The span of
\(a,b\) splits off: subtract the two pairing coordinates
to project onto its orthogonal complement. The induced
form on the complement is again perfect. Induction
gives the stated symplectic basis.
\end{proof}

\section{The de Rham comparison}\label{p202:T33}

A smooth singular simplex is smooth on the standard simplex
as a manifold with corners, with local smooth extensions in
charts. Restricting to faces gives a chain subcomplex
\(C_*^{\mathrm{sm}}(M;\mathbb R)\subset C_*(M;\mathbb R)\).

\begin{lemma}[Smooth and continuous singular chains]
\label{lem:p202-smooth-singular}
For a finite-dimensional Hausdorff second-countable smooth
manifold \(M\), this inclusion induces an isomorphism in
homology. Restriction of real cochains gives the corresponding
natural cohomology isomorphism.
\end{lemma}
\begin{proof}
Small chains and the Mayer--Vietoris sequence hold for
smooth chains by the same affine constructions on each
standard simplex: subdivision and its chain homotopy are
finite sums of affine simplex maps into that simplex.
Composing with a smooth simplex preserves smoothness,
including the face identities and the support needed for
the subdivision proof. A smooth homotopy gives the prism
identity by affine parametrizations of its prism. Thus on
a convex Euclidean open set the smooth straight-line
contraction gives the same homology of a point as for
continuous chains, compatibly with inclusion.

An open subset of a chart is a countable union of convex
open boxes; finite intersections of boxes are convex or
empty. Apply Mayer--Vietoris and the five lemma to finite
unions, inducting on the intersections with the final box.
Increasing unions commute with both homologies because
cycles and bounding chains have compact image. This proves
the result on open subsets of charts. A countable chart
cover of a manifold without boundary now gives the global
isomorphism by the same procedure, since chart intersections
are open subsets of a chart. Restriction on real cochains
is an isomorphism on cohomology because the linear dual
of a complex of vector spaces has cohomology the dual of
homology, by evaluation as in
Theorem~\ref{thm:p202-topological-uct}. All comparisons
are induced by inclusion and are natural for smooth maps.

For a smooth boundary use a collar. A smooth inward
push on the collar gives a map \(M\to M^\circ\) homotopy
inverse to the inclusion, through smooth homotopies,
also on the interior. For example, in the collar coordinate
\(t\geq0\) replace \(t\) by \(t+s\epsilon\rho(t)\), where
\(\rho\) is a nonnegative smooth cutoff equal to \(1\) near
zero and supported inside the collar, and \(\epsilon>0\)
is small. This homotopy, together with its restriction
to the interior, proves the boundary case.
\end{proof}

\begin{theorem}[Natural de Rham ring isomorphism]
\label{thm:p202-de-rham}
Integration on smooth singular simplices induces a natural
isomorphism of graded real algebras
\[
 H_{\mathrm{dR}}^*(M)\simeq H^*(M;\mathbb R),
\]
where the product on forms is wedge and the product on
singular cohomology is cup.
\end{theorem}
\begin{proof}
The map \(\mathcal I_M\) evaluates a \(q\)-form \(\omega\)
on \(\sigma:\Delta^q\to M\) by \(\int_{\Delta^q}\sigma^*\omega\).
The signs on the faces are the boundary signs of the
standard oriented simplex. Stokes gives
\(\mathcal I_M(d\omega)(\sigma)=
\mathcal I_M(\omega)(\partial\sigma)\) in degree \(q+1\).
Composition of pullbacks gives naturality. These
chain-map facts are also proved in \cite[Chapter 12]{P104}.

On a convex open set the de Rham complex and the
smooth singular complex both have cohomology \(\mathbb R\)
in degree zero and zero above it. For forms the contraction
operator is obtained by pulling back along the straight-line
homotopy, decomposing the pullback into a term without
\(dt\) and \(dt\wedge\eta_t\), and integrating \(\eta_t\)
over \(t\). Differentiation and the fundamental theorem
of calculus give \(dK+Kd=F_1^*-F_0^*\).
For chains use the prism contraction. Integration sends
the constant \(1\) to the constant \(0\)-cochain \(1\),
so it is a cohomology isomorphism there.

For the global comparison we give the open-cover argument
explicitly, including infinite covers. For an open cover
\((U_i)\), indexed by an ordered set, form the double
complex in bidegree \((p,q)\) of the product of degree-\(q\)
forms over all \((p+1)\)-fold intersections, with alternating
restrictions horizontally and exterior differentiation
vertically. Define the corresponding double complex of
smooth singular cochains with coboundary vertically.
Take the total differential to be the horizontal
differential plus \((-1)^p\) times the vertical one. In
horizontal degree \(-1\) augment by global forms on the
forms side and by cochains on small smooth chains on the
singular side. Here \emph{small} means the simplex lies in
one member of the cover.

The augmented horizontal rows are exact. For forms,
take a locally finite smooth partition of unity subordinate
to the cover. If constructed on a refinement, group its
terms by a selected covering member containing their
support. Insert the extra index of each partition term
first in the alternating restriction formula, multiply by
that partition function, extend by zero, and sum.
This gives a contracting homotopy of the augmented row:
the insertion term with the extra index deleted sums
to the identity, and the other restrictions cancel
in pairs. For cochains the row splits, by evaluation,
into a product of augmented cochain rows of full simplices.
A small singular simplex \(\sigma\) belongs to a
nonempty set of covering members; evaluate its horizontal
row just on tuples in that set, and contract it by
inserting one chosen member containing \(\sigma\).
These identities give exactness without restricting the
number of sets in the cover. Small chains compare to
all smooth chains by subdivision and the linear dual
comparison used above.

We recall the elementary total-complex consequence.
When the augmented horizontal rows are exact, the
augmentation to the total complex is a quasi-isomorphism.
Indeed, regard the augmented double complex as the cone
of this map up to a shift and eliminate a cocycle's
components in increasing vertical degree. At its first
remaining vertical degree the cocycle equation says the
horizontal differential vanishes there; exactness provides
a primitive in that row. Subtract its total differential
and continue. In a fixed total degree there are finitely
many possible vertical degrees since horizontal degree is
bounded below. The process leaves zero. Likewise a map
which is a quasi-isomorphism in each vertical column
induces a total quasi-isomorphism: its columnwise cone
has exact columns, and the same elimination in increasing
horizontal degree proves that its total complex is acyclic.
The cone only shifts the common lower bound of vertical
degrees by one. Arbitrary products of vector-space complexes
preserve exactness, since primitives may be selected
coordinate by coordinate.

Cover an open subset of Euclidean space by convex open
boxes. Their nonempty finite intersections are convex, so
integration is a quasi-isomorphism in every vertical
column of the preceding comparison. Exactness of the
augmentations gives the global result on that open subset.
Now cover a manifold without boundary by its coordinate
domains. Every finite intersection is an open subset of
one domain, so the same argument applies. For a smooth
boundary use the smooth collar push and the homotopy
formula on forms, together with
Lemma~\ref{lem:p202-smooth-singular}. This proves the
natural vector-space isomorphism. The local and chain-map
conventions agree with \cite[Chapter 18]{LeeSM}.

To prove compatibility with products, let \(\omega,\eta\)
be closed forms on \(M,N\) of degrees \(p,q\). For smooth
singular cycles \(z,w\) of these degrees,
\[
 \int_{\operatorname{sh}(z\otimes w)}
       \pi_M^*\omega\wedge\pi_N^*\eta
   =\left(\int_z\omega\right)\left(\int_w\eta\right).
\]
The shuffle triangulation, with the determinant sign in
Theorem~\ref{thm:p202-ez}, integrates the product with
its product orientation. The equality follows simplex
by simplex from Fubini. Choose ordinary singular
cocycle representatives of the two cohomology classes
corresponding to the forms. Their restrictions agree
with the smooth integrals up to smooth coboundaries,
so their evaluations on smooth cycles agree with the
integrals. On these ordinary cocycles the continuous
chain homotopy \(\operatorname{AW}\operatorname{sh}\simeq1\)
identifies the external cup evaluation with the tensor
pairing. For a product cycle of bidegree
different from \((p,q)\), both evaluations are zero:
on a product parametrization at least one form has
degree larger than its factor, and on the singular
side \(\operatorname{AW}\operatorname{sh}\simeq1\)
identifies the evaluation with the tensor pairing of
the indicated bidegree. These product cycles span
real homology in each total degree by K\"unneth over
\(\mathbb R\) and smooth singular comparison. Evaluation
on homology detects real cohomology classes, even without
finite-dimensionality. Thus integration takes exterior
products of forms to exterior products of cohomology
classes. Pulling back along the diagonal gives the
wedge--cup assertion. The degree-zero unit follows from
its constant representative.
\end{proof}

\section{Separation and invariance of domain}\label{p202:T34}

\begin{theorem}[Jordan--Brouwer separation and invariance of domain]
\label{thm:p202-invariance-domain}
For \(n\geq1\), the complement of a topologically embedded
\(S^{n-1}\) in \(S^n\) has exactly two components, each with
boundary the embedded sphere. Also, a continuous injection
from an open subset \(U\subset\mathbb R^n\) into
\(\mathbb R^n\) is a homeomorphism onto an open subset.
\end{theorem}
\begin{proof}
The embedded sphere \(K\) is compact and locally
contractible. It is proper, since the two spheres in
dimensions \(n-1,n\) have different reduced homology.
Thus Alexander duality gives
\(\widetilde H_0(S^n\setminus K;\mathbb Z)\simeq\mathbb Z\).
The complement is locally path connected, hence it has
exactly two components. If \(O\) is any open neighborhood
of \(x\in K\), take a relative open disk \(D\) in \(K\)
contained in \(O\) such that \(K\setminus D\) is a closed
disk. In the \(S^0\) case take the singleton of \(x\).
The complement of this closed disk is connected by
Alexander duality, and equals \((S^n\setminus K)\cup D\).
If a component of \(S^n\setminus K\) missed \(O\),
its closure would miss \(D\). It would be both open and
closed in this larger complement, a contradiction.
Both components therefore have \(K\) in their boundaries,
and no point of the original open complement lies
in either boundary.

For invariance of domain take a closed ball
\(B\subset U\) about any given point of \(U\), with
positive radius. The restriction of the injection
to \(B\) is an embedding by compactness and the
Hausdorff property. In \(S^n=\mathbb R^n\cup\{\infty\}\)
the complement of the embedded boundary of \(B\)
has two components. It is partitioned by the image
of the open ball and the complement of the image
of the closed ball. The first is connected and the
second is connected by Alexander duality for the
embedded closed disk. Both are nonempty, the second
containing \(\infty\). They must therefore be exactly
the two open components. In particular the image of
the open ball is open. Doing this at each point proves
that the injection is open and has continuous inverse
onto its image.
\end{proof}

See \cite[Section 2.B]{Hatcher}. The argument does not
assume that the embedding has a collar.

\section{A homological fixed point criterion}\label{p202:T35}

\begin{theorem}[Lefschetz fixed point theorem for finite polyhedra]
\label{thm:p202-fixed-point}
If \(X\) is a finite polyhedron and \(f:X\to X\) is
continuous, put
\[
 \Lambda(f)=\sum_j(-1)^j
            \operatorname{tr}\bigl(f_*:H_j(X;\mathbb Q)
                                      \to H_j(X;\mathbb Q)\bigr).
\]
If \(\Lambda(f)\ne0\), then \(f\) has a fixed point.
\end{theorem}
\begin{proof}
For a finite chain complex over a field and a chain
endomorphism \(F\), trace is additive along invariant
subspaces. Use \(B_j\subset Z_j\subset C_j\).
The quotient \(C_j/Z_j\) is equivariantly \(B_{j-1}\),
so in the alternating trace the two boundary terms
cancel. The result equals the alternating trace on
homology.

If \(f\) has no fixed point, the displacement in a fixed
geometric metric has a positive minimum \(\delta\).
Choose a triangulation \(K\) of mesh less than
\(\delta/3\), then a subdivision \(K'\) admitting a
simplicial approximation \(g:K'\to K\) to \(f\).
The approximation has pointwise error less than
\(\delta/3\), since \(g(x)\) and \(f(x)\) lie in a
common target simplex. On oriented simplicial chains
of \(K\) use the endomorphism \(F=g_\#\mathrm{Sd}_\#\).
The subdivision map sends each simplex to its oriented
sum of subsimplices and commutes with the boundary
because interior faces cancel. Under simplicial--singular
comparison it induces the identity on the realization:
on each simplex both parametrization routes have support
in that simplex and the same vertex augmentation, so
the carrier contraction used in
Theorem~\ref{thm:p202-simplicial-comparison} gives a
chain homotopy. Thus \(F\) induces \(f_*\) by the
approximation homotopy.

For any subsimplex \(\tau\) of a simplex \(\sigma\)
of \(K\), \(g(\tau)\) cannot meet \(\sigma\).
Otherwise for some \(x\in\tau\) we would have
\(g(x)\in\sigma\) as well. The mesh bound and
approximation bound would imply
\(d(x,f(x))<2\delta/3\). Thus every diagonal coefficient
of \(F\) in the oriented simplex basis is zero.
The alternating trace formula gives \(\Lambda(f)=0\).
\end{proof}

See \cite[Section 2.C]{Hatcher} and \cite[Section 2.10]{MKP}.
\part{Riemannian Geometry}
\chapter{Metrics, Connections, and Geodesics}
\label{chap:p202-metrics}

\section{Riemannian metrics and their distance}\label{p202:G01}

A Riemannian metric on \(M\) is a smooth section of
\(\operatorname{Sym}^2T^*M\) whose value is a positive-definite inner
product on every tangent space. We write \(\langle\, ,\,\rangle\)
and \(\lvert\,\cdot\,\rvert_g\) for its pairing and norm.
A local isometry is a local diffeomorphism that preserves the metric.

\begin{theorem}[Existence, pullback, and length distance]
\label{thm:p202-metric-distance}
Every smooth manifold as above admits a Riemannian metric. The
pullback of a Riemannian metric by an immersion is Riemannian, and a
local isometry preserves lengths of piecewise \(C^1\) curves.
On each connected component the formula
\[
 d_g(p,q)=\inf_{\gamma:p\leadsto q}L_g(\gamma),\qquad
 L_g(\gamma)=\int_a^b\lvert\dot\gamma(t)\rvert_g\,dt
\]
defines a metric with the original topology. The infimum is over
piecewise \(C^1\) curves on compact intervals.
\end{theorem}

\begin{proof}
Choose positive Euclidean metrics in coordinate charts and a locally
finite subordinate smooth partition of unity. The weighted sum extends
smoothly across chart boundaries and is positive on each nonzero vector
because some weight is positive there. Equivalently this is the bundle
metric construction on \(TM\) in \cite[Chapter 5]{P104}.
An immersion has injective tangent map, so its pullback is positive
and smooth. The chain rule gives the length assertion.

Any two points of a connected smooth manifold can be joined by a
piecewise smooth curve. Indeed the set reachable from a fixed point
by such curves is open, as is its complement, by the use of straight
segments in small coordinate balls. Reversal and concatenation
give symmetry and the triangle inequality. Length is additive over
the pieces and is invariant under piecewise regular increasing
reparametrizations by change of variables.

For nondegeneracy and topology, take coordinates centered at \(p\)
on a neighborhood of a closed Euclidean ball of radius \(2a\).
On that compact coordinate ball choose positive constants \(c,C\)
so that \(c\lvert v\rvert\leq\lvert v\rvert_g\leq C\lvert v\rvert\).
For \(\lvert x(q)\rvert<a\) the straight segment gives
\(d_g(p,q)\leq C\lvert x(q)\rvert\). A path that leaves the
coordinate ball of radius \(a\) has length at least \(ca\) before
first exit. A path contained in that ball has length at least
\(c\lvert x(q)\rvert\). Thus \(d_g(p,q)\geq c\lvert x(q)\rvert\)
on the inner ball, and any point outside has distance at least \(ca\).
These estimates give nondegeneracy and both directions of the
topological comparison. In dimension zero the component is a point.
\end{proof}

The same metric construction works on a manifold with smooth boundary.
For the distance argument one uses convex half-balls in boundary charts.
See \cite[Chapter 2]{LeeRM} and \cite[Chapter 1]{dC}.

\section{Connections, pullbacks, and differentiation along curves}
\label{p202:G04}

For a smooth vector bundle \(E\to M\), a connection is a bilinear
map \((X,s)\mapsto\nabla_Xs\) on vector fields and sections, linear
over smooth functions in \(X\) and satisfying
\(\nabla_X(fs)=X(f)s+f\nabla_Xs\). Write a local frame as a row \(e\),
write a section as \(e v\) for a column \(v\), and use the connection
matrix of one-forms \(\Gamma\) determined by \(\nabla e=e\Gamma\).

\begin{proposition}[Local connection calculus]
\label{prop:p202-pullback-connection}
The local formula for \(\nabla\) is \(e(dv+\Gamma v)\).
Under a change of frame \(e'=e A\), its matrix becomes
\(\Gamma'=A^{-1}\Gamma A+A^{-1}dA\).
A smooth map \(F:N\to M\) induces a unique pullback connection on
\(F^*E\) that differentiates pulled-back sections in direction \(dF\).
For a \(C^1\) curve \(\gamma\) this gives
\[
 D_t(e(\gamma(t))v(t))=
 e(\gamma(t))\bigl(v'(t)+\Gamma_{\gamma(t)}(\dot\gamma(t))v(t)\bigr).
\]
Dual and tensor bundles carry induced connections satisfying the
product rule and commuting with contraction.
\end{proposition}

\begin{proof}
Apply the Leibniz rule to the frame expansion and to \(e A\).
The matrix transformation follows by comparing coefficients of \(e'\).
On \(N\), use the same formula with \(\Gamma\) pulled back as a
matrix of one-forms. The chain rule shows that the pullback transition
matrices satisfy exactly the same transformation law. Thus the local
formulas agree, and the rule on pulled-back frame sections and the
Leibniz rule give uniqueness. The displayed curve formula is its
one-dimensional form and remains meaningful with a continuous
matrix on a \(C^1\) segment. A change of parameter differentiates in
the new tangent direction, so \(D_s\) on a reparametrized field is
\(t'(s)D_t\) on the original one.

For a dual section set
\((\nabla_X\lambda)(s)=X(\lambda(s))-\lambda(\nabla_Xs)\).
On tensor products use
\(\nabla_X(s\otimes u)=(\nabla_Xs)\otimes u+s\otimes\nabla_Xu\).
The balanced-product relations follow from the scalar Leibniz rule.
Contraction commutes with the resulting derivative by the formula
on a dual pair. Iterating gives the result for all tensor powers.
\end{proof}

For a connection on \(TM\) its torsion is
\(\mathcal T(X,Y)=\nabla_XY-\nabla_YX-[X,Y]\).
The Lie bracket product rule shows that this is a tensor.
Metric compatibility means
\[
 X\langle Y,Z\rangle
  =\langle\nabla_XY,Z\rangle+\langle Y,\nabla_XZ\rangle .
\]
This is equivalent to \(\nabla g=0\) by the dual and tensor rules.
The identical identity holds for fields along a map.
In particular, for a smooth two-parameter family \(F\) with
\(V=F_*\partial_s\) and \(T=F_*\partial_t\), direct use of the
coordinate formula gives
\[
 D_sT-D_tV=\mathcal T(V,T).
\]
Indeed the second ordinary derivatives of each coordinate of \(F\)
cancel and the remaining coefficients are the antisymmetric part
of the Christoffel symbols. This statement allows \(F\) to have
nonconstant rank. Compare \cite[Chapter 4]{LeeRM}.

\section{The Levi--Civita connection}\label{p202:G05}

\begin{theorem}[Levi--Civita and the Koszul formula]
\label{thm:p202-levi-civita}
A Riemannian metric determines a unique torsion-free metric connection.
Its formula is
\begin{align*}
 2\langle\nabla_XY,Z\rangle={}&
 X\langle Y,Z\rangle+Y\langle Z,X\rangle-Z\langle X,Y\rangle\\
 &-\langle X,[Y,Z]\rangle+\langle Y,[Z,X]\rangle
     +\langle Z,[X,Y]\rangle .
\end{align*}
In coordinates
\[
 \Gamma^k_{ij}=\frac12\sum_{\ell}g^{k\ell}
       \bigl(\partial_i g_{j\ell}+\partial_j g_{i\ell}
                          -\partial_\ell g_{ij}\bigr).
\]
A local isometry intertwines the Levi--Civita connections.
\end{theorem}

\begin{proof}
Write the three metric-compatibility identities with cyclically
permuted fields. Add the first two and subtract the third, then
replace each difference \(\nabla_XY-\nabla_YX\) by the bracket.
This gives the Koszul formula and proves uniqueness.

For existence let \(B(X,Y,Z)\) be its right-hand side.
In \(B(fX,Y,Z)\) the terms in \(Y(f)\) from the derivative
and \([fX,Y]\) cancel, as do the terms in \(Z(f)\) from
the derivative and \([Z,fX]\). In \(B(X,Y,fZ)\) the
\(X(f)\) and \(Y(f)\) terms likewise cancel. In \(B(X,fY,Z)\)
the \(Z(f)\) terms cancel and two copies of
\(X(f)\langle Y,Z\rangle\) remain. These expansions give
\(C^\infty\)-linearity in the first and last variables and
the correct Leibniz rule in the second. The metric now defines a
unique vector \(\nabla_XY\) by duality at every point. It is smooth,
as is seen in a frame by multiplication with the smooth inverse
metric matrix. Swapping \(X,Y\) in \(B\) and taking the difference
gives \(2\langle[X,Y],Z\rangle\). Adding \(B(X,Y,Z)\)
and \(B(X,Z,Y)\) gives \(2X\langle Y,Z\rangle\).
Thus the connection has both required properties.
On coordinate vector fields all brackets vanish, giving the
displayed Christoffel formula. Finally a local isometry carries
metric and brackets to those of the target in local coordinates,
so it intertwines the unique connections.
\end{proof}

\section{Parallel transport on compact intervals}\label{p202:G06}

\begin{theorem}[Parallel transport and its inverse]
\label{thm:p202-parallel}
Let \(\gamma:[a,b]\to M\) be piecewise \(C^1\), and let \(E\to M\)
have a smooth connection. Given a vector at any parameter value,
there is a unique continuous field that is \(C^1\) and parallel on
each of the smooth pieces and takes the prescribed value.
Transport \(P_{st}:E_{\gamma(s)}\to E_{\gamma(t)}\) is linear,
\(P_{tu}P_{st}=P_{su}\), and \(P_{ts}=P_{st}^{-1}\).
For a metric connection it preserves the inner product.
\end{theorem}

\begin{proof}
Refine the finite subdivision of \([a,b]\) so each piece is mapped
into one trivialization. Such a refinement exists by compactness
and a Lebesgue number of the pulled-back open cover. In that frame
parallel transport solves \(v'=-A(t)v\) with continuous coefficients
on the closed piece. The linear ODE theorem gives existence and
uniqueness with data at either end, including one-sided derivatives.
One can also obtain it directly from the Volterra integral equation
by iteration: if \(\lVert A\rVert\leq M\), the \(k\)-th iterate of
the homogeneous integral operator has supremum norm at most
\(M^k\lvert b-a\rvert^k/k!\).
The solutions are independent of frame by
Proposition~\ref{prop:p202-pullback-connection}.
Pass the final value of each piece as the initial value of the next,
in either direction. Uniqueness gives the composition and inverse
rules, and linearity follows by superposition.
For parallel fields \(U,V\) the derivative of their pairing is zero
on each piece, and continuity gives the same constant throughout.
\end{proof}

See \cite[Theorems 4.31--4.35 and the author's corrections]{LeeRM},
\cite[Chapter 2, Section 2]{dC}, and \cite[Chapter 7]{D13}.
Only the field and its one-sided values have to agree at a corner.
The two directional derivatives use the actual two path velocities.

\section{The geodesic initial-value problem}\label{p202:G07}

\begin{theorem}[Geodesic spray and smooth dependence]
\label{thm:p202-geodesic-ivp}
For each \(p\in M\) and \(v\in T_pM\) there is a unique maximal
geodesic \(\gamma_v\) with \(\gamma_v(0)=p,\ \dot\gamma_v(0)=v\).
The position and velocity depend smoothly on \((t,v)\) on their
open domain of existence. An affine change of parameter gives a
geodesic, with the corresponding scaled initial velocity.
A Levi--Civita geodesic has constant speed.
\end{theorem}

\begin{proof}
The equation \(D_t\dot\gamma=0\) is
\[
 \dot x^i=v^i,\qquad \dot v^i=-\Gamma^i_{jk}(x)v^jv^k
\]
with repeated indices summed. A solution in a chart represents a
curve with vanishing covariant acceleration by
Proposition~\ref{prop:p202-pullback-connection}. Consequently its
coordinate transforms satisfy the same first-order system in every
overlapping chart. Equivalently the displayed system defines a
smooth vector field on \(TM\), the geodesic spray. Existence,
uniqueness, continuation, and smooth dependence of a smooth local
flow give the assertions on the actual maximal flow domain
\cite[Chapter 6]{P104}\cite[Chapters 3--6]{D13}. The domain is
open because each compact segment of a solution is covered by
finitely many local flow neighborhoods. The covariant curve
formula gives acceleration multiplied by the square of the
constant scaling of the parameter. Metric compatibility gives
\((\lvert\dot\gamma\rvert_g^2)'=2\langle D_t\dot\gamma,\dot\gamma\rangle=0\).
The zero vector gives the constant solution.
\end{proof}

\section{Exponential maps and normal coordinates}\label{p202:G08}

\begin{proposition}[The exponential map near the zero section]
\label{prop:p202-exp-local}
The map \(\exp_p(v)=\gamma_v(1)\) is smooth on its open domain
in \(TM\), contains the zero section in that domain, and is
star-shaped on each fiber domain with respect to the zero vector.
In particular \(d(\exp_p)_0=\id_{T_pM}\), and \(\exp_p\) gives
normal coordinates in a neighborhood of \(p\).
In normal coordinates from an orthonormal basis the metric at \(p\)
is Euclidean and the Christoffel symbols and first derivatives of
the metric vanish at \(p\).
\end{proposition}

\begin{proof}
Use Theorem~\ref{thm:p202-geodesic-ivp} at time \(1\).
If \(v\) survives to time \(1\), then
\(\gamma_{s v}(t)=\gamma_v(st)\) for \(0\leq s\leq1\)
on the interval at issue, so \(sv\) survives as well.
For \(v\in T_pM\) the derivative at \(s=0\) of
\(\exp_p(sv)=\gamma_v(s)\) is \(v\). Apply the inverse function
theorem at \(0\). A sufficiently small open Euclidean ball in the
fiber then maps diffeomorphically onto a normal neighborhood.
In these coordinates the curves \(t\mapsto tv\) are geodesics.
At \(t=0\) this says \(\Gamma^i_{jk}(p)v^jv^k=0\) for all \(v\).
The lower indices are symmetric, hence polarization gives
\(\Gamma^i_{jk}(p)=0\). Metric compatibility in these coordinates
gives \(\partial_i g_{jk}(p)=0\).
\end{proof}

\section{Volume, divergence, and the scalar Laplacian}\label{p202:G02}

The Riemannian density is \(dV_g=\rho\,\lvert dx^1\cdots dx^n\rvert\)
in any chart, where \(\rho=\sqrt{\det(g_{ij})}\). Under a coordinate
change the metric is transformed by the Jacobian and its transpose,
so the determinant square root changes by the absolute Jacobian.
In oriented charts the analogous signed \(n\)-form is the volume
form. These are the density and form constructions of
\cite[Chapters 9--10]{P104}.

\begin{theorem}[Scalar intrinsic differential and integral formulas]
\label{thm:p202-green}
For a smooth function \(u\), let \(g(\operatorname{grad}_g u,X)=du(X)\)
and set \(\operatorname{Hess}_g u(X,Y)=g(\nabla_X\operatorname{grad}_g u,Y)\).
For a vector field \(X\) put \(\operatorname{div}_g X=\operatorname{tr}(Y\mapsto\nabla_Y X)\).
Then
\[
 \operatorname{div}_g X=\rho^{-1}\partial_i(\rho X^i),
 \qquad
 \Delta_g u=\operatorname{tr}_g\operatorname{Hess}_g u
           =\rho^{-1}\partial_i(\rho\,g^{ij}\partial_j u).
\]
On a compact manifold with smooth boundary and outward unit normal
\(\nu\), the divergence theorem holds for the density, with flux
\(\int_{\partial M}\langle X,\nu\rangle\,dS_g\).
For smooth \(u,v\) it yields
\[
 \int_M\bigl(\langle\operatorname{grad}_g u,\operatorname{grad}_g v\rangle
                      +v\Delta_g u\bigr)\,dV_g
   =\int_{\partial M} v\,\langle\operatorname{grad}_g u,\nu\rangle\,dS_g .
\]
On a manifold without boundary the analogous formula has zero
right-hand side when the vector field being integrated has compact
support.
\end{theorem}

\begin{proof}
Metric compatibility gives
\(\operatorname{Hess}_g u(X,Y)=X(Yu)-(\nabla_XY)u\).
Its antisymmetric part is zero by the torsion identity. Tracing
\(\nabla\operatorname{grad}_g u\) gives the Hessian trace.
In coordinates, the Christoffel formula gives
\(\Gamma^i_{ij}=\tfrac12g^{k\ell}\partial_jg_{k\ell}
=\partial_j\log\rho\); the other two terms cancel by interchanging
the summed indices. This proves the divergence formula. Substituting
\((\operatorname{grad}_g u)^i=g^{ij}\partial_j u\) gives the Laplacian.

On an oriented chart
\(d(\iota_X dV_g)=(\operatorname{div}_g X)dV_g\) by direct
differentiation. For an oriented manifold apply Stokes, using the
outward-first boundary orientation. In a boundary frame the pullback
of \(\iota_X dV_g\) is \(\langle X,\nu\rangle\,dS_g\).
In the nonorientable case one can instead sum this argument for
the fields \(\chi_j X\) supported in oriented coordinate patches of a
partition of unity. The terms involving \(d\chi_j\) sum to zero.
The same proof works with compact support in the boundaryless case.
Finally use the product rule with \(X=v\operatorname{grad}_g u\).
Swapping \(u,v\) and subtracting also gives the second Green identity.
\end{proof}

The sign agrees with the classical Green identity in \cite{P201}.
For the coordinate formulas see \cite[Chapters 2 and 5]{LeeRM}.

\section{First variation and corners}\label{p202:G10}

For a smooth variation \(F(s,t)\) of a curve on \([a,b]\), let
\(T=\partial_tF\), \(V=\partial_sF\), and evaluate at \(s=0\) when
writing the varied curve \(\gamma\) and its variation field \(V\).

\begin{theorem}[The first variation of energy and length]
\label{thm:p202-first-variation}
For a smooth segment the first variations are
\begin{align*}
 E'(0)&=\bigl[\langle V,\dot\gamma\rangle\bigr]_a^b
                  -\int_a^b\langle V,D_t\dot\gamma\rangle\,dt,\\
 L'(0)&=\biggl[\biggl\langle V,\frac{\dot\gamma}{\lvert\dot\gamma\rvert_g}
                         \biggr\rangle\biggr]_a^b
          -\int_a^b\biggl\langle V,
                   D_t\frac{\dot\gamma}{\lvert\dot\gamma\rvert_g}
                           \biggr\rangle\,dt .
\end{align*}
The length formula requires positive speed on the segment.
With endpoints fixed, energy criticality among all smooth
variations is equivalent to being a geodesic. Length criticality
for regular curves means being a geodesic after arc-length
reparametrization. A concatenation of nonconstant geodesic
segments that minimizes length has matching unit velocities at
its joins.
\end{theorem}

\begin{proof}
The torsion-free identity of Section~\ref{p202:G04} gives
\(D_sT=D_tV\). Differentiation under the integral and metric
compatibility therefore give \(\int\langle D_tV,T\rangle\)
for the energy, and \(\int\langle D_tV,T/\lvert T\rvert_g\rangle\)
for the length. Positive speed on a compact segment justifies
differentiating its norm. Covariant integration by parts proves
both formulas. With fixed endpoints one can test against any
compactly supported smooth field along the segment. For example,
the variation \(\exp_{\gamma(t)}(sV(t))\) is defined for uniformly
small \(s\) by compactness and Section~\ref{p202:G08}.
The fundamental lemma of the calculus of variations then gives
\(D_tT=0\) or \(D_t(T/\lvert T\rvert_g)=0\), respectively.
Conversely those equations make the first variations vanish.
No division by speed was used for the energy, so constant curves
are included.

Apply the length formula separately on each regular segment of a
piecewise variation. At a join the boundary terms sum to the
pairing of the freely chosen variation value with the left unit
velocity minus the right one. Minimality makes that pairing zero.
The velocities agree, so after unit-speed parametrization the
geodesic initial-value uniqueness glues the two pieces into one
smooth segment. The same segmentwise formulas display the
energy boundary term at a corner as the difference of the full
velocities. For moving endpoints their terms at \(a,b\) remain.
\end{proof}

For the relation between energy and length, Cauchy--Schwarz gives
\(L(\gamma)^2\leq 2(b-a)E(\gamma)\) with equality at constant
speed. A regular piecewise smooth minimizer can be parametrized by
arc length, and varying a join as above then shows it has no corner.
A stationary curve in the first variation need only be a local
candidate for minimization. See \cite[Chapter 6]{LeeRM},
\cite[Chapter 9, Section 2]{dC}, and \cite[Chapter 7]{P106}.

\section{The Gauss lemma and convex normal neighborhoods}
\label{p202:G09}

\begin{theorem}[Gauss lemma and short minimizing geodesics]
\label{thm:p202-gauss-normal}
On the domain of \(\exp_p\) one has
\[
 \bigl\langle d(\exp_p)_v(v),d(\exp_p)_v(w)\bigr\rangle
   =\langle v,w\rangle_p .
\]
On a sufficiently small normal ball about \(p\) the radial
geodesic from \(p\) to each point is the unique length-minimizing
curve up to increasing reparametrization and pauses.
Each point has arbitrarily small neighborhoods in which every
pair of points is joined by a unique minimizing geodesic that
stays in the neighborhood. The geodesic depends smoothly on the
pair of endpoints, and the squared distance is smooth there.
\end{theorem}

\begin{proof}
Set \(F(s,t)=\exp_p(t(v+s w))\) for \(0\leq t\leq1\) and
small \(s\). Each \(t\)-curve has constant squared speed
\(\lvert v+s w\rvert^2\) and the initial point is independent
of \(s\). For \(T=\partial_tF,\ V=\partial_sF\), metric compatibility
and torsion freeness give
\[
 \partial_t\langle V,T\rangle
 =\langle D_sT,T\rangle
 =\tfrac12\partial_s\lvert v+s w\rvert^2 .
\]
Integrate from \(0\) to \(1\) at \(s=0\).
This proves the formula, including orthogonality of the unit radial
direction to angular directions and preservation of radial length.

Choose \(r_0>0\) so \(\exp_p\) is a diffeomorphism on a larger
fiber ball than the closed ball of radius \(r_0\).
In this chart the speed of any piecewise \(C^1\) competitor
away from \(p\) is at least the absolute derivative of its radial
coordinate. The radial coordinate along the whole curve is
absolutely continuous. At zeros it has derivative zero almost
everywhere, so the length bound also holds across \(p\).
For an endpoint at radius \(r<r_0\), any competitor leaving
the normal ball has length at least \(r_0\) before its first exit.
Every other competitor has length at least \(r\). Equality
requires that the radius be nondecreasing and that the angular
derivative vanish wherever the radius is positive, by strict
positivity of the metric on angular directions. On intervals
with positive radius this forces a single direction. Thus
the unique minimizing image is the radial segment and the
parametrizations differ only by speed and pauses.

For the uniform statement consider
\(\Phi(q,v)=(q,\exp_q v)\) near \((p,0)\).
Its differential in a bundle chart sends
\((\delta q,\delta v)\) to \((\delta q,\delta q+\delta v)\)
at \((p,0)\). The inverse function theorem, with shrinking of
the base neighborhood, supplies a fixed radius \(r_0\) for
which \(\Phi\) is one-to-one and a local diffeomorphism on
all closed fiber balls of radius \(r_0\) over a neighborhood
of \(p\), with open domain containing those balls.
Indeed one first takes a product neighborhood inside the
inverse-function neighborhood, and then uses local comparison
of the fiber norm to the trivialization norm. Shrink the base
again so all pairs \((q,z)\) in the smaller neighborhood are
in the image and their inverse vectors have norm less than
\(r_0/2\); continuity of the local inverse at \((p,p)\) gives this.
The preceding exit argument, now centered at \(q\), gives unique
global minimizing radial segments for those pairs. Their formula
using the inverse of \(\Phi\) gives smooth dependence and
\(d_g(q,z)^2=\lvert\Phi^{-1}(q,z)_2\rvert_{g_q}^2\).

Finally \(z\mapsto d_g(p,z)^2\) has Hessian \(2g_p\) at \(p\)
by normal coordinates, hence positive Hessian on a sufficiently
small neighborhood of \(p\). Choose \(\epsilon>0\) with the
metric ball of radius \(3\epsilon\) contained there and with
\(B_g(p,\epsilon)\) in the uniform neighborhood just constructed.
For endpoints in \(B_g(p,\epsilon)\), the minimizing segment has
length less than \(2\epsilon\), hence stays in \(B_g(p,3\epsilon)\)
by the triangle inequality. The squared distance to \(p\) has
nonnegative second derivative along this geodesic. It is therefore
bounded above by the affine interpolation of the two endpoint
values, which are strictly less than \(\epsilon^2\).
This proves that the segment stays in \(B_g(p,\epsilon)\)
and gives the required convex normal neighborhoods.
\end{proof}

The distinction between normality at one center and the convex
neighborhood property is used below. The argument follows
\cite[Chapter 6, including the correction to Theorem 6.15]{LeeRM};
see also \cite[Chapter 3]{dC}.

\section{Completeness and Hopf--Rinow}\label{p202:G11}

\begin{theorem}[Hopf--Rinow]
\label{thm:p202-hopf-rinow}
Let \(M\) be a connected finite-dimensional Riemannian manifold
without boundary. The following are equivalent.
\begin{enumerate}
\item The length metric \(d_g\) is complete.
\item Every maximal geodesic is defined on the whole real line.
\item Every closed and bounded subset is compact.
\end{enumerate}
Under these conditions every pair of points is joined by a
geodesic of length equal to their distance. Completeness follows
even if \(\exp_p\) is defined on all of \(T_pM\) at a single point.
\end{theorem}

\begin{proof}
First suppose the metric is complete and a geodesic has a finite
upper endpoint \(b\) of its maximal interval. Constant speed gives
a Cauchy position as \(t\uparrow b\), hence a point \(q\in M\).
Near \(b\) use a chart on a neighborhood of a compact coordinate
ball about \(q\) that contains the remaining arc. The lower bound
of the metric norm there bounds the coordinate velocity, and
the geodesic equation with bounded Christoffel symbols bounds
the coordinate acceleration. The velocity is therefore Cauchy
as \(t\uparrow b\). The position and velocity have a limit in
\(TM\). The coordinate solution satisfies the integral equation
up to the limiting state by continuity. Local existence at that
state and uniqueness backward in time therefore extend the
original geodesic.
A finite lower endpoint is excluded by reversal.

For the converse and the minimizing assertion, fix \(p\) such
that \(\exp_p\) is defined on the whole fiber. Fix \(q\ne p\)
and put \(d=d_g(p,q)>0\). For a sufficiently small normal
radius \(0<\epsilon<d\), the sphere
\(S=\exp_p(\{\lvert v\rvert=\epsilon\})\) is compact. Every curve
from \(p\) to \(q\) crosses it, and the initial portion has
length at least \(\epsilon\). Taking a sequence of curves with
lengths tending to \(d\), and using the triangle inequality for
the reverse bound, gives
\[
 d=\epsilon+\min_{y\in S}d_g(y,q).
\]
Choose a minimizing \(y\) and let \(\gamma\) be the unit-speed
geodesic from \(p\) in its radial direction, defined for all
nonnegative times by the assumption at \(p\). The closed subset
of \([\epsilon,d]\) on which
\(d_g(\gamma(t),q)=d-t\) contains \(\epsilon\); let \(t_0\) be
its maximum. If \(t_0<d\), center a small normal sphere of
radius \(0<\delta<d-t_0\) at \(x=\gamma(t_0)\) and repeat the
sphere argument to find \(z\) on it with
\(d_g(z,q)=d-t_0-\delta\).
The two inequalities
\[
 d_g(p,z)\le t_0+\delta,\qquad
 d\le d_g(p,z)+d_g(z,q)
\]
force equality in the first one. The concatenation of the
original geodesic segment to \(x\) and the short radial segment
from \(x\) to \(z\) is thus length minimizing. By the no-corner
conclusion of Theorem~\ref{thm:p202-first-variation}, their unit
velocities agree at \(x\). Geodesic uniqueness gives
\(z=\gamma(t_0+\delta)\), contradicting the choice of \(t_0\).
Therefore \(\gamma(d)=q\) and the segment is minimizing.

It follows that the closed metric ball of radius \(R\) about \(p\)
is exactly the image under \(\exp_p\) of the closed fiber ball
of that radius. Inclusion in one direction uses the length of a
radial path, and the other uses the minimizing segment just proved.
The fiber ball is compact in finite dimension. Any closed bounded
subset is closed in such a compact ball and is compact.
Finally every Cauchy sequence is bounded, so the third condition
gives a convergent subsequence and then a limit for the whole
sequence. If all geodesics are complete the fiber condition
holds at every center; if it holds at a single center the
argument has already proved metric completeness. This proves all
the equivalences and the claim for every pair.
\end{proof}

There is no uniqueness assertion for arbitrary distant endpoints.
The normal-neighborhood result is local, and conjugate and cut
phenomena describe the possible loss of its conclusion.
Compare \cite[Chapter 6]{LeeRM} and \cite[Chapter 7]{dC}.

\chapter{Curvature and Isometric Immersions}\label{chap:p202-curvature}

\section{The curvature tensor and Bianchi identities}\label{p202:G12}

Write \(r(X,Y,Z,W)=\langle R(X,Y)Z,W\rangle\).
The convention for \(R\) is that fixed in the introduction.

\begin{theorem}[Curvature symmetries and commutators]
\label{thm:p202-curvature-symmetry}
The expression \(R\) is tensorial in its three arguments. The
four-tensor \(r\) satisfies
\[
 r(X,Y,Z,W)=-r(Y,X,Z,W)=-r(X,Y,W,Z),\qquad
 r(X,Y,Z,W)=r(Z,W,X,Y).
\]
The two Bianchi identities are
\begin{align*}
 R(X,Y)Z+R(Y,Z)X+R(Z,X)Y&=0,\\
 (\nabla_XR)(Y,Z)+(\nabla_YR)(Z,X)+(\nabla_ZR)(X,Y)&=0 .
\end{align*}
For an induced tensor connection, the operator
\(\nabla_X\nabla_Y-\nabla_Y\nabla_X-\nabla_{[X,Y]}\)
acts on a tensor as \(R(X,Y)\) in each vector factor and its
negative transpose in each covector factor.
\end{theorem}

\begin{proof}
The connection product rule and the Lie bracket product rule
cancel the derivative terms when \(X\) or \(Y\) is multiplied
by a function. For the third argument the remaining second-order
term is \(X(Yf)-Y(Xf)-[X,Y]f=0\); the first-order mixed terms
cancel in pairs. Thus \(R\) is a tensor and is manifestly
skew in \(X,Y\). Since \(\nabla g=0\), applying the commutator
to \(\langle Z,W\rangle\) gives
\(r(X,Y,Z,W)+r(X,Y,W,Z)=0\).

For the remaining algebraic identities take normal coordinates
at an arbitrary point. Set
\(r_{ijkl}=\langle R(\partial_i,\partial_j)\partial_k,\partial_l\rangle\)
there. The Christoffel symbols and first metric derivatives vanish
at the point. The coordinate curvature is
\(R^l{}_{kij}=\partial_i\Gamma^l_{jk}-\partial_j\Gamma^l_{ik}\)
there, hence the Christoffel formula gives
\[
 r_{ijkl}=\tfrac12
   \bigl(\partial_i\partial_k g_{jl}-\partial_i\partial_l g_{jk}
          -\partial_j\partial_k g_{il}+\partial_j\partial_l g_{ik}\bigr).
\]
This is invariant under interchanging the two ordered pairs.
Cyclically permuting \(i,j,k\) makes every term cancel with its
opposite by commutation of partial derivatives and symmetry of \(g\).
This proves the first Bianchi identity.

For the second Bianchi identity keep the same coordinates.
In \(\partial_m R^l{}_{kij}\) at the point the derivative of every
quadratic Christoffel term vanishes. A covariant derivative there
is the ordinary derivative of the components. Thus the cyclic sum
in \(m,i,j\) is the sum of differences
\(\partial_m\partial_i\Gamma^l_{jk}-\partial_m\partial_j\Gamma^l_{ik}\)
and vanishes by commutation. Both statements hold at every point
by tensoriality.

The tensor product Leibniz rule gives the commutator on a product:
the mixed derivatives cancel and each factor contributes its own
commutator. For a dual factor the scalar contraction has zero
commutator, which forces the negative transpose rule. This
also gives the Ricci commutation rule on tensors without a choice
of frame.
\end{proof}

The algebraic and differential identities are compared with
\cite[Chapter 7]{LeeRM} and \cite[Chapter 4]{dC}.

\section{Sectional, Ricci, and scalar curvature}\label{p202:G13}

\begin{theorem}[Recovery and contraction of curvature]
\label{thm:p202-sectional-recovery}
At a point, the sectional curvatures determine \(r\).
In an orthonormal frame \((e_i)\) define
\(\operatorname{Ric}(Y,Z)=\sum_i r(e_i,Y,Z,e_i)\) and
\(\operatorname{Scal}=\operatorname{tr}_g\operatorname{Ric}\).
Then \(\operatorname{Ric}\) is symmetric and
\[
   \nabla^i\operatorname{Ric}_{ij}=\tfrac12\partial_j\operatorname{Scal}.
\]
On a connected manifold of dimension \(n\ge3\), if all sectional
curvatures at each point have a common value, that value is constant.
\end{theorem}

\begin{proof}
Let \(S\) be the difference of two tensors with the curvature
symmetries and the same sectional curvatures.
Then \(S(X,Y,Y,X)=0\) for all \(X,Y\), including linearly dependent
ones by skewness. Polarizing in \(X\) gives
\(S(X,Y,Y,W)+S(W,Y,Y,X)=0\). Pair interchange and the two
skew symmetries make the two terms equal. Hence
\(S(X,Y,Y,W)=0\). Polarizing this in \(Y\) shows skewness in
the second and third entries. Combined with skewness in the first
two entries, this makes \(S\) alternating in its first three
entries. Their cyclic sum is \(3S\), and the first Bianchi identity
forces \(S=0\).

The Ricci tensor is independent of frame by contraction, and pair
interchange and skewness show its symmetry. In a normal orthonormal
coordinate frame at a point, the second Bianchi identity, summed
over \(i,m\), gives
\[
 \sum_{i,m}\bigl(\nabla_m r_{ijmi}
           +\nabla_i r_{jmmi}+\nabla_j r_{mimi}\bigr)=0 .
\]
The first sum is the divergence of Ricci with index \(j\).
Pair interchange turns the second into the same divergence.
Skewness in the last two entries turns the last sum into
\(-\partial_j\operatorname{Scal}\).
This proves the contracted identity.
If the common sectional value is \(\kappa(p)\), recovery gives
\[
 r(X,Y,Z,W)=\kappa(p)
     \bigl(\langle Y,Z\rangle\langle X,W\rangle
                         -\langle X,Z\rangle\langle Y,W\rangle\bigr).
\]
Consequently \(\operatorname{Ric}=(n-1)\kappa g\) and
\(\operatorname{Scal}=n(n-1)\kappa\). The function \(\kappa\) is
smooth by the scalar formula. Contracted Bianchi gives
\((n-1)(1-n/2)d\kappa=0\), so \(d\kappa=0\) for \(n\ge3\).
Connectedness makes the value constant.
In dimension two there is only one sectional plane at a point
and the same conclusion about constancy does not follow.
\end{proof}

\begin{theorem}[Local Euclidean geometry and flatness]
\label{thm:p202-flat}
A Riemannian metric has zero curvature on an open set if and only
if every point there has a neighborhood isometric to a Euclidean
open set.
\end{theorem}

\begin{proof}
A local isometry intertwines the connection and its curvature,
so the Euclidean direction is immediate.
For the converse, work on a sufficiently small coordinate box
and use the matrices \(\Gamma_i\) of the connection.
Writing \(\nabla_i=\partial_i+\Gamma_i\), the zero-curvature
condition says
\[
 \partial_i\Gamma_j-\partial_j\Gamma_i+
                    \Gamma_i\Gamma_j-\Gamma_j\Gamma_i=0 .
\]
We construct a matrix \(P\) whose columns are parallel.
In one coordinate direction this is the linear initial-value
problem \(\partial_1P=-\Gamma_1P\).
Inductively solve \(\partial_iP=-\Gamma_iP\) for \(i<n\)
on the face \(x^n=0\), with the prescribed matrix at the origin,
then solve \(\partial_nP=-\Gamma_nP\) off that face.
The solutions are smooth in the remaining coordinates by smooth
dependence of linear ODEs, and the matrices stay invertible by
uniqueness, or by also solving the inverse matrix equation.
For \(i<n\) the residual \(Q_i=\partial_iP+\Gamma_iP\) satisfies
\[
 \partial_n Q_i=-\Gamma_n Q_i+
   \bigl(\partial_n\Gamma_i-\partial_i\Gamma_n+
                           [\Gamma_n,\Gamma_i]\bigr)P .
\]
Its initial value on the face is zero, so curvature zero and ODE
uniqueness give \(Q_i=0\). The columns therefore form a parallel
frame. If the initial frame is orthonormal, metric compatibility
keeps it orthonormal on the connected box. Torsion freeness makes
its fields commute. For the dual coframe \((\theta^i)\) this implies
\(d\theta^i=0\), by evaluating on frame fields.
On a smaller coordinate box, each closed one-form is \(dy^i\):
the primitive is obtained by integrating along straight segments
from the origin, and differentiation using closedness verifies its
derivative. The inverse function theorem applies to \((y^i)\)
since their differentials form a coframe. The metric becomes
\(\sum_i(dy^i)^2\), as required.
\end{proof}

See \cite[Theorems 7.10 and 7.12--7.19]{LeeRM}.

\section{The tangent and normal connections of an immersion}
\label{p202:G14}

Let \(F:M\to\overline M\) be an isometric immersion into a
Riemannian manifold. Identify \(TM\) with its image in
\(F^*T\overline M\) and write \(NM\) for the orthogonal complement.
Use the pullback of the ambient Levi--Civita connection, denoted
\(\overline\nabla\). The calculations are local and apply to
immersions with self-intersection in the ambient space.

\begin{proposition}[Gauss and Weingarten decompositions]
\label{prop:p202-gauss-weingarten}
For tangent fields \(X,Y\) and a normal field \(\xi\) one has
\[
 \overline\nabla_XY=\nabla_XY+B(X,Y),\qquad
 \overline\nabla_X\xi=-A_\xi X+\nabla^\perp_X\xi ,
 \qquad
 \langle A_\xi X,Y\rangle=\langle B(X,Y),\xi\rangle .
\]
Here \(\nabla\) is the intrinsic Levi--Civita connection,
\(B\) is a symmetric normal-valued tensor, and \(\nabla^\perp\)
is a metric connection on \(NM\).
\end{proposition}

\begin{proof}
Project the ambient derivative onto the two summands.
The tangential projection on tangent fields is a connection
and is metric by orthogonality and ambient metric compatibility.
The coordinate pullback calculation from Section~\ref{p202:G04}
shows
\(\overline\nabla_X(dF\,Y)-\overline\nabla_Y(dF\,X)=dF[X,Y]\).
Its tangential part proves torsion freeness, so the tangential
connection is Levi--Civita. The normal part proves symmetry of
\(B\), and the connection product rule makes \(B\) tensorial
since the derivative of a scalar times a tangent field is tangent.
The normal projection on normal fields is a metric connection by
the same argument. Differentiate \(\langle \xi,Y\rangle=0\)
to identify the tangential projection of
\(\overline\nabla_X\xi\), including its minus sign.
\end{proof}

\section{Gauss, Codazzi, and Ricci equations}\label{p202:G15}

\begin{theorem}[Curvature equations in arbitrary codimension]
\label{thm:p202-gauss-codazzi-ricci}
For the above immersion let \(R^\perp\) be the curvature of the
normal connection. With tangential \(X,Y,Z,W\) and normal
\(\xi,\eta\) one has
\begin{align*}
 r(X,Y,Z,W)
 &=\overline r(X,Y,Z,W)
   +\langle B(X,W),B(Y,Z)\rangle
   -\langle B(Y,W),B(X,Z)\rangle ,\\
 (\overline R(X,Y)Z)^\perp
 &=(\nabla_X B)(Y,Z)-(\nabla_Y B)(X,Z),\\
 \langle R^\perp(X,Y)\xi,\eta\rangle
 &=\langle\overline R(X,Y)\xi,\eta\rangle+
                           \langle[A_\xi,A_\eta]X,Y\rangle .
\end{align*}
In the second formula \(\nabla B\) uses the normal connection
on its value and the tangent connection on its two inputs.
\end{theorem}

\begin{proof}
The curvature of the pullback connection is ambient curvature
evaluated on \(dF X,dF Y\); this follows in a frame by pulling
back \(d\Gamma+\Gamma\wedge\Gamma\).
In \(\overline\nabla_X\overline\nabla_YZ\), the tangential
part is \(\nabla_X\nabla_Y Z-A_{B(Y,Z)}X\).
Subtract the corresponding term with \(X,Y\) switched and the
bracket term. Pairing with \(W\) gives the Gauss equation.
The normal part before subtraction is
\(B(X,\nabla_YZ)+\nabla_X^\perp(B(Y,Z))\).
After subtraction, the bracket term is \(B([X,Y],Z)\).
Replace \([X,Y]\) by \(\nabla_XY-\nabla_YX\); grouping then
gives the asserted two tensor derivatives of \(B\).

For a normal field \(\xi\), the normal part of
\(\overline\nabla_X\overline\nabla_Y\xi\) is
\(-B(X,A_\xi Y)+\nabla_X^\perp\nabla_Y^\perp\xi\).
Subtracting and pairing with \(\eta\), the two \(B\)-terms
give
\[
 -\langle A_\eta X,A_\xi Y\rangle+
        \langle A_\eta Y,A_\xi X\rangle
    =-\langle[A_\xi,A_\eta]X,Y\rangle,
\]
since the shape operators are self-adjoint. Moving this term
to the other side gives the Ricci equation.
\end{proof}

\begin{corollary}[The intrinsic meaning of Gaussian curvature]
\label{cor:p202-egregium}
For an immersed surface in Euclidean three-space, with a local
unit normal and its shape operator \(A\), the intrinsic Gaussian
curvature is \(\det A\). In any parametrization it is the
determinant of the second form divided by the determinant of the
first form. In particular it is invariant under local isometry
of the induced metrics.
\end{corollary}

\begin{proof}
Set \(X,W=e_1,\ Y,Z=e_2\) in the Gauss equation for an orthonormal
tangent basis. Ambient curvature is zero and the right-hand side
is \(\langle A e_1,e_1\rangle\langle A e_2,e_2\rangle-
\langle A e_1,e_2\rangle^2\), the determinant of \(A\).
In any basis the shape matrix is \(g^{-1}B\) with \(B\) interpreted
as the scalar second form in that normal. This gives the ratio.
The left side depends only on the metric and its Levi--Civita
connection. Reversing the normal changes \(A\) to \(-A\) and does
not change its two-dimensional determinant.
\end{proof}

\begin{proposition}[Mean curvature and first variation of volume]
\label{prop:p202-volume-variation}
Let \(F_s\) be a smooth variation through immersions of a compact
manifold, with a smooth boundary if present, and use the induced
metrics. Write \(V=\partial_s F_s|_0=dF(W)+V^\perp\) and
\(H=\operatorname{tr}_g B\). Then
\[
 \frac{d}{ds}\operatorname{Vol}(M,F_s^*\overline g)\bigg|_{s=0}
 =-\int_M\langle V^\perp,H\rangle\,dV_g+
                    \int_{\partial M}\langle W,\nu\rangle\,dS_g .
\]
The same formula for a local change of volume holds for a compactly
supported variation of a noncompact immersion.
\end{proposition}

\begin{proof}
Use a fixed coordinate frame on \(M\). Differentiating
\(g_{ij}(s)=\langle\partial_iF_s,\partial_jF_s\rangle\) with
the ambient metric connection gives
\(\dot g_{ij}=\langle\overline\nabla_iV,\partial_jF\rangle+
\langle\partial_iF,\overline\nabla_jV\rangle\), using
torsion freeness along the variation.
The determinant differentiation formula
\(\dot\rho/\rho=\tfrac12\operatorname{tr}(g^{-1}\dot g)\)
therefore gives the sum of
\(\langle\overline\nabla_{e_i}V,e_i\rangle\)
over an orthonormal basis.
By Proposition~\ref{prop:p202-gauss-weingarten} this is
\(\operatorname{div}_g W-\langle V^\perp,H\rangle\).
Integrate and apply the divergence theorem
(Theorem~\ref{thm:p202-green}). For compact support, differentiate
the difference of the densities, which vanishes outside a compact
set, and apply the same divergence formula to the compactly
supported field \(W\). Only the original boundary can contribute.
\end{proof}

Here the mean curvature is the trace, with no division by dimension.
See \cite[Chapter 8, including the volume variation]{LeeRM}
and \cite[Chapter 6]{dC}.

\section{Euclidean, spherical, and hyperbolic models}\label{p202:G03}

\begin{theorem}[The three simply connected curvature models]
\label{thm:p202-models}
For \(n\ge2\), Euclidean space, the sphere \(S^n_a\) of radius
\(a>0\), and the upper sheet
\[
 \mathbb H^n_a=\{(x,\tau)\in\R^n\times\R:
           \lvert x\rvert^2-\tau^2=-a^2,\ \tau>0\}
\]
with the induced metric from \(\langle(x,\tau),(y,\sigma)\rangle_L
=x\cdot y-\tau\sigma\), are complete and simply connected. Their
sectional curvatures are \(0,\ a^{-2},\ -a^{-2}\), respectively.
For a unit tangent vector \(u\) at \(p\), the unit-speed geodesic
with this initial data is
\[
 \begin{cases}
 p+t u &\text{in Euclidean space},\\
 \cos(t/a)p+a\sin(t/a)u &\text{on the sphere},\\
 \cosh(t/a)p+a\sinh(t/a)u &\text{on the hyperboloid}.
 \end{cases}
\]
The upper half-space with coordinates \(y\in\R^{n-1}\), \(z>0\)
and metric \(a^2z^{-2}(dy^2+dz^2)\) is isometric to the hyperboloid.
\end{theorem}

\begin{proof}
The Euclidean statements follow from the constant metric.
On the sphere \(N=p/a\) is the unit normal and its derivative in
a tangent direction \(X\) is \(X/a\). The second fundamental form
is \(-\langle X,Y\rangle N/a\). The Gauss equation gives curvature
\(a^{-2}\). The displayed great circle has unit speed, correct
initial data, and normal acceleration, hence is a geodesic by
Proposition~\ref{prop:p202-gauss-weingarten}. It is defined for
all time.

On the hyperboloid write \(p=(x,\tau)\). A tangent vector
\((u,s)\) satisfies \(s=x\cdot u/\tau\). Thus its squared norm
is \(\lvert u\rvert^2-(x\cdot u)^2/\tau^2\ge
(a^2/\tau^2)\lvert u\rvert^2\), so the induced metric is positive.
Here \(N=p/a\) has squared norm \(-1\).
Let \(D\) be the ordinary flat derivative in the ambient
vector space. Its tangential projection on tangent fields
is metric and torsion-free by the same projection argument as
in Proposition~\ref{prop:p202-gauss-weingarten}.
Since \(D_XN=X/a\), differentiating \(\langle Y,N\rangle_L=0\)
gives
\[
 D_XY=\nabla_XY+a^{-1}\langle X,Y\rangle N .
\]
In the tangential part of the flat commutator the two normal
terms contribute
\(a^{-2}(\langle Y,Z\rangle X-\langle X,Z\rangle Y)\).
Hence the intrinsic \(R(X,Y)Z\) is the negative of this
expression, giving curvature \(-a^{-2}\).
The displayed hyperbolic curve has Lorentz squared norm \(-a^2\)
at all times. It lies in the upper sheet by continuity, and
its unit speed and normal acceleration follow by differentiation.
Its tangential acceleration is therefore zero. The formulas and
geodesic uniqueness, including affine changes of speed, give
geodesic completeness and hence metric completeness by
Theorem~\ref{thm:p202-hopf-rinow}.
The sphere is simply connected for \(n\ge2\) by
Section~\ref{p202:T02}, and projection to \(x\) identifies the
upper sheet diffeomorphically with \(\R^n\).

For the half-space isometry use null coordinates and set
\[
 x'=\frac{a y}{z},\qquad
 \tau-x_n=\frac{a}{z},\qquad
 \tau+x_n=\frac{a(\lvert y\rvert^2+z^2)}{z},
\]
where \(x'=(x_1,\ldots,x_{n-1})\).
Then \(\lvert x'\rvert^2-(\tau-x_n)(\tau+x_n)=-a^2\).
Conversely the upper sheet has \(\tau-x_n>0\), so the first two
formulas determine \(z>0\) and \(y\) smoothly and the last
follows from the hyperboloid equation. Substituting the
differentials into
\(dx'\cdot dx'-d(\tau-x_n)d(\tau+x_n)\) cancels the mixed terms
and yields \(a^2z^{-2}(dy^2+dz^2)\).
\end{proof}

The model metrics are compared in \cite[Chapters 3 and 5]{LeeRM}
and \cite[Chapter 8]{dC}. Their radial Jacobi fields and the
classification of complete constant-curvature manifolds are
developed in Sections~\ref{p202:G20} and \ref{p202:G22}.
\chapter{Jacobi Fields, Cut Points, and Comparison}
\label{ch:p202-jacobi-comparison}

Throughout this chapter the Riemannian metric is smooth and
positive definite. The geodesic and completeness results are used
on manifolds without boundary. Our curvature convention is
\[
 R(X,Y)Z=\nabla_X\nabla_YZ-\nabla_Y\nabla_XZ-\nabla_{[X,Y]}Z.
\]
Write \(D_t\) for covariant differentiation along a curve and
\(T=\dot\gamma\) along a geodesic. Then the round sphere has positive
sectional curvature and, for a unit \(T\), \(V\mapsto R(V,T)T\)
is the self-adjoint radial curvature operator on \(T^\perp\).
The energy on an interval \([a,b]\) is
\(E(\gamma)=\frac12\int_a^b|\dot\gamma|^2\,dt\).

\section{The Jacobi equation and geodesic variations}
\label{p202:G16}

\begin{theorem}[Linearized geodesic equation]\label{thm:p202-jacobi-variation}
Let \(\gamma\) be a geodesic. The variation field \(J\) of a smooth
variation through affinely parametrized geodesics satisfies
\[
 D_t^2J+R(J,T)T=0.
\]
A solution, called a Jacobi field, is uniquely specified by
\(J(t_0)\) and \(D_tJ(t_0)\). Conversely, on any compact parameter
segment in the interval of \(\gamma\), each Jacobi field is realized
by a geodesic variation. For a nonconstant geodesic its tangential
component is \((a+bt)T\); the normal component is a Jacobi field
as well.
\end{theorem}

\begin{proof}
For the two-parameter map \(F(s,t)\), put \(S=\partial_sF\)
and \(T_s=\partial_t F\), and use the pullback connection.
Torsion freeness gives \(D_sT_s=D_tS\).
The curvature identity is
\[
 D_sD_t Z-D_tD_s Z=R(S,T_s)Z
\]
for every vector field \(Z\) along the map. Differentiate
\(D_t T_s=0\) in \(s\) and set \(s=0\). This gives the asserted
equation with the stated sign.

In a parallel orthonormal frame along \(\gamma\) the equation is
a linear second-order ordinary differential equation with smooth
coefficients. The first-order system for the field and its
derivative has a unique solution throughout the interval, since
its coefficients are bounded on each compact subinterval.
For the converse, choose a curve \(p(s)\) through
\(p=\gamma(t_0)\) with \(p'(0)=J(t_0)\). Take a smooth vector
\(v(s)\) along it with \(v(0)=T(t_0)\) and
\(D_sv(0)=D_tJ(t_0)\). For example, use parallel transport of
\(v(0)+sD_tJ(t_0)\) along \(p(s)\). The geodesics with initial data
\((p(s),v(s))\) are all defined on the prescribed compact segment
for small \(s\). Indeed, the geodesic flow is smooth on its open
survival domain and that segment times the initial point is
compact. The resulting variation field has the required two
initial values, so it equals \(J\).

Along a nonconstant geodesic \(T\) is parallel and has nonzero
constant norm. Curvature skew-symmetry in the metric pairing
gives \(\langle R(J,T)T,T\rangle=0\), so
\(\langle J,T\rangle''=0\). The projection of \(J\) on \(T\)
is therefore \((a+bt)T\), itself a solution of the Jacobi
equation. Subtracting it leaves the normal solution.
\end{proof}

See \cite[Chapter 10, Proposition 10.4]{LeeRM} and
\cite[Chapter 5, Sections 2--3]{dC}.

\section{The differential of the exponential and conjugate points}
\label{p202:G17}

\begin{theorem}[Exponential differential]\label{thm:p202-exp-jacobi}
Let \(\gamma(t)=\exp_p(tv)\) and \(w\in T_pM\). Wherever
\(\gamma\) is defined, the Jacobi field with initial values
\(J(0)=0\), \(D_tJ(0)=w\) satisfies
\[
 J(t)=d(\exp_p)_{tv}(tw).
\]
For \(t>0\) the space of fields vanishing at both \(0\) and \(t\)
has dimension equal to the nullity of \(d(\exp_p)_{tv}\).
In particular \(\gamma(t)\) is conjugate to \(p\) along
\(\gamma\) exactly when that differential is singular.
\end{theorem}

\begin{proof}
Differentiate \(\exp_p(t(v+sw))\) at \(s=0\), with \(t\) fixed.
At \(t=0\) the base point is independent of \(s\), and
\(D_tJ(0)=D_s\partial_tF(0,0)=w\). Uniqueness from
Theorem~\ref{thm:p202-jacobi-variation} identifies the field.
When \(t>0\), multiplication by \(t\) on \(T_pM\) is an
isomorphism. Thus the displayed formula identifies the kernel
with fields having the two zero values. At \(t=0\) one instead
has \(d(\exp_p)_0=\operatorname{id}\).
A nonzero field vanishing at both ends of a nonconstant geodesic
is necessarily normal, since its affine tangential part has
two zeros.
\end{proof}

The multiplicity here refers to the given geodesic, rather than
to the endpoint as a point reached by all possible geodesics.
See \cite[Chapter 10, Proposition 10.10 and Proposition 10.20]{LeeRM}.

\section{Second variation, the index lemma, and minimality}
\label{p202:G18}

\begin{theorem}[Second variation of energy]\label{thm:p202-second-variation}
Let \(F(s,t)\) be a smooth variation of a geodesic
\(\gamma:[a,b]\to M\), with variation field \(V\).
Put \(S=\partial_sF\). Then
\[
 \begin{aligned}
 \left.\frac{d^2}{ds^2}E(F(s,\cdot))\right|_{s=0}
 &= I(V,V)+\big[\langle D_sS,T\rangle\big]_{t=a}^{t=b}\big|_{s=0},\\
 I(V,W)
 &=\int_a^b
  \big(\langle D_tV,D_tW\rangle-\langle R(V,T)T,W\rangle\big)\,dt .
 \end{aligned}
\]
For a fixed-endpoint variation the boundary term vanishes.
The formula also applies to variations continuous and smooth
in \(s\) and piecewise smooth in \(t\), with fixed break
parameters and smooth matching traces in \(s\).
\end{theorem}

\begin{proof}
Differentiate the integral for \(E\), use torsion freeness,
and obtain \(E'(s)=\int_a^b\langle D_tS,T_s\rangle\,dt\).
A second derivative gives
\[
 E''(0)=\int_a^b
       \big(|D_tV|^2+\langle D_sD_tS,T\rangle\big)\,dt.
\]
Commuting derivatives replaces \(D_sD_tS\) by
\(D_tD_sS+R(V,T)V\) at \(s=0\).
Metric skew-symmetry gives
\(\langle R(V,T)V,T\rangle=-\langle R(V,T)T,V\rangle\).
Integration by parts in \(t\) leaves the endpoint term because
\(D_tT=0\). For a piecewise smooth variation, apply the same
calculation on each subinterval. The internal boundary terms
cancel since the trace \(F(s,t_i)\) has the same \(s\)-derivatives
from both sides and \(T\) is smooth along the central geodesic.
\end{proof}

\begin{theorem}[Index lemma]\label{thm:p202-index-lemma}
Let \(b>0\) and let \(\gamma:[0,b]\to M\) have no point conjugate to
\(\gamma(0)\) for \(0<t\le b\). For each terminal vector
\(z\in T_{\gamma(b)}M\) there is a unique Jacobi field \(J\)
with \(J(0)=0\) and \(J(b)=z\). Every continuous, piecewise
\(C^1\) field \(V\) with the same endpoints satisfies
\(I(V,V)\ge I(J,J)\), with equality exactly when \(V=J\).
In particular \(I(V,V)>0\) for a nonzero field vanishing at
both endpoints.
\end{theorem}

\begin{proof}
Use parallel coordinates and let \(A(t)\) solve
\[
 A''+\mathcal R A=0,\qquad A(0)=0,\quad A'(0)=\operatorname{id},
 \qquad \mathcal R(t)W=R(W,T)T.
\]
The preceding theorem says that \(A(t)\) is invertible on
\((0,b]\). In particular \(J(t)=A(t)A(b)^{-1}z\).
Since \(\mathcal R\) is symmetric, differentiating
\(A^{\mathsf T}A'-(A')^{\mathsf T}A\) gives zero, and its
initial value is zero. If \(W=V-J\) and \(W=A f\) on \((0,b]\),
then
\[
 |W'|^2-\langle\mathcal R W,W\rangle
 =|Af'|^2+\frac{d}{dt}\langle Af,A'f\rangle .
\]
Indeed the derivative of the last pairing contributes the
terms \(|A'f|^2+\langle Af,A''f\rangle\) and the two mixed
terms, which are equal by the Wronskian identity. The boundary
term at \(b\) is zero. At \(0\), \(A(t)=t\,\operatorname{id}
+O(t^3)\), \(A'(t)=\operatorname{id}+O(t^2)\), and \(W(t)=O(t)\),
so \(f\) stays bounded and the boundary term tends to zero.
Integrating on \([\varepsilon,b]\), summing across breaks,
and letting \(\varepsilon\downarrow0\) proves
\(I(W,W)=\int_0^b|Af'|^2\,dt\ge0\).
Equality makes \(f\) constant on each smooth piece, hence on
\((0,b]\) by continuity. As \(W(b)=0\), that constant is zero.
Finally \(I(J,W)=0\) by integration by parts and the Jacobi
equation, using \(W(0)=W(b)=0\). This proves the result.
\end{proof}

\begin{corollary}[A minimizing segment has no interior conjugate point]
\label{cor:p202-no-interior-conjugate}
If a geodesic segment minimizes length between its endpoints,
no point strictly between the endpoints is conjugate to its
initial point along that segment.
\end{corollary}

\begin{proof}
The geodesic has constant speed, and
\(\int_0^b|\dot\sigma|^2\,dt\ge L(\sigma)^2/b\) by
Cauchy--Schwarz. Thus it minimizes energy among competing curves
on the same interval. The second variation makes \(I\)
nonnegative on continuous piecewise smooth fields with zero
endpoints: realize such a field \(W\) by
\(F(s,t)=\exp_{\gamma(t)}(sW(t))\), using compactness for a
common range of small \(s\).

If there were a conjugate point at \(c\in(0,b)\), choose a
nonzero Jacobi field \(J\) vanishing at \(0\) and \(c\). Extend
it by zero on \([c,b]\) to get a piecewise smooth field \(W\).
Integration by parts gives \(I(W,W)=0\) and, for any smooth
zero-endpoint field \(Z\),
\[
 I(W,Z)=\langle D_tJ(c),Z(c)\rangle .
\]
The vector \(D_tJ(c)\) is nonzero by Jacobi uniqueness. Choose
\(Z\) with \(Z(c)=-D_tJ(c)\), using a parallel frame and a
cutoff. Then \(I(W+\varepsilon Z,W+\varepsilon Z)<0\) for
all sufficiently small positive \(\varepsilon\), a
contradiction.
\end{proof}

Conjugacy at the terminal point is consistent with minimality.
For example a great semicircle to the antipode of the round
sphere minimizes. A Jacobi field rotating the great circle about
the antipodal axis vanishes at both ends. This distinction
will also occur in the cut-point criterion. See
\cite[Chapter 10, Theorems 10.22 and 10.28]{LeeRM} and
\cite[Chapter 9, Section 2]{dC}.

\section{Cut time and the regular domain of distance}
\label{p202:G19}

Fix \(p\) on a connected complete Riemannian manifold of positive
dimension, and write \(S_pM\) for the unit sphere. For \(v\in S_pM\)
put
\[
 c(v)=\sup\{t\ge0:d(p,\exp_p(tv))=t\}\in(0,\infty].
\]
The times for which equality holds form an initial interval,
closed in \([0,\infty)\): every initial segment of a minimizing
segment minimizes, and the distance is continuous. The interval
has positive length by the Gauss lemma and the local
minimality of radial geodesics. When \(c(v)<\infty\),
\(\exp_p(c(v)v)\) is called the cut point in direction \(v\).
Their union is the cut locus \(\operatorname{Cut}(p)\).

\begin{theorem}[Cut-point criterion and polar domain]
\label{thm:p202-cut-polar}
If \(q=\exp_p(c(v)v)\) is a cut point, then either \(q\) is
conjugate to \(p\) along the minimizing segment or there are
two distinct minimizing geodesics from \(p\) to \(q\), or both.
At any \(0<t<c(v)\) the minimizing segment to \(\exp_p(tv)\)
is unique and nonconjugate. The function \(c\) is continuous
with values in the extended positive half-line, and
\[
 (t,v)\longmapsto \exp_p(tv),\qquad 0<t<c(v),
\]
is a diffeomorphism onto \(M\setminus(\{p\}\cup
\operatorname{Cut}(p))\). On this open set the distance from
\(p\) is smooth with gradient the unit radial vector.
\end{theorem}

\begin{proof}
Let \(0<t<c(v)\) and extend the segment to a larger minimizing
time. If a distinct minimizing segment reached the same point
at \(t\), append the rest of the first segment to it. The
concatenation would be minimizing and have constant speed
away from its join. A constant-speed minimizing path has no
corner: the first variation across the join would otherwise
give a shortening, or one may replace the two short pieces
by the unique minimizing geodesic in a convex normal
neighborhood. Thus the terminal velocities of the two first
segments agree. Uniqueness of the geodesic equation run
backward would make the two segments identical.
Nonconjugacy follows from
Corollary~\ref{cor:p202-no-interior-conjugate}.

We also need a stability observation. Suppose \(q\ne p\) has
a unique minimizing initial vector \(w\), and \(\exp_p\) is
nonsingular at \(w\). For a sequence \(q_i\to q\) choose
minimizing initial vectors \(w_i\), given by Hopf--Rinow and
parametrized on \([0,1]\). Since \(|w_i|=d(p,q_i)\), all lie
in a compact Euclidean ball of \(T_pM\). Any limit \(w'\) of
a subsequence has \(|w'|=d(p,q)\) and \(\exp_p w'=q\), so
\(w'=w\). Thus for every neighborhood of \(w\), all minimizing
initial vectors to all sufficiently nearby endpoints lie in
it. Take that neighborhood to be one on which the exponential
is a diffeomorphism. Every sufficiently nearby endpoint then
has its unique minimizing vector given by that local inverse.

At a finite cut time \(c(v)\), if the minimizing vector were
unique and nonconjugate, apply this observation at
\(w=c(v)v\) to points \(\exp_p((c(v)+\varepsilon)v)\).
The local inverse is \((c(v)+\varepsilon)v\), so the original
geodesic would remain minimizing beyond \(c(v)\), which is
impossible. This proves the dichotomy.

If \(r<c(v)\), the stability observation at \(rv\) shows that
\(\exp_p(rv')\) is reached by its unique nonconjugate
minimizing vector \(rv'\) for all \(v'\) near \(v\).
By the cut criterion, \(r<c(v')\). If \(r>c(v)\) is finite,
the strict inequality \(d(p,\exp_p(rv))<r\) holds and persists
for nearby directions, so \(c(v')<r\). The two conclusions
give continuity, including when \(c(v)=\infty\). The polar
domain is open. Every point \(q\ne p\) has a minimizing
segment by Hopf--Rinow. If it is not a cut point, its length
is strictly below its cut time. Conversely an endpoint
strictly before its cut time is unique and nonconjugate,
so it is no cut point in any direction. The polar exponential
is therefore bijective onto the stated complement and is
a local diffeomorphism by Theorem~\ref{thm:p202-exp-jacobi}.
The inverse is smooth. The radial coordinate equals the
distance; the Gauss lemma makes its gradient the unit radial
vector.
\end{proof}

See \cite[Chapter 10, Section on cut points]{LeeRM} and
\cite[Chapter 13, Section 2]{dC}.

\section{Constant-curvature Jacobi fields and Rauch comparison}
\label{p202:G20}

Define, for \(k\in\mathbb R\),
\[
 s_k(t)=
 \begin{cases}
 \sin(\sqrt{k}t)/\sqrt{k},& k>0,\\
 t,& k=0,\\
 \sinh(\sqrt{-k}t)/\sqrt{-k},& k<0.
 \end{cases}
\]
Thus \(s_k''+k s_k=0\), \(s_k(0)=0\), and \(s_k'(0)=1\).

\begin{proposition}[Constant-curvature polar metric]\label{prop:p202-model-polar}
On an \(n\)-manifold of constant sectional curvature \(k\),
the normal Jacobi field along a unit geodesic with
\(J(0)=0\) and \(D_tJ(0)=w\perp T(0)\) is
\(J(t)=s_k(t)P_t w\), where \(P_t\) is parallel transport.
On a normal ball the polar-coordinate metric is
\[
 dt^2+s_k(t)^2g_{S^{n-1}},
\]
where the angular metric is that of the unit sphere in \(T_pM\).
On the whole exponential domain, this same identity holds for
the pullback away from the origin.
\end{proposition}

\begin{proof}
The algebraic characterization of constant sectional curvature
gives \(R(W,T)T=kW\) for \(W\perp T\) and unit \(T\).
In a parallel frame each normal component solves the scalar
initial value problem for \(s_k\). For a direction \(v\) on the
unit sphere and \(w\in T_v S_pM\), the differential of the polar
exponential in the angular direction \(w\) is
\(d(\exp_p)_{tv}(tw)=s_k(t)P_t w\) by
Theorem~\ref{thm:p202-exp-jacobi}. Its radial derivative is \(T\).
These derivatives are orthogonal, so their metric pairings give
the claimed tensor identity. At a zero of \(s_k\) the angular
part of the pullback degenerates, as the formula indicates.
\end{proof}

\begin{theorem}[Rauch comparison for fields vanishing initially]
\label{thm:p202-rauch}
Let \(\gamma,\bar\gamma:[0,b]\to M,\bar M\) be unit-speed
geodesics in Riemannian manifolds of the same dimension \(n\ge2\).
Suppose for every \(t\) and every pair of unit normals
\(u\perp\dot\gamma(t)\), \(\bar u\perp\dot{\bar\gamma}(t)\) that
\[
 K(\dot\gamma(t),u)\le
 \bar K(\dot{\bar\gamma}(t),\bar u).
\]
Assume that along \(\bar\gamma\) there is no point conjugate to
\(\bar\gamma(0)\) for \(0<t\le b\). Then there is no such point
along \(\gamma\) for \(0<t\le b\), and for normal Jacobi fields
with
\[
 J(0)=\bar J(0)=0,\qquad
 |D_tJ(0)|=|\bar D_t\bar J(0)|
\]
one has \(|J(t)|\ge|\bar J(t)|\) for \(0\le t\le b\).
\end{theorem}

\begin{proof}
First take a time \(r>0\) such that both segments on \((0,r]\)
are free of conjugate points. If both initial derivatives
are zero the fields vanish. Otherwise the fields are nonzero
for each positive time in this interval. For any \(t\in(0,r]\)
normalize \(J\) and \(\bar J\) on \([0,t]\) by their respective
terminal norms, writing \(V=J/|J(t)|\) and
\(\bar V=\bar J/|\bar J(t)|\). Choose a fiberwise isometry \(P\)
between the normal bundles on \([0,t]\), parallel in both
variables and taking \(V(t)\) to \(\bar V(t)\). Such an isometry
is obtained by choosing it at \(t\) and extending it in
parallel frames. Its action preserves the derivative term in
the index form, while the sectional curvature bound gives
\(\bar I(PV,PV)\le I(V,V)\). The index lemma applied in
\(\bar M\) gives \(\bar I(\bar V,\bar V)\le \bar I(PV,PV)\).
For a Jacobi field starting at zero, integration by parts
in the index form gives only the terminal pairing. Hence
\[
 \frac{d}{dt}\log|\bar J(t)|=\bar I(\bar V,\bar V)
       \le I(V,V)=\frac{d}{dt}\log|J(t)|.
\]
As \(t\downarrow0\) the norm ratio tends to \(1\), because
in parallel coordinates each field is its initial derivative
times \(t\) plus \(O(t^3)\). Integration gives the comparison
up to \(r\).

The exponential differential is nonsingular for small
positive time. If \(\gamma\) has a conjugate time in
\((0,b]\), there is a first one \(t_*\), since the relevant
Jacobi matrix has continuous determinant. Take a nonzero
normal field \(J\) vanishing at \(0\) and \(t_*\), and choose
\(\bar J\) with initial derivative of the same norm.
The comparison holds on every interval ending below \(t_*\).
By continuity, \(0=|J(t_*)|\ge|\bar J(t_*)|>0\), impossible.
This proves both claims.
\end{proof}

The constant-curvature comparison is a specialization in
which \(\bar K\) or \(K\) is constant. Compare
\cite[Chapter 11, Proposition 11.3 and Theorem 11.9]{LeeRM}
and \cite[Chapter 10]{dC}.

\section{Polar volume and Bishop--Gromov comparison}
\label{p202:G21}

\begin{lemma}[The cut locus has volume zero]\label{lem:p202-cut-null}
On a complete connected Riemannian manifold of positive dimension,
\(\operatorname{Cut}(p)\) has zero Riemannian volume.
\end{lemma}

\begin{proof}
The cut-time function is continuous by
Theorem~\ref{thm:p202-cut-polar}. Its finite graph in
\((0,\infty)\times S_pM\) is measurable, and each section
at fixed direction consists of at most one point in the
\(t\)-variable. Fubini gives product measure zero. In coordinate
patches the map \((t,v)\mapsto\exp_p(tv)\) is smooth. A smooth
map between manifolds of the same dimension sends a null set
to a null set: cover the domain by countably many relatively
compact coordinate neighborhoods on which the coordinate
expression has bounded derivative, and apply the Lipschitz
null-set estimate on each. The volume densities are smooth
and positive, so coordinate null sets are volume null.
The image of the graph is exactly the cut locus.
\end{proof}

\begin{theorem}[Bishop--Gromov volume comparison]
\label{thm:p202-bishop-gromov}
Let \(M^n\) be complete and connected, with \(n\ge2\) and
\(\operatorname{Ric}\ge(n-1)k g\). Let \(V_k^n(r)\) denote
the volume of a radius-\(r\) ball in the simply connected
\(n\)-dimensional constant-curvature-\(k\) model. Then for
\(r>0\)
\[
 \frac{\operatorname{Vol} B(p,r)}{V_k^n(r)}
\]
is nonincreasing as a function of \(r\), on the interval
\(0<r<\pi/\sqrt{k}\) when \(k>0\), and on the whole positive
half-line when \(k\le0\). Its limit at \(r=0\) is \(1\).
For \(k>0\) the comparison also holds at the upper endpoint
by taking the one-sided limit.
\end{theorem}

\begin{proof}
Put \(m=n-1\). Fix a unit direction \(v\), and identify the
normal spaces along \(\gamma(t)=\exp_p(tv)\) in a parallel
frame. The normal Jacobi matrix \(A(t)\) has the initial
conditions \(A(0)=0\), \(A'(0)=\operatorname{id}\) and
satisfies \(A''+\mathcal R A=0\), where
\(\mathcal R W=R(W,T)T\). For \(0<t<c(v)\) the determinant
\(j(t,v)=\det A(t)\) is positive, since it is nonzero there
and \(A=t\,\operatorname{id}+O(t^3)\) near zero. It is the
density of Riemannian volume relative to \(dt\,d\omega(v)\)
in the polar domain. Indeed, the angular differentials
of the exponential are the columns of \(A\) in orthonormal
angular directions and the radial differential is a unit
vector orthogonal to them.

The Wronskian identity in the index lemma makes
\(U=A'A^{-1}\) self-adjoint. Differentiation gives
\[
 U'+U^2+\mathcal R=0,\qquad
 (\log j)'=\operatorname{tr}U.
\]
Write \(\theta=\operatorname{tr}U/m\). Since
\(\operatorname{tr}(U^2)\ge(\operatorname{tr}U)^2/m\) and
\(\operatorname{tr}\mathcal R=\operatorname{Ric}(T,T)\), we get
\(\theta'+\theta^2+k\le0\). On the interval where
\(s_k>0\), \(\theta_k=s'_k/s_k\) solves
\(\theta_k'+\theta_k^2+k=0\). Subtracting gives
\[
 (\theta-\theta_k)' +(\theta+\theta_k)(\theta-\theta_k)\le0.
\]
The integrating factor \(j^{1/m}s_k\) is positive and has
logarithmic derivative \(\theta+\theta_k\).
As \(t\downarrow0\), \(A'= \operatorname{id}+O(t^2)\)
and \(s'_k=1+O(t^2)\), so
\(\theta=1/t+O(t)\), \(\theta_k=1/t+O(t)\), and the
product of the integrating factor with \(\theta-\theta_k\)
tends to zero. Thus \(\theta\le\theta_k\). Equivalently
\(j(t,v)/s_k(t)^m\) is nonincreasing before the cut time.
Its initial limit is \(1\).

Define \(j_*(t,v)=j(t,v)\) for \(0<t<c(v)\) and set it
equal to zero after and at the cut time. The ratio
\(j_*(t,v)/s_k(t)^m\) stays nonincreasing on the comparison
interval in each direction, since its only change at the
cut time is a downward jump. By the polar diffeomorphism
and the nullity of the omitted cut locus,
\[
 \operatorname{Vol} B(p,r)=\int_0^r\int_{S_pM}
                         j_*(t,v)\,d\omega(v)\,dt,\qquad
 V_k^n(r)=\omega_{n-1}\int_0^r s_k(t)^m\,dt .
\]
The second identity follows from
Proposition~\ref{prop:p202-model-polar} and the normal polar
domain of the model up to the prescribed radii.
The volume ratio is therefore a weighted average over
\((0,r)\) of the nonincreasing function
\[
 a(t)=\frac{1}{\omega_{n-1}}
       \int_{S_pM}\frac{j_*(t,v)}{s_k(t)^m}\,d\omega(v)
\]
with weight \(s_k(t)^m>0\). Extending the interval can only
lower the average: all new values are no larger than the
previous average. The directionwise limit \(1\) and the
bound \(0\le j_*/s_k^m\le1\) give the initial limit of
\(a\), hence of the volume ratio. At a finite upper endpoint
the increasing union of the open balls gives the endpoint
value of each volume by monotone convergence.
\end{proof}

See \cite[Chapter 11, Theorem 11.19]{LeeRM}.

\chapter{Complete Models and Global Consequences}
\label{ch:p202-global-geometry}

\section{Local isometries, covers, and space forms}
\label{p202:G22}

\begin{theorem}[Complete local isometries]\label{thm:p202-complete-cover}
Let \(F:X\to Y\) be a local isometry of connected Riemannian
manifolds without boundary. If \(X\) is complete, then \(F\)
is a covering map. Conversely, the pullback Riemannian metric
on a smooth covering of a complete Riemannian manifold is
complete on each connected component.
\end{theorem}

\begin{proof}
A local isometry preserves the Levi--Civita connection, by
metric compatibility, torsion freeness, and uniqueness.
It therefore sends geodesics to geodesics. From any point
\(F(x)\) every point of \(Y\) is reachable by a finite chain
of short geodesic segments, using connectedness and local
normal neighborhoods. Each segment can be lifted from its
initial point by lifting its velocity and solving the geodesic
equation in \(X\) for the full segment. Completeness and
Hopf--Rinow supply existence for its entire length of time,
and geodesic uniqueness in \(Y\) identifies the projection.
Thus \(F\) is onto.

Fix \(y\in Y\) and a normal ball \(U\) about \(y\). For
\(x\in F^{-1}(y)\) define
\[
 s_x(z)=\exp_x\big((dF_x)^{-1}\exp_y^{-1}(z)\big),
 \qquad z\in U.
\]
This is the endpoint of the lift of the radial segment, so
\(F s_x(z)=z\). Because \(F\) is a local diffeomorphism, the
image of this smooth section is open and maps diffeomorphically
to \(U\). From any point of \(F^{-1}(U)\) the reversed radial
segment lifts and ends in \(F^{-1}(y)\). Its reversal is one of
the radial lifts above. Uniqueness of the reversed lift also
shows that the images of distinct sections are disjoint. This
is an evenly covered neighborhood.

For the converse, the projection of a geodesic in a pullback
metric is a geodesic of the base. Extend that projection for
all real times. By path lifting on successive compact
intervals and uniqueness on their overlaps, it has a lift
for all real times through the given initial point. The lift
is smooth and geodesic, as seen in covering charts, and it
extends the original geodesic. Hopf--Rinow gives metric
completeness in each component.
\end{proof}

\begin{theorem}[Killing--Hopf classification]\label{thm:p202-killing-hopf}
Let \(M^n\) be complete and connected, without boundary, with
\(n\ge2\) and constant sectional curvature \(k\).
Its Riemannian universal cover is isometric to the simply
connected constant-curvature model \(X_k^n\): Euclidean
space, a round sphere of radius \(1/\sqrt{k}\), or hyperbolic
space, according as \(k=0\), \(k>0\), or \(k<0\).
In particular a simply connected \(M\) is isometric to the
model. In general \(M\) is the quotient of the model by its
group of covering transformations, acting by isometries
freely and properly discontinuously.
\end{theorem}

\begin{proof}
The simply connected, complete curvature models and their
geodesics were proved in Theorem~\ref{thm:p202-models}.
At the hyperboloid pole \((0,a)\), the spatial component
\(a\sinh(t/a)v\) of the geodesic formula makes the
exponential a global diffeomorphism. On the sphere the
formula gives the polar chart up to the antipode.
These are the model polar domains used in the volume
comparison. The radial coordinate \(r\) is continuous
at the pole and also at the antipode on the sphere,
where its value is \(\pi a\). Its gradient has unit norm
in the chart. A path from the pole to radius \(r_0>0\)
first reaches every level \(\varepsilon<r_0\). From its
last visit to level \(\varepsilon\) before its first visit
to \(r_0\), integrating \(|dr|\) bounds its length below
by \(r_0-\varepsilon\). Use \(r_0<\pi a\) on the sphere,
and take the limit for its antipode. Thus radial length
equals distance on the claimed domains. The isometry
groups are transitive: use translations in Euclidean
space, orthogonal transformations on the sphere, and
translations and positive dilations in the hyperbolic
half-space model.

Choose \(x\in X_k^n\), \(y\in M\), and a linear isometry
\(A:T_xX_k^n\to T_yM\). On a sufficiently small normal ball
about \(x\),
\[
 z\longmapsto\exp_y\big(A\exp_x^{-1}(z)\big)
 \tag{*}
\]
is a local isometry by the identical polar metric formulas
in Proposition~\ref{prop:p202-model-polar}, and by smoothness
at the center. More generally it is a local isometry on any
model normal ball of radius less than the first positive
zero of \(s_k\), since \(\exp_y\) is everywhere defined by
completeness. The pullback polar formula proves metric
preservation even if \(\exp_y\) ceases to be one-to-one on
that ball. Every germ of a local isometry from the model to
\(M\) is determined by its value and differential at its
center, since it takes radial geodesics to the geodesics
specified by those data. It therefore extends to a local
isometry on such a ball by \((*)\).

Here are details of globalizing the germs. Form the space
of all germs of local isometries from \(X_k^n\) to \(M\),
projecting a germ to its point of definition. Give it the
topology whose basic open sets are the germs of a single
local isometry defined on an open set. The basis property
follows from the definition of equality of germs.
By homogeneity of each model choose a fixed \(r>0\) so that
radius-\(2r\) balls are normal and \(2r\) is below any first
positive zero of \(s_k\). For any \(x\) put \(U=B(x,r)\).
Every germ at \(x\) extends to a map on \(U\).
Two such maps have disjoint sets of germs on \(U\) unless
the center germs agree: the set where two local isometries
have equal value and differential is closed in their
connected common domain and is open by radial uniqueness.
Every germ at a point \(z\in U\) extends to \(B(z,2r)\),
which contains \(U\), so comes from one of the center germs.
The projection is consequently a covering. Since \(X_k^n\)
is simply connected, the component through the chosen germ
projects homeomorphically to \(X_k^n\).
Reading the values of its germs gives a globally defined
local isometry \(X_k^n\to M\), smooth by the local
representatives.

By Theorem~\ref{thm:p202-complete-cover} this map is an
isometric covering. Its simply connected source makes it
the universal cover. If \(M\) is simply connected the map
is a diffeomorphism and hence an isometry. In general the
covering transformation group acts transitively on each
fiber and freely, by the universal covering classification
in Sections~\ref{p202:T07}--\ref{p202:T08}. The evenly covered
neighborhoods give proper discontinuity and identify the
quotient with \(M\). A covering transformation preserves
the pullback metric because its composition with the
covering projection is the projection itself.
\end{proof}

Compare \cite[Chapter 12, Theorem 12.4]{LeeRM}
and \cite[Chapters 7--8]{dC}.

\section{The Cartan--Hadamard theorem}
\label{p202:G23}

\begin{theorem}[Cartan--Hadamard]\label{thm:p202-cartan-hadamard}
Let \(M\) be a complete connected Riemannian manifold of
nonpositive sectional curvature. For each \(p\in M\),
\(\exp_p:T_pM\to M\) is a covering map. If \(M\) is simply
connected, each such exponential is a diffeomorphism and
any two points are joined by a unique minimizing geodesic,
up to parametrization.
\end{theorem}

\begin{proof}
Along a unit geodesic consider a normal Jacobi field
\(J(0)=0\), \(D_tJ(0)=w\ne0\).
The Jacobi equation gives
\[
 (|J|^2)''=2|D_tJ|^2-2\langle R(J,T)T,J\rangle\ge0 .
\]
The function is strictly increasing for small positive
time, since \(J(t)=tw+O(t^3)\) in a parallel frame, and
convexity prevents any later zero. On \(t>0\), differentiate
\(|J|\) twice. Cauchy--Schwarz gives
\[
 |J|\, (|J|)''
 = |D_tJ|^2-\frac{\langle D_tJ,J\rangle^2}{|J|^2}
       -\langle R(J,T)T,J\rangle\ge0 .
\]
The right derivative of \(|J|\) at zero is \(|w|\), so
\(|J(t)|\ge t|w|\). A Jacobi field starting at zero and
vanishing again is normal unless it is zero. There are
therefore no conjugate times and the exponential is a
local diffeomorphism.

Pull the metric back to \(T_pM\). At \(tv\) with
\(|v|=1\), decompose a tangent vector into its Euclidean
radial and orthogonal parts. The differential sends the
radial part isometrically to the geodesic direction and
the orthogonal part to a Jacobi field value orthogonal
to that direction. For an orthogonal vector \(u\) the
initial derivative is \(u/t\), so its image norm is at
least \(|u|\). This uses the Gauss lemma and
Theorem~\ref{thm:p202-exp-jacobi}. At the origin the
differential is the identity. The pullback norm thus
dominates the Euclidean norm everywhere. A Cauchy
sequence for the pullback distance is Euclidean Cauchy.
Near its Euclidean limit the pullback norm has a uniform
upper bound by a constant times the Euclidean norm on
a small closed ball, so straight segments show convergence
also in the pullback distance. The pullback metric is
complete, and the exponential is a local isometry for
this metric. Theorem~\ref{thm:p202-complete-cover} makes
it a covering. When \(M\) is simply connected it is a
diffeomorphism. Hopf--Rinow supplies a minimizing geodesic,
and injectivity of \(\exp_p\) then gives uniqueness.
\end{proof}

In particular the universal cover in this setting is
diffeomorphic to \(\mathbb R^n\); completeness is part of
the assertion. See \cite[Chapter 12, Theorem 12.8]{LeeRM}.

\section{The Bonnet--Myers theorem}
\label{p202:G24}

\begin{theorem}[Bonnet--Myers]\label{thm:p202-bonnet-myers}
Let \(M^n\) be a complete connected Riemannian manifold
without boundary, \(n\ge2\). If
\(\operatorname{Ric}\ge(n-1)k g\) for \(k>0\), then
\(\operatorname{diam}M\le\pi/\sqrt{k}\), \(M\) is compact,
and its fundamental group is finite.
\end{theorem}

\begin{proof}
Let \(\gamma:[0,L]\to M\) be a unit-speed minimizing
geodesic of positive length, which exists between any
distinct points by Hopf--Rinow. Take a parallel orthonormal
normal frame \(E_1,\ldots,E_{n-1}\) and consider
\(V_i(t)=\sin(\pi t/L)E_i(t)\). Minimality and the fixed-endpoint
second variation give \(I(V_i,V_i)\ge0\) for every \(i\).
Summing and using the Ricci bound gives
\[
 \begin{aligned}
 0\le \sum_i I(V_i,V_i)
 &\le (n-1)\int_0^L
 \left((\pi/L)^2\cos^2(\pi t/L)-k\sin^2(\pi t/L)\right)\,dt\\
 &=\frac{(n-1)L}{2}\big((\pi/L)^2-k\big).
 \end{aligned}
\]
Since \(n-1>0\), this implies the diameter bound.
The whole manifold is closed and bounded in itself, so
Hopf--Rinow gives compactness.

The universal covering, with the pullback metric, is
complete by Theorem~\ref{thm:p202-complete-cover} and
satisfies the same Ricci bound by local isometry.
It too is compact by the preceding argument. Its
fiber over \(p\) is both closed and discrete, hence
finite. The cardinality of that fiber equals the order
of the fundamental group by universal covering theory.
\end{proof}

The dimension assumption matters for this normalization:
when \(n=1\) the lower bound is vacuous and the Euclidean
line has unbounded diameter. See
\cite[Chapter 12, Theorem 12.24]{LeeRM}
and \cite[Chapter 9, Section 3]{dC}.

\section{The local Gauss--Bonnet formula with corners}
\label{p202:G25}

Let \(S\) be an oriented Riemannian surface and \(J\) the
rotation through \(+\pi/2\) in its oriented tangent planes.
For a regular oriented curve with unit tangent \(T\), the
signed geodesic curvature is \(k_g=\langle D_TT,JT\rangle\).
The positive boundary orientation puts the interior on
the left, so \(JT\) is the inward unit normal to each
smooth boundary arc, viewed as a vector in the surface.
At a corner of interior angle \(0<\alpha<2\pi\), the
signed exterior angle is \(\varepsilon=\pi-\alpha\).
In particular a reflex corner contributes a negative angle.

\begin{lemma}[Euler characteristic of a planar region]
\label{lem:p202-planar-region}
Let \(D\) be the compact closure of a nonempty connected planar
domain, whose boundary consists of \(b\ge1\) disjoint simple closed
piecewise smooth regular curves with finitely many noncuspidal
corners. Assume that locally at each boundary point the region
occupies one side of the curve. Then \(\chi(D)=2-b\).
Exactly one boundary curve has the bounded side on the left in
the induced boundary orientation. Every other boundary curve
has the bounded side on the right.
\end{lemma}

\begin{proof}
By Jordan separation, Theorem~\ref{thm:p202-invariance-domain}
in dimension two, each boundary curve \(C\) separates \(S^2\)
into two open components with common boundary \(C\).
The connected interior of \(D\) lies on one side. The other
side is disjoint from \(D\), since \(D\) is the closure of its
interior. That side is an entire component of \(S^2\setminus D\):
it is connected and its boundary lies in \(D\).
Every component of \(S^2\setminus D\) has a boundary point in
\(\partial D\). Near such a point the local outside of the
boundary curve is connected and belongs to its outside component
in \(S^2\setminus C\). Thus the components just described account
for the whole complement. They are distinct because their
boundaries are the distinct curves. Exactly one contains the
point at infinity. It is therefore precisely its boundary
curve whose bounded side contains the interior of \(D\).
This gives the orientation statement.

For duality, \(D\) is a topological surface with collared boundary.
At a smooth boundary point the local graph description is
immediate. At a corner the two consistently directed one-sided
tangents are not opposite. The sum of their inward normals
has positive dot product with each inward normal. The same
strict inequalities hold along both arcs in a neighborhood of
the corner. Use such local directions and a smooth partition
of unity to choose a smooth ambient vector field strictly inward
along all boundary arcs, including the one-sided inequalities
at corners. In a flow box for this field
\cite[Theorem 6.5]{P104}, the projection transverse
to the field is strictly monotone along the locally oriented
boundary on both sides of a corner. Consequently the boundary is
a graph in the flow box, and positive flow time gives the local
inward product. These products give a single collar by flowing
the entire boundary for a uniform short time. For injectivity,
if collisions occurred for arbitrarily small nonnegative times,
compactness would give two boundary base points converging to the
same point. They would eventually be in one graph flow box, where
the transverse coordinate determines the base point and then the
flow coordinate determines the time. The local graph products show
that the collar contains a neighborhood of the boundary in \(D\).
Pushing the collar parameter a little inward also shows that the
interior and \(D\) are homotopy equivalent.

Now \(D\) is compact and locally contractible, so Alexander duality
over \(\Q\) gives \(\dim H^1(D;\Q)=b-1\). For a chain complex over a
field evaluation identifies cohomology with the dual of homology:
a functional on cycles killing boundaries extends to all chains by
choosing a complement, and a cocycle zero on cycles is a coboundary
by extending the induced functional on boundaries in the previous
degree. A vector space with finite-dimensional dual has the same
dimension as its dual, since functionals separating any finite
independent family extend from its span. Thus
\(\dim H_1(D;\Q)=b-1\). Theorem~\ref{thm:p202-lefschetz}
gives \(H_i(D;\Q)=0\) for \(i>2\) and identifies \(H_2(D;\Q)\) with
\(H^0(D,\partial D;\Q)=0\), since \(D\) is connected with nonempty
boundary. The degree-zero homology has dimension one, so the
alternating sum is \(2-b\).
\end{proof}

\begin{lemma}[Turning of a positively oriented simple plane boundary]
\label{lem:p202-plane-turning}
Let a piecewise smooth regular Jordan curve in an oriented plane
have its bounded side on the left. Assume it has finitely many
corners and no cusps. Along each smooth arc choose a continuous
tangent angle. Its total change along the arcs, with the signed
exterior angles added at the corners, equals \(2\pi\).
Reversing the orientation changes this total to \(-2\pi\).
\end{lemma}

\begin{proof}
The moving tangent, with the shorter signed rotation
inserted at each corner, is a continuous loop in the unit
circle. The desired sum is the final minus initial angle
of a lift of this loop to \(\mathbb R\), and so is \(2\pi\)
times its winding number. Approximate the boundary by a
polygon inscribed at its corners and at a fine subdivision
of its smooth arcs. Here are two details of this reduction.

First, the polygon is simple for all sufficiently fine
subdivisions. Near each boundary parameter the curve is
a graph in some direction with strictly monotone projection
in the oriented parameter. At a smooth point use a direction
close to its tangent. At a corner there is a direction
having positive dot product with both oriented one-sided
tangents, since they are not opposite. For a fine enough
mesh, any pair of intersecting polygon edges would have
endpoints close in the plane and therefore close in the
boundary parameter, since the inverse of the Jordan
parametrization on its image is uniformly continuous.
They would then lie in one of those graph neighborhoods.
In that neighborhood the projection intervals of different
edge interiors are disjoint, ruling out such intersections.
The integral of \(x\,dy-y\,dx\) around the polygon converges
by Riemann sums on the smooth arcs to the corresponding
positive integral on the curve. For that integral, the planar
area formula follows by cutting out small disks around the
corners and applying Stokes before taking the limit, as
detailed below. For the polygon the same area formula follows
by summing over a straight triangulation. Its positive integral
therefore gives positive orientation for sufficiently fine meshes.

Second, the winding of the tangent loop is unchanged by
passing to the polygon. On a small smooth subarc, the chord
direction is uniformly close to all the unit tangents on
the subarc, by expressing the chord as the integral of the
velocity. The short angular interpolation between consecutive
chord directions therefore lies in a common angle chart
with the corresponding smooth tangents. At each original
corner the two limiting chord directions approach the two
one-sided tangent directions. Since those directions are
not opposite, the shorter rotations vary continuously
under this perturbation. Aligning the time intervals of
the arc traversals and rotations gives a homotopy of the
two direction loops, by interpolation in these angle
charts.

It remains to compute the turning for a simple polygon.
Triangulate its interior by straight triangles. For an
elementary construction one can choose a coordinate
direction in which the vertices have distinct first
coordinates and the edges are nonvertical, cut by the
vertical lines through the vertices, and subdivide each
resulting nondegenerate convex strip piece by joining an
interior point to the vertices on its subdivided boundary.
In a strip the edges stay in fixed vertical order, since
they do not meet in its interior. The intervals of the
bounded region between consecutive edges therefore sweep
out exactly those strip pieces. Use common side subdivisions
along the cut lines. Additional boundary vertices on
straight edges have exterior angle zero. At internal
vertices the triangle angles sum to \(2\pi\), and at
boundary vertices they sum to the polygon's interior
angles. Write \(F,E_i,E_b,V_i,V_b\) for the face count
and the internal and boundary edge and vertex counts.
The sum of the exterior angles is
\(\pi V_b-\pi F+2\pi V_i\), since the sum of all triangle
angles is \(\pi F\). Now \(3F=2E_i+E_b\), \(E_b=V_b\),
and the triangulated closed region has Euler characteristic
\(V_i+V_b-E_i-E_b+F\) by Section~\ref{p202:T21}.
The bounded region is connected and has the polygon as its
single boundary curve, by Jordan separation. Its Euler
characteristic is therefore \(1\) by Lemma~\ref{lem:p202-planar-region},
which gives \(2\pi\). Reversing the boundary traversal reverses
the direction loop up to a fixed rotation by \(\pi\), so its
winding number changes sign.
\end{proof}

\begin{theorem}[Local Gauss--Bonnet]\label{thm:p202-local-gb}
Let \(D\) be a compact disk in an oriented coordinate
neighborhood of a Riemannian surface, with simple piecewise
smooth regular boundary and finitely many noncuspidal
corners. Use the positive boundary orientation. Then
\[
 \int_D K\,dA+\int_{\partial D}k_g\,ds+
   \sum_{\text{corners}}\varepsilon=2\pi .
\]
The boundary integral is over the smooth arcs. More generally,
for a connected compact chart region as in
Lemma~\ref{lem:p202-planar-region}, the same left-hand side is
\(2\pi\chi(D)\), using the induced orientation on every boundary
component.
\end{theorem}

\begin{proof}
Choose a positive orthonormal frame \(e_1,e_2\) on the
coordinate neighborhood by applying Gram--Schmidt to a
positive coordinate frame. Define its connection form
\(\omega\) by
\(\nabla_Xe_1=\omega(X)e_2\) and
\(\nabla_Xe_2=-\omega(X)e_1\).
The curvature definition gives
\[
 R(X,Y)e_1=d\omega(X,Y)e_2,\qquad d\omega=-K\,dA .
\]
In the first identity the terms quadratic in \(\omega\)
cancel. The second follows from the definition of \(K\)
and metric skew-symmetry.

Along a smooth arc write \(T=\cos\theta\,e_1+\sin\theta\,e_2\).
Differentiating with respect to arc length gives
\(k_g=\theta'+\omega(T)\). In the chart coordinates the
direction loop with short corner rotations has winding number
\(+1\) on the outer boundary and \(-1\) on each inner boundary,
by the preceding two lemmas. The changes
of vector coordinates from the positive coordinate frame
to its Gram--Schmidt frame are upper triangular with
positive diagonal. Their straight interpolation to the
identity stays invertible and orientation preserving.
Applying this interpolation to tangent directions,
normalizing, and interpolating the short rotations
preserves winding, since independent one-sided tangent
directions cannot become opposite under these maps, and
equal one-sided directions remain equal. Thus the sum
of the smooth angle changes and signed corner changes
in the orthonormal frame is \(2\pi(2-b)=2\pi\chi(D)\).
For a disk there is one boundary curve and the sum is \(2\pi\).

To apply Stokes also at reflex corners, remove small closed
coordinate disks of radius \(\epsilon\) about the corner
points. On each incident arc the derivative of the radial
distance from the corner, measured away from the corner,
tends to its positive initial speed. The arcs therefore
cross the small circles transversely. The truncated domain
has ordinary transverse corners, where Stokes applies.
The new circular portions have total length \(O(\epsilon)\),
and the removed areas are \(O(\epsilon^2)\). A smooth
one-form has bounded coefficients on these compact chart
neighborhoods, so the extra boundary integrals and the
integrals of its derivative over the removed regions tend
to zero. The retained boundary arcs converge to the original
integrals. This proves the Stokes identity for the region,
including reflex corners.

Now integrate \(k_g=\theta'+\omega(T)\) on the arcs and
add the corners. With the positive boundary orientation,
the Stokes identity gives
\[
 \int_{\partial D}k_g\,ds+\sum\varepsilon
 =2\pi\chi(D)+\int_{\partial D}\omega
 =2\pi\chi(D)-\int_DK\,dA .
\]
This proves the formula.
\end{proof}

The connection-form proof and its turning-number input can
be compared with \cite[Chapter 9, Theorems 9.1 and 9.3]{LeeRM}.
For the Stokes convention, see \cite[Chapter 10]{P104}.

\section{Global Gauss--Bonnet for compact surfaces}
\label{p202:G26}

\begin{lemma}[Planar regions cut out on a compact surface]
\label{lem:p202-surface-regions}
Let \(S\) be a compact smooth surface with smooth boundary, possibly
empty. There is a finite embedded graph \(G\) containing \(\partial S\),
with smooth regular open edges, whose complementary components have
closures \(D_j\) of the type in Lemma~\ref{lem:p202-planar-region}
in individual charts. Near a vertex these regions occupy the sectors
of the graph. If \(G\) has \(V\) vertices and \(E\) edges, then
\[
 \chi(S)=V-E+\sum_j\chi(D_j).
\]
\end{lemma}

\begin{proof}
If \(\partial S\ne\varnothing\), work in the double \(\widehat S\).
To construct it, glue two copies of \(S\) by the identity on the
boundary. The quotient is compact Hausdorff, since the identification
is a closed equivalence relation on the compact disjoint union.
The smooth collar of \cite[Theorem 6.5]{P104} combines its two
copies into signed collar charts, giving a smooth surface across
the seam. The seam is a smooth compact one-dimensional submanifold
without boundary. It has finitely many circle components. Indeed,
cover it by interiors of finitely many closed coordinate intervals
and mark their endpoints. Each component of the unmarked set
lies in one such interval, since it cannot leave that interval
without crossing an endpoint. The cuts yield a finite graph with
two incident rays at every marked point, hence a disjoint union
of circles.
If \(\partial S=\varnothing\), use \(\widehat S=S\) with no seam.

Cover \(S\) by the interiors of finitely many closed coordinate
disks in \(\widehat S\). Choose them initially with room to increase
each radius to a larger radius whose closed disk still lies in the
same chart. Such inner and outer radii exist around each point by
the openness of the coordinate chart, and compactness gives the
finite cover by the inner disks. In each chart choose the final
radius between its two bounds. Do this successively so that the
new circle is transverse to all earlier circles and the seam and
avoids their intersection points. In fact the squared radius
function restricted to the portions of those curves in the chart
has a measure-zero set of critical values by the one-dimensional
case of \cite[Theorem 4.1]{P104}. Each regular value gives the
required transversality. There are only finitely many previously
formed intersections: each pair of compact curves intersects in a
compact discrete set. Excluding their finitely many radius values
still leaves choices. The final disks continue to cover \(S\).

Intersect the resulting arrangement of circles and seam with \(S\).
Put vertices at the crossings in \(S\), and place a vertex on any
whole circle that has no crossings. Cutting the curves at these
vertices gives a finite graph \(G\), with loops allowed. Each
component of \(S\setminus G\) lies entirely on one side of the
boundary of each coordinate disk. Since the disks cover \(S\),
the component lies in the interior of one disk and its closure
lies in that closed disk, still in the chart. At an interior
crossing the four sectors have distinct pairs of disk memberships.
At a crossing on the seam the two sectors inside \(S\) have
distinct memberships in the disk whose boundary crosses the seam.
Along a single curve the sides are distinguished by one membership,
or there is just one side at the boundary of \(S\).
It follows that one component can occupy at most one sector at
any given boundary point of its closure. A sector meeting that
component is wholly in it, by connectedness of the sector.
Consequently every such closure is a planar region with regular
boundary arcs and ordinary corners. Its boundary is made of graph
edges, and is a disjoint union of simple closed curves, since it
has exactly two boundary rays at every boundary vertex. Its boundary
is nonempty: otherwise the closure would be an open compact surface
inside a coordinate chart. Each component therefore has at least
one edge side on its boundary, so there are only finitely many
components.

Use the collars from the proof of Lemma~\ref{lem:p202-planar-region}
in all the \(D_j\). Let \(W\) consist of \(G\) and a short
half-open inner collar of every region. The compact inner cores
are its complement, so \(W\) is open in \(S\).
Contracting the collar parameter to zero gives a deformation
retraction of \(W\) onto \(G\). The contractions agree on the graph,
and continuity follows by gluing over the finite closed cover
of \(W\) by the intersections with \(G\) and the \(D_j\).
The second open set is the disjoint union of the interiors of
the \(D_j\), each homotopy equivalent to \(D_j\) by the collar.
The overlap is a disjoint union of open collars, homotopy
equivalent to the union of the boundary circles, counted once
for each region they bound. All these spaces have finite-dimensional
rational homology in only finitely many degrees. The Mayer--Vietoris
sequence proves the same for \(S\) and gives the Euler identity:
the overlap contributes zero, and the graph contributes \(V-E\).
\end{proof}

\begin{theorem}[Gauss--Bonnet for compact Riemannian surfaces]
\label{thm:p202-global-gb}
Let \(S\) be a compact smooth Riemannian surface, either
closed or with smooth boundary. On the boundary, let
\(k_g=\langle D_TT,N_{\mathrm{in}}\rangle\), with \(T\) either
unit tangent and \(N_{\mathrm{in}}\) the inward unit vector
normal to the boundary in the surface. Then
\[
 \int_S K\,dA+\int_{\partial S}k_g\,ds=2\pi\chi(S).
\]
Here \(dA\) is the Riemannian area density and \(\chi(S)\)
is the topological Euler characteristic. The assertion also
holds for nonorientable or disconnected \(S\).
\end{theorem}

\begin{proof}
Choose the graph and regions from Lemma~\ref{lem:p202-surface-regions}.
Apply Theorem~\ref{thm:p202-local-gb} in each chart with an
orientation chosen for that chart. Its curvature integral uses
the area density, and its boundary curvature is the pairing of
\(D_TT\) with the inward normal of the region. Both interpretations
are independent of the chosen chart orientation.
All interior edge integrals cancel, since the inward normals of
the two regions are opposite and \(D_TT\) is unchanged by replacing
\(T\) by \(-T\). On edges in \(\partial S\) the inward normal is the
one for \(S\). Thus only the asserted boundary integral remains.

At a vertex \(v\) let \(d_v\) denote the number of incident half-edges.
At an interior vertex the \(d_v\) sectors fill a tangent plane and
their angles sum to \(2\pi\). At a boundary vertex the \(d_v-1\)
sectors fill a tangent half-plane and their angles sum to \(\pi\).
The sum of exterior angles at either type of vertex is therefore
\(\pi(d_v-2)\). This also covers an added vertex in a smooth edge,
where the smooth sector angles are \(\pi\). Every edge has two
half-edges, including a loop, so the exterior-angle sum over all
regions is \(2\pi(E-V)\). The graph has area zero, being a finite
union of smooth arcs and points. Summing the region formulas gives
\[
 \int_S K\,dA+\int_{\partial S}k_g\,ds
 =2\pi\bigl(V-E+\textstyle\sum_j\chi(D_j)\bigr)
 =2\pi\chi(S).
\]
The argument applies on every component without requiring a global
orientation. The empty surface gives zero on both sides.
\end{proof}

For the closed surface formula see
\cite[Chapter 9, Theorem 9.7]{LeeRM}. For the orientation
cover and the distinction between volume forms and
densities, see \cite[Chapter 9]{P104}.

\backmatter
\begingroup
\small
\begin{thebibliography}{99}
\bibitem{Hatcher} A. Hatcher, \emph{Algebraic Topology}, Cambridge
University Press, 2002, with the author's subsequent corrections.
\bibitem{MKP} J. R. Munkres, S. G. Krantz, and H. R. Parks,
\emph{Elements of Algebraic Topology}, 2nd ed., CRC Press, 2025.
\bibitem{MT} J. R. Munkres, \emph{Topology}, 2nd ed.,
Prentice Hall, 2000.
\bibitem{Weibel} C. A. Weibel, \emph{An Introduction to Homological
Algebra}, Cambridge University Press, 1994.
\bibitem{LeeSM} J. M. Lee, \emph{Introduction to Smooth Manifolds},
2nd ed., Springer, 2013.
\bibitem{LeeRM} J. M. Lee, \emph{Introduction to Riemannian Manifolds},
2nd ed., Springer, 2018, with the author's corrections.
\bibitem{dC} M. P. do Carmo, \emph{Riemannian Geometry}, translated
by F. Flaherty, Birkh\"auser, 1992.
\bibitem{D10} MathGenome, \emph{Point-Set Topology}, unified notes D10.
\bibitem{D13} MathGenome, \emph{Ordinary Differential Equations and
Dynamical Systems}, unified notes D13.
\bibitem{P103} MathGenome, \emph{Rings, Modules, Representations, and
Homological Algebra}, unified notes P1-03.
\bibitem{P104} MathGenome, \emph{Smooth Manifolds and Intrinsic
Calculus}, unified notes P1-04.
\bibitem{P106} MathGenome, \emph{Variational Methods, Convex Analysis,
and Optimality}, unified notes P1-06.
\bibitem{P201} MathGenome, \emph{Classical and Weak Linear Partial
Differential Equations}, unified notes P2-01.
\end{thebibliography}
\endgroup
\end{document}

